\RequirePackage{xr-hyper}
\documentclass[lang=cn,scheme=chinese,openany]{elegantbook}
\usepackage{mathtools,bm,algorithm,algpseudocode,longtable}
\usepackage{cleveref,xurl}
\makeatletter
\providecommand{\theHALG@line}{}
\renewcommand{\theHALG@line}{\thealgorithm.\arabic{ALG@line}}
\makeatother
\newcommand{\E}{\mathbb E}
\newcommand{\KL}{\operatorname{KL}}
\newcommand{\Var}{\operatorname{Var}}
\newcommand{\Cov}{\operatorname{Cov}}
\newcommand{\TV}{\operatorname{TV}}
\newcommand{\osc}{\operatorname{osc}}
\newcommand{\sg}{\operatorname{stopgrad}}
\newcommand{\diag}{\operatorname{diag}}
\newcommand{\R}{\mathbb R}
\newcommand{\JS}{\operatorname{JS}}
\DeclareMathOperator*{\argmin}{arg\,min}
\DeclareMathOperator*{\argmax}{arg\,max}
\ifdefined\CoreOnly
  \externaldocument[full-]{2026-10-15-MirrorOPD核心推导}[2026-10-15-MirrorOPD核心推导.pdf]
  \newcommand{\appref}[1]{完整稿附录~\ref*{full-#1}}
\else
  \newcommand{\appref}[1]{附录~\ref{#1}}
\fi
\addbibresource{references.bib}
\hypersetup{pdftitle={MirrorOPD 核心推导：目标教师、内层拟合与参数性质},
  pdfauthor={MirrorOPD Research Notes},bookmarksnumbered=true}
\setcounter{tocdepth}{1}
\setcounter{secnumdepth}{2}
\allowdisplaybreaks[2]
\title{MirrorOPD 核心推导}
\author{MirrorOPD Research Notes}
\date{2026 年 10 月 15 日}
\begin{document}
\hypersetup{pageanchor=false}
\begin{titlepage}
\centering
\vspace*{28mm}
{\Huge\bfseries MirrorOPD 核心推导\par}
\vspace{12mm}
{\Large 目标教师、内层拟合与参数性质\par}
\vspace{12mm}
{\color{structurecolor}\rule{0.78\textwidth}{1pt}\par}
\vspace{12mm}
{\large 从加权散度最优化到 $(\kappa,\eta)$ 镜像更新\par}
\vspace{7mm}
\ifdefined\CoreOnly
{\large 正文阅读版；证明扩展见同源完整稿\par}
\else
{\large 正文与八组技术附录\par}
\fi
\vfill
{\large MirrorOPD Research Notes\par}
\vspace{5mm}
{\large 2026 年 10 月 15 日\par}
\vspace{18mm}
\end{titlepage}
\frontmatter
\hypersetup{pageanchor=true}
\chapter*{阅读约定}
本文从每轮外层的 Mirror Descent 子问题出发，依次定义冻结信号 $f$ 与 advantage，求出加权教师 $m$ 的显式表达式，再反推产生该镜像方向的教师目标 $J$。前三章完成这条推导；第四章介绍 inner 如何拟合已经确定的 $m$，第五、六章分析参数性质与有限拟合。附录对每项独立结论先陈述命题与条件，再给完整证明。

镜像下降（Mirror Descent，MD）在概率空间通过局部线性目标与 Bregman 正则定义更新。本文沿用来源中的 MDPO、BPO、PMD-Part、PMD-Mean 名称，依次表示镜像目标优化、中心化残差回归、完整最优性方程回归和均值近似回归；第~\ref{ch:inner} 章给出它们的精确定义。

精确分布更新、总体拟合目标、有限样本估计和神经网络参数更新是四个层次。定理在哪一层成立，就在哪一层表述。完整词表求和消除条件动作积分误差，仍保留前缀覆盖与模型表示误差。本文的解析结论不直接给出任务正确率排序。

\begin{center}
\begin{tabular}{lp{100mm}}
\toprule
记号 & 固定含义\\
\midrule
$T,q_k,p_{k,n}$ & 教师、第 $k$ 轮冻结锚点、第 $n$ 步内层学生\\
$m_k$ & 该轮由最优化问题定义的目标教师\\
$\kappa,\eta$ & 目标形状、教师插值权重；正文取 $0\le\kappa\le1,\ 0<\eta<1$\\
$\alpha_{\rm opt},\tau$ & 参数优化器学习率、固定 $J_{\kappa,\eta}$ 后另选的镜像步长\\
$t,t_{\rm in}$ & 内层梯度流时间、每轮固定内层运行时长\\
$F_k,J_{\kappa,\eta},L_k$ & 重心目标、教师目标、内层拟合损失\\
$f,A=f/\eta,K_T$ & 冻结信号、标准化教师信号、冻结教师中心\\
$\mathcal A_k=C_qA$ & 本轮零条件均值的 advantage；$g_k=-\mathcal A_k$\\
$C_qh=h-\E_qh$ & 按同一前缀的锚点分布中心化\\
$d_k,\widehat d_k$ & 总体、经验前缀访问测度\\
$K_{\rm head},M,N$ & 头部动作数、同前缀辅助动作数、内层步数\\
\bottomrule
\end{tabular}
\end{center}

所有对数均为自然对数。$\TV(P,Q)=\frac12\sum_a|P_a-Q_a|$；平方 Hellinger 距离固定为 $H^2(P,Q)=\sum_a(\sqrt{P_a}-\sqrt{Q_a})^2$，不带 $1/2$。向量方差表示协方差矩阵的迹。正文省略前缀条件时，所有分布仍条件于同一个前缀。

附录按数学依赖组织：A 给出测度和基础工具；B 完整识别教师目标；C 处理内层投影和优化；D 分析形状与矩；E 处理积分信息与估计；F 将中心估计接回拟合；G 分析外层误差与序列边界；H 比较 TOP-D 的目标教师。短阅读版与完整稿使用同一份正文源码。

\tableofcontents
\mainmatter
\chapter{加权最优化与目标教师}
\label{ch:target}
\section{每轮外层的最优化问题}
固定一个前缀，词表 $\mathcal V$ 有限，$|\mathcal V|\ge2$。教师 $T$ 与本轮冻结学生 $q=q_k$ 都严格为正且归一化。先允许每个前缀的候选分布独立变化；参数化约束在第~\ref{ch:inner} 章引入。令 $\Delta$ 为闭概率单纯形，$\Delta^\circ$ 为其内部。

采用如下幂散度约定\footfullcite{cichocki2010}：
\begin{equation}
D_\kappa(P\Vert Q)=
\frac{1-\sum_aP_a^\kappa Q_a^{1-\kappa}}{\kappa(1-\kappa)},
\qquad 0<\kappa<1.
\label{eq:power-divergence}
\end{equation}
两个连续端点为 $D_0(P\Vert Q)=\KL(Q\Vert P)$、$D_1(P\Vert Q)=\KL(P\Vert Q)$。因此以教师为第一自变量时，$\kappa=0$ 给出学生逆向 KL，$\kappa=1$ 给出学生正向 KL。

从一开始引入形状 $\kappa\in[0,1]$ 和教师权重 $\eta\in(0,1)$：
\begin{equation}
\boxed{m_k=\argmin_{p\in\Delta}F_k(p),\qquad
F_k(p)=(1-\eta)D_\kappa(q_k\Vert p)+\eta D_\kappa(T\Vert p).}
\label{eq:barycenter-problem}
\end{equation}
这个问题同时惩罚偏离锚点和偏离教师。$\eta$ 指定两者的权重，$\kappa$ 指定散度形状。所得 $m_k$ 称为本轮目标教师；实际回复仍可由 $q_k$ 生成。

\section{唯一解与拟合误差}
\begin{theorem}[归一化幂均值目标]
\label{thm:target}
对 $0<\kappa\le1$，定义
\begin{equation}
h_a=[(1-\eta)q_a^\kappa+\eta T_a^\kappa]^{1/\kappa},
\qquad Z=\sum_ah_a,\qquad m_a=h_a/Z.
\label{eq:target}
\end{equation}
它是式~\eqref{eq:barycenter-problem} 的唯一解。$\kappa=0$ 的连续端点为 $m_a=q_a^{1-\eta}T_a^\eta/Z$。当 $0<\kappa<1$ 时，
\begin{equation}
F(p)-F(m)=Z^\kappa D_\kappa(m\Vert p).
\label{eq:barycenter-gap}
\end{equation}
端点 $\kappa=0,1$ 的超额损失分别为 $\KL(p\Vert m)$、$\KL(m\Vert p)$。
\end{theorem}
\begin{proof}
令 $S_a=(1-\eta)q_a^\kappa+\eta T_a^\kappa$。对内点 $p$，
\begin{equation}
F(p)=\frac{1-\sum_aS_ap_a^{1-\kappa}}{\kappa(1-\kappa)},\quad
\partial_{p_a}F=-\frac{S_a}{\kappa}p_a^{-\kappa},\quad
\partial^2_{p_a}F=S_ap_a^{-\kappa-1}>0.
\end{equation}
加入约束乘子 $\nu$，一阶条件为 $S_ap_a^{-\kappa}=\kappa\nu$，归一化即得 $p=m$。为同时覆盖边界并证明全局唯一性，利用 $S_a=Z^\kappa m_a^\kappa$ 直接写出
\begin{equation}
F(p)=\frac{1-Z^\kappa}{\kappa(1-\kappa)}
+Z^\kappa D_\kappa(m\Vert p).
\end{equation}
Hölder 不等式给出 $\sum_am_a^\kappa p_a^{1-\kappa}\le1$，等号要求两个归一化向量相同；边界上也只有 $p=m$ 可以取等。因此得到唯一解与超额损失。

在 $\kappa=0$ 时，展开加权 KL 得 $F(p)=\KL(p\Vert m)-\log Z$。在 $\kappa=1$ 时，令 $m=(1-\eta)q+\eta T$，逐项相减得
\begin{equation}
F(p)-F(m)=\sum_am_a\log(m_a/p_a).
\end{equation}
两端的 KL 非负性给出相同的唯一性结论。最后对 $q_a^\kappa,T_a^\kappa$ 在零处展开，得到 $\log h_a\to(1-\eta)\log q_a+\eta\log T_a$，证明连续端点。
\end{proof}

幂均值重心属于已有的 $\alpha$ 散度模型汇合框架\footfullcite{amari2007}。本文使用它组织外层目标、内层投影与误差分析。最优化问题与拟合误差的精确关系是式~\eqref{eq:barycenter-gap}，不只是一条最优点相同的结论。

\section{三个代表形状}
\begin{equation}
\begin{aligned}
m_{0,\eta}&\propto q^{1-\eta}T^\eta,\\
m_{1/2,\eta}&\propto[(1-\eta)\sqrt q+\eta\sqrt T]^2,\\
m_{1,\eta}&=(1-\eta)q+\eta T.
\end{aligned}
\label{eq:three-targets}
\end{equation}
几何端点在 log-odds 坐标中走直线；算术端点在概率坐标中走直线；平方根点在平方根概率单位球面上走短弧。$\eta$ 对平方根点通常不是弧长比例。所有形状在 $\eta=0,1$ 分别到达 $q,T$，端点通过连续延拓理解。

只有算术端点的归一化常数恒为一。一般情况下不能删除 $Z$，也不能把幂均值目标当成概率的算术加权。其几何、覆盖界与多教师推广见\appref{app:objective}、\appref{app:shape}。

\chapter{教师目标与精确镜像更新}
\label{ch:objective}
\section{由目标分布识别教师信号}
为解释外层在优化哪个教师目标，定义概率比函数
\begin{equation}
A_{\kappa,\eta}(x)=
\begin{cases}
\dfrac{\log(1-\eta+\eta x^\kappa)}{\kappa\eta},&\kappa>0,\\[1mm]
\log x,&\kappa=0,
\end{cases}
\quad f_a=\eta A_{\kappa,\eta}(T_a/q_a).
\label{eq:signal}
\end{equation}
则 $q_ae^{f_a}=h_a$，目标教师写为 $m_a=q_ae^{f_a}/Z$。整个内层中 $q,T,f,Z$ 冻结。

本轮重心目标 $F_k$ 的最小点是 $m_k$。下面定义一个不依赖锚点、最小点为原教师 $T$ 的目标。两个函数的作用不同。

\begin{theorem}[教师目标的构造与镜像解释]
\label{thm:J}
定义
\begin{equation}
\omega_{\kappa,\eta}(x)=x\int_1^x\frac{A_{\kappa,\eta}(u)}{u^2}\,du,
\qquad
J_{\kappa,\eta}(p;T)=\sum_ap_a\omega_{\kappa,\eta}(T_a/p_a).
\label{eq:J}
\end{equation}
则 $J\ge0$，且在概率单纯形上唯一零点为 $p=T$。在正正交象限中，
\begin{equation}
\partial_{p_a}J(p;T)=-A_{\kappa,\eta}(T_a/p_a).
\label{eq:J-gradient}
\end{equation}
对 $J$ 在 $q$ 处执行步长 $\eta$ 的负熵镜像下降，恰好得到目标教师：
\begin{equation}
\boxed{m=\argmin_{p\in\Delta}
\left\{\langle\nabla J(q;T),p-q\rangle+\eta^{-1}\KL(p\Vert q)\right\}.}
\label{eq:MD}
\end{equation}
\end{theorem}
\begin{proof}
对积分求导，
\begin{align}
\omega'(x)&=\int_1^x A(u)u^{-2}\,du+A(x)/x,\\
\omega(x)-x\omega'(x)&=-A(x),\\
\omega''(x)&=A'(x)/x
=\frac{x^{\kappa-2}}{1-\eta+\eta x^\kappa}>0.
\end{align}
最后一式在 $\kappa=0$ 时取 $x^{-2}$。由 $A(1)=0$，有 $\omega(1)=\omega'(1)=0$，所以严格凸的 $\omega$ 只在一处取零。正权重求和给出 $J\ge0$，且 $J=0$ 要求 $T_a/p_a=1$ 对每个 $a$ 成立。链式法则给出
\begin{equation}
\partial_{p_a}[p_a\omega(T_a/p_a)]
=\omega(x)-x\omega'(x)=-A(x).
\end{equation}
式~\eqref{eq:MD} 的一阶条件为 $-A(T_a/q_a)+\eta^{-1}[\log(p_a/q_a)+1]+\nu=0$。消去公共乘子并归一化，即得 $p_a\propto q_ae^{\eta A(T_a/q_a)}$。KL 正则严格凸，故解唯一。
\end{proof}

\section{下降保证与参数角色}
负熵势记为 $\Phi(p)=\sum_ap_a\log p_a$。令 $x_a=T_a/p_a$，则
\begin{equation}
\nabla^2J(p;T)
=\diag\left(\frac1{p_a}\frac{x_a^\kappa}{1-\eta+\eta x_a^\kappa}\right)
\preceq\eta^{-1}\nabla^2\Phi(p).
\label{eq:relative-smooth}
\end{equation}
沿线段积分 Hessian，得到相对光滑不等式
\begin{equation}
J(p)\le J(q)+\langle\nabla J(q),p-q\rangle+\eta^{-1}\KL(p\Vert q).
\end{equation}
镜像最优性条件与 $m-q$ 作内积，给出
\begin{equation}
\langle\nabla J(q),m-q\rangle
=-\eta^{-1}[\KL(m\Vert q)+\KL(q\Vert m)].
\end{equation}
所以 $J(q)-J(m)\ge\eta^{-1}\KL(q\Vert m)$，非平凡精确更新严格下降。对同一条件分布反复精确刷新，三点恒等式进一步给出
\begin{equation}
J(q_K;T)\le\frac{\KL(T\Vert q_0)}{\eta K}.
\label{eq:J-rate}
\end{equation}
完整望远镜求和、边界排除和可变镜像步长证明见\appref{app:objective}。

在整个递推中固定 $\kappa,\eta,T$，$J_{\kappa,\eta}$ 才是同一个函数。改变 $\eta$ 一般同时改变目标形状和镜像步长。若要保持 $J$ 不变而调整步长，应另用 $\tau$，更新 $p^+\propto q\exp[\tau A_{\kappa,\eta}(T/q)]$。参数优化器的学习率 $\alpha_{\rm opt}$ 又是第三个量。

\section{散度、教师目标与中心}
三个容易混用的对象分别是
\begin{equation}
D_\kappa(T\Vert p),\qquad J_{\kappa,\eta}(p;T),\qquad
K_T(q,T)=-\E_q A_{\kappa,\eta}(T/q).
\end{equation}
它们依次用于定义重心、解释精确镜像更新、中心化本轮监督。$K_T$ 在内层冻结。端点关系为
\begin{equation}
J_{0,\eta}(p;T)=\KL(p\Vert T),\quad
\lim_{\eta\downarrow0}J_{\kappa,\eta}(p;T)=D_\kappa(T\Vert p),\quad
\lim_{\eta\uparrow1}J_{\kappa,\eta}(p;T)=\KL(p\Vert T).
\end{equation}
当 $\kappa=1$、$0<\eta<1$，$J$ 是归一化加权 JS，而 $D_1(T\Vert p)$ 是正向 KL。有限步幂均值一般不是对固定 $D_\kappa$ 做同一步长的熵镜像下降；两者在小步极限才一阶一致。

\chapter{固定目标教师的内层拟合}
\label{ch:inner}
\section{最优性方程与四种损失}
本轮固定 $q,T,m$ 与 $f=\eta A(T/q)$，当前学生为 $p=p_\theta$。记
\begin{equation}
U_a=\log(p_a/q_a),\qquad G=C_qf,\qquad
K_S=\KL(q\Vert p)=-\E_qU,\qquad K_T=-\E_qA.
\end{equation}
目标教师满足 $U^\star=f-\log Z$。将该方程转成可训练目标，有以下四种方式：
\begin{align}
L_{\rm MDPO}(p)&=\frac1\eta\{\KL(p\Vert q)-\E_pf\}
=\frac1\eta\KL(p\Vert m)-\frac{\log Z}{\eta},\label{eq:mdpo}\\
L_{\rm Part}(p)&=\frac12\E_q(U-f+\log Z)^2,\label{eq:part}\\
L_{\rm BPO}(p)&=\frac1{2\eta}\E_q[C_q(U-f)]^2
=\frac1{2\eta}\Var_q\log(p/m),\label{eq:bpo}\\
L_{\rm Mean}(p)&=\frac12\E_q(U-G)^2.\label{eq:mean}
\end{align}
MDPO 保留原镜像子问题；Part 回归完整的对数比方程；BPO 回归中心化后的方程；Mean 用 $\E_qf$ 近似 $\log Z$。Mean 的定义与 PMD-mean 原始分析一致，整体损失尺度在此统一\footfullcite{pmd2026}。

\begin{proposition}[共同目标与 Mean 的偏离]
\label{prop:inner-roots}
在完整支持、精确期望和可实现的分布类中，MDPO、Part、BPO 的唯一最优分布均为 $m$。若 $f$ 非恒定，Mean 的最优分布一般不等于 $m$。
\end{proposition}
\begin{proof}
MDPO 的结论由式~\eqref{eq:mdpo} 的 KL 分解给出。Part 的零损失要求 $\log(p/m)=0$。BPO 的零损失要求 $\log(p/m)$ 在全部动作上为常数，归一化再强制常数为零。

Mean 满足约束 $\E_qe^U=1$，其拉格朗日一阶条件是
\begin{equation}
U_a-G_a+\lambda e^{U_a}=0.
\label{eq:mean-kkt-main}
\end{equation}
Gibbs 解满足 $U_a=G_a-\Delta$，其中 $\Delta=\log\E_qe^G>0$。代入后要求 $\lambda e^{G_a-\Delta}=\Delta$ 对所有 $a$ 成立；非恒定 $G$ 使其不可能。Mean 的存在唯一性及 Lambert-$W$ 解见\appref{app:inner}。
\end{proof}

Mean 与 BPO 的差异还可精确写成
\begin{equation}
L_{\rm Mean}=\eta L_{\rm BPO}+\frac12K_S^2.
\label{eq:mean-bpo}
\end{equation}
因为 $U-G=C_q(U-f)-K_S$，而第一项的 $q$ 均值为零，平方交叉项消失。额外的 $K_S^2/2$ 改变最优分布与梯度；它不能当成不影响优化的常数。

\section{精确梯度与采样测度}
令 $\psi_a=\nabla_\theta\log p_\theta(a)$、$\bar\psi_q=\E_q\psi$，则 $\E_p\psi=0$，但 $\bar\psi_q$ 通常不为零。定义
\begin{equation}
R_P=U-f+\log Z,\quad R_M=U-G,\quad
R_B=C_q(U-f)=U+K_S-\eta(A+K_T).
\end{equation}
直接求导得到
\begin{align}
\nabla L_{\rm Part}&=\E_q[R_P\psi],&
\nabla L_{\rm Mean}&=\E_q[R_M\psi],\\
\nabla L_{\rm BPO}
&=\eta^{-1}\E_q[R_B(\psi-\bar\psi_q)]
=\eta^{-1}\E_q[R_B\psi],\label{eq:bpo-gradient}\\
\nabla L_{\rm MDPO}
&=\eta^{-1}\E_p[(U-f)\psi]
=\eta^{-1}\E_q[(p/q)(U-f)\psi].\label{eq:mdpo-gradient}
\end{align}
BPO 第二个等号使用 $\E_qR_B=0$；MDPO 删除常数项使用 $\E_p\psi=0$。这两次抵消发生在不同测度下。

完整词表求和时，可直接对真实损失求导。使用解析 score 形式时，把整个系数停止梯度，再乘以 $\log p_\theta(a)$，实现上述向量。这个代理标量的数值不等于真实损失。Part、Mean、BPO 的动作测度是 $q$，MDPO 的动作测度是 $p$；从 $q$ 采样 MDPO 时需乘 $p/q$。

在 $p=q$ 处，BPO 和 MDPO 的梯度相同，Part、Mean 的梯度等于它们的 $\eta$ 倍。离开锚点后，损失的度量、中心与测度不同，不能用起点相等推断整个轨迹相等。

\section{接入前缀与外层同步}
恢复条件变量 $s$，内层实际优化
\begin{equation}
\mathcal L_k(\theta)=\E_{s\sim d_k}L_k(p_\theta(\cdot\mid s)).
\end{equation}
$d_k$ 在本轮冻结；有限 rollout 用经验测度 $\widehat d_k$ 代替。教师信号 $f,K_T$ 整轮冻结，学生中心 $K_S$ 随每次参数更新重算。每个前缀的一条实际动作观测，不等于该前缀全部动作的期望；辅助动作抽样与新增完整回复是两种不同操作。

\begin{algorithm}[ht]
\caption{固定前缀的 MirrorOPD\\{\small\textit{Semantics: 根据条件目标重构；动作期望使用完整词表}}}
\begin{algorithmic}[1]
\Require teacher $T$, initial student $\theta_0$, $\kappa\in[0,1]$, $0<\eta<1$, finite budgets $K,N$
\For{$k=0,\ldots,K-1$}
    \State $q_k\gets p_{\theta_k}$; freeze $q_k$
    \State sample prompts and one response per prompt from $q_k$ with full action support
    \State form $\widehat d_k$ from valid response tokens, including EOS
    \If{the batch has no valid response tokens}
        \State terminate with an invalid-batch error
    \EndIf
    \State cache frozen $f_k=\sg(\eta A_{\kappa,\eta}(T/q_k))$ and required teacher integrals
    \State $\theta_{k,0}\gets\theta_k$
    \For{$n=0,\ldots,N-1$}
        \State recompute $p_{\theta_{k,n}}$ and all student-dependent integrals
        \State $g_{k,n}\gets\nabla_\theta\E_{\widehat d_k}L_k(p_{\theta_{k,n}})$
        \State $\theta_{k,n+1}\gets\theta_{k,n}-\alpha_{k,n}g_{k,n}$
    \EndFor
    \State $\theta_{k+1}\gets\theta_{k,N}$; discard the old rollout for the next outer
\EndFor
\State \Return $\theta_K$
\end{algorithmic}
\end{algorithm}

算法以普通梯度下降明确更新语义。替换为其他优化器时，目标定义保持不变，数值轨迹需另行分析。共享参数可能无法同时表示所有 $m_k(\cdot\mid s)$；此时三种具有相同自由分布解的方法，也可能具有不同的受限投影。

\chapter{形状与插值强度的性质}
\label{ch:parameters}
\section{路径、覆盖与局部速度}
固定 $q,T$ 时，$\eta$ 从零增加到一，把目标教师从锚点移向教师。对固定 $\kappa$，任意凸生成散度 $D_\varphi(m_{\kappa,\eta}\Vert q)$ 随 $\eta$ 非减；不同 $\kappa$ 下，相同 $\eta$ 不代表相同距离。其交点证明见\appref{app:shape}。

正形状有一个分布无关的覆盖保证：
\begin{equation}
m_a\ge(1-\eta)^{1/\kappa}q_a+\eta^{1/\kappa}T_a,\qquad 0<\kappa\le1.
\label{eq:coverage-main}
\end{equation}
证明使用 $h_a\le(1-\eta)q_a+\eta T_a$，故 $Z\le1$；以及 $(x+y)^{1/\kappa}\ge x^{1/\kappa}+y^{1/\kappa}$。于是 $m_a=h_a/Z\ge h_a$ 给出下界。几何端点没有统一的正教师覆盖系数。

对固定严格正分布，小步展开为
\begin{equation}
m_a=q_a+\eta v_{\kappa,a}+O(\eta^2),\qquad
v_{\kappa,a}=\frac{q_a}{\kappa}
\left[(T_a/q_a)^\kappa-\E_q(T/q)^\kappa\right].
\label{eq:velocity-main}
\end{equation}
$\kappa=0$ 时取连续极限 $v_{0,a}=q_aC_q\log(T_a/q_a)$。这是 $D_\kappa(T\Vert q)$ 在 Fisher 度量下的负自然梯度。形状改变高维方向，不能通常由一个学习率缩放替代。

\section{监督尾部与方差阈值}
当 $\kappa>0,\eta<1$ 时，
\begin{equation}
A_{\kappa,\eta}(x)\ge-c_{\kappa,\eta},\qquad
c_{\kappa,\eta}=-\frac{\log(1-\eta)}{\kappa\eta}.
\end{equation}
正尾部仍无逐点统一上界，但在 $q$ 下利用 $\E_q(T/q)=1$ 可以控制矩。令 $V_{\kappa,\eta}=\Var_q A_{\kappa,\eta}(T/q)$。对每个固定正 $\eta<1$ 和 $\kappa>0$，$V$ 在所有有限正概率分布上有统一有限上界。

更强的小步问题需要另一个结论：
\begin{equation}
\sup_{q,T}V_{\kappa,\eta}
=\Theta_\kappa\!\left(\eta^{\min\{2,1/\kappa\}-2}\right),
\qquad \eta\downarrow0,\quad\kappa>0.
\label{eq:variance-threshold-main}
\end{equation}
所以保持统一小步方差有界的区间为 $0<\kappa\le1/2$。几何端点对任意固定正 $\eta$ 都没有分布无关界。这里上确界允许分布随 $\eta$ 改变；固定一对正分布时，每种形状的小步极限都有限。

对 $0<\kappa\le1/2$，还有误差相关界
\begin{equation}
V_{\kappa,\eta}
\le\frac{D_\kappa(T\Vert q)}{\kappa(1-\eta)^2}
\le\frac{K_T}{\kappa(1-\kappa)(1-\eta)^2}.
\label{eq:error-variance-main}
\end{equation}
若 score 范数统一不超过 $S$，中心化教师梯度的协方差迹不超过 $S^2V$。标量监督方差与完整参数梯度方差由这个额外条件连接；精确均值中心本身不保证 score 梯度的方差最小。

\section{平方根点的作用与限制}
令 $a=\sum_v\sqrt{q_vT_v}$、$w=1-\eta$。平方根目标满足
\begin{equation}
\sqrt m=\frac{w\sqrt q+\eta\sqrt T}{\sqrt Z},\qquad
Z=w^2+\eta^2+2w\eta a.
\end{equation}
因为 $Z\le1$，
\begin{equation}
\sum_v\sqrt{m_vT_v}=\frac{wa+\eta}{\sqrt Z}\ge wa+\eta,
\quad
H^2(m,T)\le(1-\eta)H^2(q,T).
\label{eq:hellinger-contraction-main}
\end{equation}
这是精确分布更新的全局收缩。所有形状在固定正教师附近却共享
\begin{equation}
m_{\kappa,\eta}(T+\delta,T)
=T+(1-\eta)\delta+O(\|\delta\|^2).
\label{eq:local-rate-main}
\end{equation}
因此平方根的全局 Hellinger 保证不代表更好的局部一阶概率收缩率。

平方根点唯一最小化小步监督方差相对于自身教师散度的最坏比值；若改为最小化绝对最坏方差，最优形状约为 $0.441928$。两种准则的推导在\appref{app:shape}。固定 $q,\eta$ 时，BPO 的独立 logits Hessian 为
\begin{equation}
\nabla_z^2L_{\rm BPO}=\eta^{-1}[\diag(q)-qq^\top],
\end{equation}
其中不含 $\kappa$。形状改变目标和初始残差，不能由此宣称改善该二次问题的条件数。

\chapter{有限拟合与外层进展}
\label{ch:finite}
\section{从本轮拟合到教师接近度}
\begin{theorem}[重心进展恒等式]
\label{thm:progress}
对任意候选 $p$，令 $c_\kappa=Z^\kappa$，端点 $\kappa=0,1$ 取 $c_\kappa=1$，则
\begin{align}
\eta[D_\kappa(T\Vert q)-D_\kappa(T\Vert p)]
={}&(1-\eta)D_\kappa(q\Vert p)\nonumber\\
&+c_\kappa[D_\kappa(m\Vert q)-D_\kappa(m\Vert p)].
\label{eq:progress}
\end{align}
因此 $D_\kappa(m\Vert p)\le D_\kappa(m\Vert q)$ 是教师散度非增的充分条件。精确更新 $p=m$ 在 $q\ne T$ 时严格改善。
\end{theorem}
\begin{proof}
重心定义给出 $F(q)=\eta D_\kappa(T\Vert q)$。另一方面，定理~\ref{thm:target} 给出 $F(q)-F(p)=c_\kappa[D_\kappa(m\Vert q)-D_\kappa(m\Vert p)]$。代入 $F(p)$ 并移项得到恒等式。条件成立时右侧两项非负。若 $p=m$ 且 $q\ne T$，则 $m\ne q$，右侧第一项严格为正。
\end{proof}

精确幂均值更新还有更强性质：它同时降低所有凸生成的教师 $f$-散度。证明使用概率比的共同交点，见\appref{app:shape}。有限 inner 降低某个拟合损失却不足以自动继承这一性质。

\section{拟合误差证书}
MDPO 的超额损失 $\delta$ 精确对应 $\KL(p\Vert m)=\eta\delta$，由 Pinsker 不等式得到 $\TV(p,m)\le\sqrt{\eta\delta/2}$。BPO 控制的是按 $q$ 加权的对数比方差。若另有 $q_a\ge q_{\min}>0$，则
\begin{equation}
d_H(p,m):=\osc_a\log(p_a/m_a)
\le2\sqrt{\eta L_{\rm BPO}/q_{\min}},
\quad
\TV(p,m)\le\tanh[d_H(p,m)/4].
\label{eq:bpo-certificate-main}
\end{equation}
最小锚点概率极小时，这个证书可能很弱；推导与锐 TV 界见\appref{app:measure}、\appref{app:dynamics}。

这种限制确实会发生。取
\begin{equation}
q=(0.4999995,0.4999995,10^{-6}),\quad
T=(0.8,0.199999,10^{-6}),\quad
\kappa=\tfrac12,\quad\eta=0.3,
\end{equation}
并令 $p\propto m\exp(0,0,12)$。直接求和得到
\begin{equation}
L_{\rm BPO}(q)\simeq0.062277,\quad L_{\rm BPO}(p)\simeq0.000240,
\qquad
\KL(q\Vert T)\simeq0.223145,\quad\KL(p\Vert T)\simeq1.654219.
\end{equation}
稀有动作的对数误差在 $q$ 下权重很小，却可以改变归一化后的学生质量。因此需要区分 loss 下降与目标教师拟合证书。

\section{有效保证的条件}
从 $q$ 出发，使用精确锚点 Fisher 预条件的 BPO 产生理想路径
\begin{equation}
p_\xi\propto q\exp(\xi f),\qquad 0\le\xi\le1.
\end{equation}
沿这条路径，所有凸生成的教师散度非增。该结论依赖路径、完整期望与自由 logits；普通神经网络优化器不自动位于这条路径上，证明见\appref{app:inner}。

固定教师的精确外层收敛，允许通过输出误差递推扩展到近似拟合。应先由内层损失、中心误差和表示限制建立输出误差，再分析多轮累积；梯度估计的方差不能直接替代输出误差。

最后，正文的目标是固定前缀测度下的条件目标。完整序列 KL 使用随策略变化的访问权重；条件 KL 逐前缀下降，也可能同时让学生更频繁进入误差较大的分支。序列反例、后缀价值修正和同步边界见\appref{app:dynamics}。TOP-D 的目标教师由算术端点精确恢复，其收缩与方差性质见\appref{app:topd}。

\ifdefined\CoreOnly\else
\appendix
\chapter{测度、归一化与基础工具}
\label{app:measure}
\section{变长轨迹如何诱导条件目标}
设提示 $x\sim\mathcal D$，锚点 $q$ 生成轨迹，回复长度 $H$ 包含 EOS，不包含提示与 padding。预算终止时保留已生成动作。假设 $0<\E H<\infty$，定义
\begin{equation}
d_q(B)=\frac{\E_q\sum_{t\ge1}\mathbf1\{t\le H,\ s_t\in B\}}{\E_qH}.
\end{equation}
\begin{proposition}[访问测度的条件分解]
在上述停止规则及 $0<\E_qH<\infty$ 下，总体前缀--动作测度为 $d_q(s)q(a\mid s)$。对满足 $\E_q\sum_{t\le H}|h(s_t,a_t)|<\infty$ 的函数，
\begin{equation}
\E_q\sum_{t\le H}h(s_t,a_t)=\E_qH\;\E_{d_q,q}h.
\end{equation}
\end{proposition}
\begin{proof}
这里 $s_t$ 是抽取第 $t$ 个动作之前的历史。存活事件 $\{t\le H\}$ 关于这段历史可测。对于绝对可积的 $h$，塔式期望给出
\begin{align}
\E_q\sum_{t\ge1}\mathbf1\{t\le H\}h(s_t,a_t)
&=\E_q\sum_{t\ge1}\mathbf1\{t\le H\}
\sum_aq(a\mid s_t)h(s_t,a)\\
&=\E_qH\ \E_{s\sim d_q,a\sim q(\cdot\mid s)}h(s,a).
\end{align}
第一步交换求和与期望由非负性或绝对可积性保证。因此总体前缀--动作测度是 $d_qq$。
\end{proof}

\begin{proposition}[全局 token 比率估计的一致性]
独立同分布轨迹具有有限正平均长度及可积的轨迹总损失时，全局 token 均值
$\widehat I_B=(\sum_{i,t\le H_i}h_{it})/(\sum_iH_i)$ 几乎处处收敛到 $\E_{d_q,q}h$。该结论不要求轨迹内 token 独立，也不保证有限样本无偏。
\end{proposition}
\begin{proof}
对 $B$ 条独立轨迹，全局 token 均值为
\begin{equation}
\widehat I_B=\frac{\sum_{i=1}^B\sum_{t=1}^{H_i}h(s_{it},a_{it})}{\sum_{i=1}^BH_i}.
\end{equation}
分子分母除以 $B$，分别由大数定律趋于对应期望；分母极限严格为正，所以比值一致收敛到 $\E_{d_q,q}h$。有限 $B$ 时随机分母一般产生比率估计偏差。同一轨迹内的 token 相关，置信区间不能把全部 token 当成独立样本。
\end{proof}

\begin{proposition}[逐序列长度权重改变动作分布]
存在有限长度生成过程，使按每条轨迹的 $1/H$ 加权后，同一前缀下的动作分布不同于生成策略。因此逐序列均值不能一般替代 $d_qq$ 下的 token 目标。
\end{proposition}
\begin{proof}
逐序列先除以 $H_i$ 会改变测度。取初始两动作各以 $1/2$ 发生，动作一立即终止，动作二随后再生成一个确定 EOS。按 $1/H$ 加权后，初始两动作质量为 $1/2,1/4$，条件概率成为 $2/3,1/3$。这个权重依赖当前与未来动作，不能移出条件期望。
\end{proof}

\begin{proposition}[微批次梯度累加不变性]
固定有效 token 总数 $\mathsf T>0$，微批次构成有效位置的不交划分。若每批使用总分母 $\mathsf T$，并在全部梯度累加后进行裁剪和更新，则所得梯度等于全 batch 目标的梯度。
\end{proposition}
\begin{proof}
固定 batch 后，总有效 token 数 $\mathsf T$ 不依赖内层参数。若微批次互不相交，使用 $L_m=\mathsf T^{-1}\sum_{i\in\mathcal B_m}\ell_i$，则
\begin{equation}
\sum_m\nabla L_m=\mathsf T^{-1}\sum_i\nabla\ell_i.
\end{equation}
每个微批次各自取均值时，需再乘 $|\mathcal B_m|/\mathsf T$。裁剪和优化器更新放在梯度累加以后，才保留这条线性恒等式。
\end{proof}

\section{积分微分与支持条件}
默认有限词表上的概率严格为正。若允许零概率，$\KL(P\Vert Q)$ 有限至少要求 $P\ll Q$；重要性抽样则需要提案覆盖非零被积贡献。两项支持要求由不同对象决定。

\begin{lemma}[score 恒等式与积分微分]
有限词表上，若 $p_\theta>0$ 归一化且 $p_\theta,h_\theta$ 可微，则
\begin{equation}
\E_p\psi=0,\qquad
\nabla_\theta\E_ph_\theta=\E_p(\nabla_\theta h_\theta+h_\theta\psi).
\end{equation}
若积分权重 $q$ 冻结，则 $\nabla_\theta\E_qh_\theta=\E_q\nabla_\theta h_\theta$。
\end{lemma}
\begin{proof}
令 $p_\theta$ 可微，$\psi_a=\nabla_\theta\log p_\theta(a)$。归一化给出
\begin{equation}
\E_p\psi=\sum_a\nabla p_\theta(a)=0.
\end{equation}
对参数依赖函数 $h_\theta$，
\begin{equation}
\nabla_\theta\E_{p_\theta}h_\theta
=\E_p[\nabla_\theta h_\theta+h_\theta\psi].
\label{eq:measure-derivative}
\end{equation}
固定采样索引的自动微分只产生第一项。冻结测度 $q$ 下则有 $\nabla\E_qh_\theta=\E_q\nabla h_\theta$，且 $\E_q\psi$ 一般不为零。有限词表使逐项微分合法；连续空间需另行提供可交换微分与积分的支配条件。
\end{proof}

\section{散度规范化与局部二阶项}
\begin{proposition}[幂散度的生成函数与对偶]
对 $0<\kappa<1$，幂散度的规范化生成函数为
\begin{equation}
g_\kappa(x)=\frac{\kappa x+1-\kappa-x^\kappa}{\kappa(1-\kappa)}.
\end{equation}
其曲率为 $x^{\kappa-2}$；连续端点为 $g_0=x-1-\log x$、$g_1=x\log x-x+1$。交换分布时 $D_\kappa(P\Vert Q)=D_{1-\kappa}(Q\Vert P)$，且 $D_{1/2}=2H^2$。
\end{proposition}
\begin{proof}
对凸函数 $g$，定义 $D_g(P\Vert Q)=\sum_aQ_ag(P_a/Q_a)$。加入线性项 $b(x-1)$ 不改变散度。采用 $g(1)=g'(1)=0$ 的规范化后，
\begin{equation}
g_\kappa(x)=\frac{\kappa x+1-\kappa-x^\kappa}{\kappa(1-\kappa)},
\qquad g_\kappa''(x)=x^{\kappa-2}.
\end{equation}
$g_0(x)=x-1-\log x$、$g_1(x)=x\log x-x+1$ 由洛必达法则得到。对偶关系由逐点恒等式 $xg_{1-\kappa}(1/x)=g_\kappa(x)$ 给出：
\begin{equation}
D_\kappa(P\Vert Q)=D_{1-\kappa}(Q\Vert P).
\end{equation}
在 $\kappa=1/2$，$D_{1/2}(P\Vert Q)=4(1-\sum\sqrt{PQ})=2H^2(P,Q)$。
\end{proof}

\begin{proposition}[规范化散度的 Fisher 二阶项]
设 $Q>0$ 固定，$g$ 在一附近三阶连续可微，$g(1)=g'(1)=0,g''(1)=1$。若 $\sum_a\delta_a=0$ 且 $\max_a|\delta_a/Q_a|\to0$，则
\begin{equation}
D_g(Q+\delta\Vert Q)=\tfrac12\sum_a\delta_a^2/Q_a
+O\!\left(\sum_a|\delta_a|^3/Q_a^2\right).
\end{equation}
\end{proposition}
\begin{proof}
若 $g$ 在一附近三阶可微，$g''(1)=1$，并且 $P=Q+\delta$、$\sum_a\delta_a=0$、$\max_a|\delta_a/Q_a|\to0$，则
\begin{equation}
D_g(Q+\delta\Vert Q)
=\frac12\sum_a\frac{\delta_a^2}{Q_a}
+O\left(\sum_a\frac{|\delta_a|^3}{Q_a^2}\right).
\end{equation}
证明是对 $g(1+\delta_a/Q_a)$ 逐项 Taylor 展开；常数和一次项消失，余项乘以 $Q_a$ 后求和。这个共同 Fisher 二阶尺度不表示不同散度在远处相等。
\end{proof}

\section{中心化坐标与归一化恢复}
\begin{proposition}[中心化误差与概率分布的双射]
固定正分布 $q,m$。映射 $p\mapsto e=C_q\log(p/m)$ 将正单纯形双射到 $\{e:\E_qe=0\}$，逆映射为
\begin{equation}
p_a=m_a\exp[e_a-\log\E_me^e].
\end{equation}
\end{proposition}
\begin{proof}
固定 $q,m>0$。定义 $e=C_q\log(p/m)$，则 $\E_qe=0$。写 $\log(p/m)=e+c$，归一化给出 $e^c\E_me^e=1$，故
\begin{equation}
p_a=m_a e^{e_a-\Lambda(e)},\qquad
\Lambda(e)=\log\E_me^e.
\label{eq:quotient-inverse-core}
\end{equation}
反过来，任何 $\E_qe=0$ 的有限向量经该式得到正分布，并恢复原 $e$。因此归一化分布与 $q$ 中心化坐标一一对应。$\Lambda(e)$ 可以为负，因为中心化使用 $q$，配分函数使用 $m$。
\end{proof}

欧氏中心坐标 $P_\perp v$ 使用 $P_\perp=I-\mathbf1\mathbf1^\top/|\mathcal V|$，而 $C_qv$ 使用 $q$。两者相差常数，代表同一个归一化自由度；它们的欧氏范数及所在子空间不能直接混用。

\section{从对数误差到概率距离}
\begin{lemma}[加权误差控制跨度]
若 $q_i,q_j>0$，则
\begin{equation}
(v_i-v_j)^2\le(q_i^{-1}+q_j^{-1})\E_qv^2.
\end{equation}
特别地，若 $q_a\ge q_{\min}>0$，则 $(\osc v)^2\le2\E_qv^2/q_{\min}$。
\end{lemma}
\begin{proof}
将差写为 $(\sqrt{q_i}v_i,\sqrt{q_j}v_j)\cdot(q_i^{-1/2},-q_j^{-1/2})$，应用 Cauchy--Schwarz，再把两个加权平方放大为全部平方和。选择最大与最小动作即得跨度界。
\end{proof}

\begin{theorem}[Hilbert 距离到 TV 的锐界]
\label{thm:hilbert-tv}
对正概率分布，
\begin{equation}
\TV(p,q)\le\tanh\left(\frac{\osc\log(p/q)}4\right).
\end{equation}
任意给定正跨度都有二元分布取等。
\end{theorem}
\begin{proof}
令 $L=p/q$、$l=\min L$、$u=\max L$、$h=u/l=e^\delta$。$\E_qL=1$ 保证 $l\le1\le u$。凸函数 $|L-1|$ 不超过端点连线，所以
\begin{equation}
\TV(p,q)\le\frac{(u-1)(1-l)}{u-l}
=\frac{h+1-hl-l^{-1}}{h-1}
\le\frac{\sqrt h-1}{\sqrt h+1}.
\end{equation}
最后一步使用 $hl+l^{-1}\ge2\sqrt h$。取 $L$ 的两值为 $e^{\delta/2},e^{-\delta/2}$，在 $q$ 下的概率为 $1/(1+e^{\delta/2})$ 和 $e^{\delta/2}/(1+e^{\delta/2})$。它们满足均值一，且弦线及算术--几何均值不等式同时取等。
\end{proof}

\begin{theorem}[Pinsker 不等式]
对有限词表上的正概率分布，$\TV(p,q)\le\sqrt{\KL(p\Vert q)/2}$。
\end{theorem}
\begin{proof}
另一个常用界是 Pinsker 不等式 $\TV(p,q)\le\sqrt{\KL(p\Vert q)/2}$。为给出常数，取 $B=\{a:p_a\ge q_a\}$，则 $p(B)-q(B)=\TV(p,q)$。log-sum 不等式把 KL 下界为二元 KL。固定 $v=q(B)$，函数
\begin{equation}
b(u)=u\log(u/v)+(1-u)\log[(1-u)/(1-v)]
\end{equation}
满足 $b(v)=b'(v)=0$、$b''(u)=1/[u(1-u)]\ge4$。在 $v,u$ 之间积分两次，得到 $b(u)\ge2(u-v)^2$。非平凡正支持情形中 $0<v<1$；距离为零时直接成立。
\end{proof}

\chapter{重心、教师目标与镜像下降}
\label{app:objective}
\section{密度比矩、散度与几何重加权}
固定教师 $T$，当前学生 $p=p_\theta$，令 $\rho=T/p$，定义 $g(u)=\E_p\rho^u=\sum_aT_a^up_a^{1-u}$。有限正支持下可逐项求导：
\begin{equation}
g(0)=g(1)=1,\quad g'(u)=\E_p[\rho^u\log\rho],\quad
g''(u)=\E_p[\rho^u\log^2\rho]\ge0.
\end{equation}
两个端点斜率为 $g'(0)=-\KL(p\Vert T)$、$g'(1)=\KL(T\Vert p)$。因此 $D_\kappa(T\Vert p)=[1-g(\kappa)]/[\kappa(1-\kappa)]$ 在两个端点连续恢复对应 KL；严格凸性或 Hölder 给 $g(\kappa)<1$，除非 $p=T$。

对学生参数求导，$T$ 冻结，概率项 $p_a^{1-\kappa}$ 的导数为 $(1-\kappa)p_a^{1-\kappa}\psi_a$。所以
\begin{equation}
-\nabla_\theta D_\kappa(T\Vert p)
=\frac1\kappa\E_p[\rho^\kappa\psi]
=\E_p\left[\frac{\rho^\kappa-1}{\kappa}\psi\right].
\end{equation}
第二式使用 $\E_p\psi=0$。$\kappa\downarrow0$ 给逆向 KL 的 log-ratio 方向；$\kappa=1$ 给 $\E_T\psi$。令 $\pi_\kappa\propto T^\kappa p^{1-\kappa}$，相同梯度也写成 $g(\kappa)\E_{\pi_\kappa}\psi/\kappa$。这个几何重加权测度随当前学生变化，解释每次冻结测度后的局部最大似然方向。

在 $\kappa\downarrow0$，$\pi_\kappa=p[1+\kappa C_p\log(T/p)]+O(\kappa^2)$，所以 $\E_{\pi_\kappa}\psi$ 的零阶项为零，一阶项与外部 $1/\kappa$ 抵消。这里的 $\pi_\kappa$ 是梯度积分的测度，通常不等于重心问题的 $m_{\kappa,\eta}$；后者还包含独立的插值参数 $\eta$。

\section{加权重心的两个端点与多教师形式}
对 $\kappa=0$，令 $Z_0=\sum_aq_a^{1-\eta}T_a^\eta$。展开各对数可得
\begin{equation}
(1-\eta)\KL(p\Vert q)+\eta\KL(p\Vert T)
=\KL(p\Vert m_{0,\eta})-\log Z_0.
\end{equation}
对 $\kappa=1$，$m=(1-\eta)q+\eta T$，因此
\begin{align}
&(1-\eta)\KL(q\Vert p)+\eta\KL(T\Vert p)\\
&\quad=\KL(m\Vert p)+(1-\eta)\KL(q\Vert m)+\eta\KL(T\Vert m).
\end{align}
两式既确定最优分布，也确定有限误差应使用哪个 KL 方向。

对任意有限个正分布 $T^{(j)}$ 与非负权重 $w_j$、$\sum_jw_j=1$，令 $S_a=\sum_jw_j(T_a^{(j)})^\kappa$。正文证明逐项保持不变，得到
\begin{equation}
m_a=\frac{S_a^{1/\kappa}}{\sum_bS_b^{1/\kappa}},
\qquad
\sum_jw_jD_\kappa(T^{(j)}\Vert p)
-\sum_jw_jD_\kappa(T^{(j)}\Vert m)
=Z^\kappa D_\kappa(m\Vert p).
\end{equation}
几何端点取 $m\propto\prod_j(T^{(j)})^{w_j}$，算术端点取加权概率和。锚点可作为一个分布加入。多教师重心并不意味着每个教师距离都下降；各教师目标可能冲突。

对 $0<\kappa<1$，重心目标在最优点的概率 Hessian 为 $Z^\kappa\diag(1/m)$。若参数化模型恰好实现 $m$，令概率 Jacobian 为 $B_\theta=\partial p/\partial\theta$，则
\begin{equation}
\nabla_\theta^2F(p_\theta)\big|_{p=m}
=Z^\kappa B_\theta^\top\diag(1/m)B_\theta.
\end{equation}
链式法则的另一项为 $\sum_a(\partial F/\partial p_a)\nabla_\theta^2p_a$；最优点的概率梯度是常数向量，该项由 $\sum_ap_a=1$ 消失。离开可实现最优点时不能删除这项。

\section{逆向构造的唯一性与通用信号}
若要求 $D_\omega(T\Vert p)$ 的逐坐标梯度是 $-A(T/p)$，透视函数求导强制 $\omega-x\omega'=-A$。于是
\begin{equation}
(\omega(x)/x)'=A(x)/x^2.
\end{equation}
加上 $\omega(1)=0$ 唯一确定正文积分式；$A(1)=0$ 再保证 $\omega'(1)=0$。若只指定单纯形上的梯度，允许添加不改变散度的线性规范项；这里的规范化消除了该自由度。

更一般，若 $S(1)=0$、$S'(x)>0$，定义
\begin{equation}
\omega_S(x)=x\int_1^xS(u)u^{-2}\,du,\qquad
J_S(p;T)=D_{\omega_S}(T\Vert p).
\end{equation}
则 $\omega_S''=S'/x>0$、$\nabla_pJ_S=-S(T/p)$，熵镜像步为 $p^+\propto q e^{\tau S(T/q)}$。若 $L_S=\sup_{x>0}xS'(x)<\infty$，则
\begin{equation}
\nabla_p^2J_S=\diag[x_aS'(x_a)/p_a]\preceq L_S\nabla^2\Phi.
\end{equation}
这证明相对光滑性。它本身不保证统一方差、可估计中心或有限神经网络拟合。

\section{完整的下降与收敛推导}
固定 $\kappa,\eta,T$，记 $J=J_{\kappa,\eta}$，另取 $0<\tau\le\eta$。$p^+$ 最小化 $\langle\nabla J(q),p-q\rangle+\KL(p\Vert q)/\tau$。一阶条件与概率差作内积，得
\begin{equation}
\langle\nabla J(q),p^+-q\rangle
=-\tau^{-1}[\KL(p^+\Vert q)+\KL(q\Vert p^+)].
\end{equation}
将它代入相对光滑式：
\begin{equation}
J(q)-J(p^+)\ge\tau^{-1}\KL(q\Vert p^+)
+(\tau^{-1}-\eta^{-1})\KL(p^+\Vert q).
\label{eq:J-descent-app}
\end{equation}
凸性给出 $J(q)\le\langle\nabla J(q),q-T\rangle$，因为 $J(T)=0$。因此
\begin{align}
J(p^+)
&\le\langle\nabla J(q),p^+-T\rangle+\eta^{-1}\KL(p^+\Vert q)\\
&=\tau^{-1}[\KL(T\Vert q)-\KL(T\Vert p^+)]
-(\tau^{-1}-\eta^{-1})\KL(p^+\Vert q).
\label{eq:J-three-point}
\end{align}
第二行可直接核对：最优性条件使梯度与 $-\tau^{-1}\log(p^+/q)$ 相差常数，常数与 $p^+-T$ 的内积为零；展开三个 KL 即得到该式。

每步重新计算信号并刷新 $q_{k+1}=p^+(q_k)$，对式~\eqref{eq:J-three-point} 求和：
\begin{equation}
\sum_{k=0}^{K-1}\tau_kJ(q_{k+1})\le\KL(T\Vert q_0).
\end{equation}
由式~\eqref{eq:J-descent-app}，$J(q_K)\le J(q_{k+1})$，故
\begin{equation}
J(q_K)\le\frac{\KL(T\Vert q_0)}{\sum_{k<K}\tau_k}.
\end{equation}
固定正 $\tau$ 或 $\sum_k\tau_k=\infty$ 时目标趋于零。对 $\kappa=0$，KL 的连续延拓与唯一零点给出 $q_k\to T$。对 $\kappa>0$，还需排除边界极限。令
\begin{equation}
d_\infty=\int_1^\infty A(u)u^{-2}\,du.
\end{equation}
$0<A(u)\le(\log u)/\eta$ 在 $u>1$ 上成立，故 $0<d_\infty\le1/\eta$。若某坐标 $p_a\to0$，则 $p_a\omega(T_a/p_a)\to T_ad_\infty>0$。因此目标趋零的序列不可能在边界聚集；内点零点唯一为 $T$，结合有限单纯形紧性得到整个序列收敛。

上述求和要求同一个 $J$。逐轮改变 $\eta$ 或 $\kappa$ 时，不能直接把各轮不同函数的下降式望远镜相消。固定前缀测度的正加权和保留推导；策略变化引起的访问测度刷新另见附录~\ref{app:dynamics}。

\section{端点、JS 与曲率插值}
在固定正 $x$ 的紧积分区间上，
\begin{equation}
A_{0,\eta}(x)=\log x,\qquad
\lim_{\eta\downarrow0}A_{\kappa,\eta}(x)=\frac{x^\kappa-1}{\kappa},
\qquad
\lim_{\eta\uparrow1}A_{\kappa,\eta}(x)=\log x.
\end{equation}
分母在紧区间内有统一正下界，故可在 $\omega$ 的有限积分中取极限。对应二阶导数及在一处的两项规范化条件唯一确定生成函数，得到正文的三个 $J$ 端点。

对 $\kappa=1$，令 $b=1-\eta$、$B=b+\eta x$。分部积分给出
\begin{align}
\omega_{1,\eta}(x)
&=-\eta^{-1}\log B+x\int_1^x\frac{du}{u(b+\eta u)}\\
&=\frac{x\log x}{b}-\left(\eta^{-1}+\frac xb\right)\log B
=\frac{\eta x\log x-B\log B}{\eta b}.
\end{align}
代入 $x=T_a/p_a$ 并求和，记 $M_p=bp+\eta T$：
\begin{align}
J_{1,\eta}(p;T)
&=\frac{\eta\KL(T\Vert p)-\KL(M_p\Vert p)}{\eta b}\\
&=\frac{b\KL(p\Vert M_p)+\eta\KL(T\Vert M_p)}{\eta b}
=\frac{\JS_\eta(p,T)}{\eta(1-\eta)}.
\label{eq:JS}
\end{align}
第二个等号由 $M_p=bp+\eta T$ 展开对数得到。一般 $\kappa$ 的曲率满足
\begin{equation}
\frac1{\omega_{\kappa,\eta}''(x)}
=(1-\eta)x^{2-\kappa}+\eta x^2
=\frac{1-\eta}{g_\kappa''(x)}+\frac{\eta}{g_0''(x)}.
\end{equation}
因此这里插值的是生成函数曲率的倒数，通常不是两个散度的线性组合。各生成函数均有 $\omega''(1)=1$，故不同 $\eta$ 的目标也不能一般通过整体缩放互相替代。

\section{统一有界性与教师中心}
当 $\kappa>0,\eta<1$，令 $c=-\log(1-\eta)/(\kappa\eta)$。由 $A(x)\to-c$ 当 $x\downarrow0$，可得
\begin{equation}
\lim_{x\downarrow0}\omega(x)=c,\qquad
\lim_{x\to\infty}\omega(x)/x=d_\infty.
\end{equation}
第一式把积分拆为 $[x,b]$ 与 $[b,1]$：选小 $b$ 使 $|A(u)+c|\le\epsilon$，前段误差不超过 $\epsilon x(1/x-1/b)$，后段乘 $x$ 趋零。凸函数导数不超过无穷远渐近斜率，所以 $\omega(x)\le c+d_\infty x$。于是
\begin{equation}
0\le J(p;T)\le c+d_\infty.
\end{equation}
取 $p=(\epsilon,1-\epsilon)$、$T=(1-\epsilon,\epsilon)$，两坐标分别趋于 $d_\infty,c$，证明这是锐上确界。

教师中心有生成函数 $\phi(x)=x-1-A(x)$，因为 $\E_q(T/q-1)=0$：
\begin{equation}
K_T=D_\phi(T\Vert q).
\end{equation}
计算得
\begin{equation}
\frac{\phi''(x)}{\omega''(x)}
=1-\frac{\kappa(1-\eta)}{1-\eta+\eta x^\kappa}\in[1-\kappa,1].
\end{equation}
两生成函数在一处的值和一阶导数均为零。对二阶导数差积分两次，可知 $\phi-(1-\kappa)\omega$ 与 $\omega-\phi$ 均非负。因此
\begin{equation}
(1-\kappa)J(q;T)\le K_T\le J(q;T).
\label{eq:center-sandwich}
\end{equation}
中心作为当前分布函数时，梯度为
\begin{equation}
\partial_{p_a}K_T(p,T)
=-A(x_a)+\frac{x_a^\kappa}{1-\eta+\eta x_a^\kappa}-1.
\end{equation}
额外项一般随动作变化，所以中心不是教师目标 $J$ 的替代梯度。

\section{固定幂散度的镜像步与小步流}
若直接固定目标 $D_\kappa(T\Vert p)$，可取概率梯度代表 $(1-(T/p)^\kappa)/\kappa$。其步长 $\tau$ 的熵镜像解是
\begin{equation}
\widetilde m_a\propto q_a
\exp\left[\frac{\tau}{\kappa}\big((T_a/q_a)^\kappa-1\big)\right].
\end{equation}
它与幂均值目标仅一阶一致。对后者展开 $h_a=q_a[1+\eta(\rho_a^\kappa-1)/\kappa+O(\eta^2)]$，除以其总和得到式~\eqref{eq:velocity-main}。在 Fisher 内积 $\langle v,w\rangle_q=\sum_av_aw_a/q_a$ 下，任意切向量 $w$ 满足
\begin{equation}
\langle v_\kappa,w\rangle_q
=-\sum_aw_a\,\partial_{q_a}D_\kappa(T\Vert q),
\end{equation}
因为 $\sum_aw_a=0$ 消去常数。连续流 $\dot q=v_\kappa$ 因而有
\begin{equation}
\frac d{dt}D_\kappa(T\Vert q)
=-\|v_\kappa\|_q^2=-\frac{\Var_q(\rho^\kappa)}{\kappa^2}.
\end{equation}
这些展开固定正分布，余项不对任意极端密度比统一。

\chapter{内层投影与优化几何}
\label{app:inner}
\section{统一坐标中的四种目标}
\begin{proposition}[四种内层损失的统一分解]
固定 $m\propto qe^f$，令 $e=C_q\log(p/m)$、$\Lambda=\log\E_me^e$、$\Delta=\KL(q\Vert m)$。四种总体损失满足
\begin{align}
L_{\rm Part}&=\tfrac12\E_qe^2+\tfrac12\Lambda^2,&
\eta L_{\rm BPO}&=\tfrac12\E_qe^2,\\
L_{\rm Mean}&=\tfrac12\E_qe^2+\tfrac12(\Lambda+\Delta)^2,&
L_{\rm MDPO}&=\eta^{-1}[\KL(p\Vert m)-\log Z].
\end{align}
因而 $L_{\rm Part}-\eta L_{\rm BPO}=(\Delta-K_S)^2/2$，$L_{\rm Mean}-\eta L_{\rm BPO}=K_S^2/2$。
\end{proposition}
\begin{proof}
固定 $q,m,f$，其中 $m\propto qe^f$。令 $G=C_qf$、$\Delta=\log Z-\E_qf=\KL(q\Vert m)$，采用附录~\ref{app:measure} 的 $e=C_q\log(p/m)$ 与 $\Lambda(e)=\log\E_me^e$。则
\begin{equation}
U=G-\Delta+e-\Lambda(e),\qquad K_S=\Delta+\Lambda(e).
\end{equation}
代入四种损失：
\begin{align}
R_P&=e-\Lambda,& R_B&=e,&R_M&=e-\Lambda-\Delta,\\
L_{\rm Part}&=\tfrac12\E_qe^2+\tfrac12\Lambda^2,\\
\eta L_{\rm BPO}&=\tfrac12\E_qe^2,\\
L_{\rm Mean}&=\tfrac12\E_qe^2+\tfrac12(\Lambda+\Delta)^2,\\
L_{\rm MDPO}&=\eta^{-1}\KL(p(e)\Vert m)-\eta^{-1}\log Z.
\end{align}
三个平方分解使用 $\E_qe=0$，没有忽略任何均值项。它们也给出
\begin{equation}
L_{\rm Part}=\eta L_{\rm BPO}+\tfrac12(\Delta-K_S)^2,\quad
L_{\rm Mean}=\eta L_{\rm BPO}+\tfrac12K_S^2.
\end{equation}
求导后，额外项分别为 $(K_S-\Delta)\nabla K_S$ 与 $K_S\nabla K_S$。在锚点 $K_S=0$ 且 $\nabla K_S=0$，这些差异暂时消失。
\end{proof}

\section{Mean 的完整约束解}
\begin{lemma}[Mean 的归一化 KKT 候选]
设 $q>0$、$\E_qG=0$。Mean 的约束为 $\E_qe^U=1$。若归一化乘子 $\lambda\ge0$，其 KKT 方程唯一确定各坐标
\begin{equation}
U_a-G_a+\lambda e^{U_a}=0,\qquad U_a=G_a-W_0(\lambda e^{G_a}).
\end{equation}
\end{lemma}
\begin{proof}
Mean 优化 $\frac12\E_q(U-G)^2$，约束为 $\E_qe^U=1$。取拉格朗日函数
\begin{equation}
\mathscr L(U,\lambda)=\tfrac12\E_q(U-G)^2+\lambda(\E_qe^U-1).
\end{equation}
逐坐标一阶条件给出 $U_a-G_a+\lambda e^{U_a}=0$。当 $\lambda\ge0$ 时，解为
\begin{equation}
U_a=G_a-W_0(\lambda e^{G_a}),
\label{eq:mean-solution}
\end{equation}
其中 $W_0(x)e^{W_0(x)}=x$，取非负实主支。
\end{proof}

\begin{theorem}[Mean 全局解存在且唯一]
函数
\begin{equation}
\Psi(\lambda)=\sum_aq_a\exp[G_a-W_0(\lambda e^{G_a})]
\end{equation}
在 $\lambda\ge0$ 连续、严格递减，$\Psi(0)=\E_qe^G\ge1$ 且 $\Psi(\lambda)\to0$。故存在唯一 $\lambda\ge0$ 使 $\Psi(\lambda)=1$；由式~\eqref{eq:mean-solution} 构造的分布是唯一全局最优解 $p^{\rm M}$。此外
\begin{equation}
\lambda=\KL(q\Vert p^{\rm M}).
\end{equation}
\end{theorem}
\begin{proof}
$W_0$ 在非负轴连续严格递增，所以每个求和项严格递减。Jensen 给出起点界，有限项分别趋零给出无穷远极限。若 $G$ 非恒定，Jensen 严格，故根 $\lambda>0$；若 $G=0$，根为零。

固定这个乘子，$\nabla_U^2\mathscr L=\diag[q_a(1+\lambda e^{U_a})]\succ0$，所以构造点唯一最小化整个欧氏空间上的 $\mathscr L$。在所有可行 $U$ 上，乘子项为零；构造点又可行，所以它唯一最小化原目标。这一论证不需要指数等式约束集合凸。对 KKT 取 $q$ 均值，$\E_qG=0$、$\E_qe^U=1$ 给出 $\E_qU+\lambda=0$，即所述 KL 恒等式。
\end{proof}

\begin{proposition}[Mean 的混合正则变分解释]
固定由 Mean 全局解确定的 $\lambda=\KL(q\Vert p^{\rm M})$。同一 $p^{\rm M}$ 唯一最大化
\begin{equation}
\E_pf-\KL(p\Vert q)-\frac{\lambda}{2}\chi^2(p\Vert q).
\end{equation}
该命题断言最优分布相同，不断言两损失的梯度处处相同。
\end{proposition}
\begin{proof}
这个最优分布还唯一最大化
\begin{equation}
\E_pf-\KL(p\Vert q)-\frac{\lambda}{2}\chi^2(p\Vert q).
\label{eq:mean-mixed}
\end{equation}
证明是其概率梯度为 $f-U-1-\lambda(p/q-1)$，在 Mean 解处由 KKT 化为动作无关常数；Hessian 为 $-\diag(1/p+\lambda/q)\prec0$。这里 $\lambda$ 由 Mean 解决定，并非预先独立选择的正则系数。
\end{proof}

\begin{proposition}[Mean 超额损失的精确分解]
设 $p^{\rm M}$ 是精确 Mean 解，$\lambda$ 是其归一化乘子。对任意正分布 $p$，
\begin{equation}
L_{\rm Mean}(p)-L_{\rm Mean}(p^{\rm M})
=\tfrac12\E_q\log^2(p/p^{\rm M})+\lambda\KL(p^{\rm M}\Vert p).
\end{equation}
\end{proposition}
\begin{proof}
对任意正分布 $p$，还有精确超额损失分解
\begin{equation}
L_{\rm Mean}(p)-L_{\rm Mean}(p^{\rm M})
=\tfrac12\E_q\log^2(p/p^{\rm M})+\lambda\KL(p^{\rm M}\Vert p).
\label{eq:mean-excess}
\end{equation}
令 $v=U-U^\star$，平方展开的交叉项为 $\E_qv(U^\star-G)$。KKT 将它变成 $-\lambda\E_{p^{\rm M}}v=\lambda\KL(p^{\rm M}\Vert p)$，即得结论。
\end{proof}

\section{精确曲率与唯一驻点}
\begin{proposition}[平方回归曲率及锚点差异]
固定 $q,m,f$，以独立 logits 表示当前 $p$，$Q_p=\diag(p)-pp^\top$。Part 或 Mean 的残差记为 $R$，则
\begin{equation}
\nabla_z^2\tfrac12\E_qR^2=\E_q[(\mathbf e_a-p)(\mathbf e_a-p)^\top]-(\E_qR)Q_p.
\end{equation}
在 $p=q$ 处，Mean、Part、BPO 的 Hessian 分别为 $Q_q,(1-\Delta)Q_q,Q_q/\eta$，其中 $\Delta=\KL(q\Vert m)$。
\end{proposition}
\begin{proof}
以独立 logits $z$ 表示 $p$，令 $Q_p=\diag(p)-pp^\top$，以 $\mathbf e_a$ 表示第 $a$ 个坐标基向量，与中心化误差 $e$ 区分。于是
\begin{equation}
\nabla_z\log p_a=\mathbf e_a-p,\qquad \nabla_z^2\log p_a=-Q_p.
\end{equation}
对 Part 或 Mean 的残差 $R$，
\begin{equation}
\nabla_z\tfrac12\E_qR^2=q\odot R-(\E_qR)p,\qquad
\nabla_z^2\tfrac12\E_qR^2
=\E_q[(\mathbf e_a-p)(\mathbf e_a-p)^\top]-(\E_qR)Q_p.
\end{equation}
在锚点，Mean 残差均值为零、Part 为 $\Delta$，故
\begin{equation}
\nabla_z^2L_{\rm Mean}|_q=Q_q,\quad
\nabla_z^2L_{\rm Part}|_q=(1-\Delta)Q_q,\quad
\nabla_z^2L_{\rm BPO}=Q_q/\eta.
\end{equation}
例如 $q=(1/2,1/2)$、$f=(2,-2)$，则 $\Delta=\log\cosh2>1$，Part 在全部非常数方向有负曲率。
\end{proof}

\begin{theorem}[Part 的唯一内部驻点]
若 $q,m>0$，则 Part 在独立有限 logits 参数化中的唯一概率内部驻点是 $p=m$。该结论允许 Part 在其他位置具有负曲率。
\end{theorem}
\begin{proof}
Part 仍只有一个有限内部驻点。记 $R=\log(p/m)$、$c=\E_qR$，驻点条件为 $q_aR_a=cp_a$。若 $c>0$，每个 $R_a>0$，于是每个 $p_a>m_a$，与归一化矛盾；若 $c<0$，同理每个 $p_a<m_a$ 也矛盾。故 $c=0$，正 $q$ 强制 $R=0$。内部驻点唯一不等于任意优化器全局收敛。
\end{proof}

\begin{proposition}[MDPO 的精确方向 Hessian]
固定正目标 $m$，沿任意 logit 方向 $v$，令 $p_t\propto pe^{tv}$、$h=\log(p/m)$。则
\begin{equation}
\frac{d^2}{dt^2}\KL(p_t\Vert m)\big|_0
=\Var_pv+\Cov_p\!\left(h,(v-\E_pv)^2\right).
\end{equation}
MDPO 损失的对应导数为右侧除以 $\eta$；在 $p=m$ 时协方差项为零。
\end{proposition}
\begin{proof}
MDPO 的完整二阶方向导数也可以写出，无需以未定义的残差项代替。固定 logit 方向 $v$，沿 $p_t\propto p e^{tv}$，令 $h=\log(p/m)$、$\bar v=\E_pv$。指数倾斜的微分恒等式为 $\frac d{dt}\E_{p_t}u=\Cov_{p_t}(u,v)$，若 $u$ 依赖 $t$ 再加 $\E_{p_t}\dot u$。因此
\begin{align}
\frac d{dt}\KL(p_t\Vert m)\big|_0&=\Cov_p(h,v),\\
\frac{d^2}{dt^2}\KL(p_t\Vert m)\big|_0
&=\Var_p(v)+\Cov_p\bigl(h,(v-\bar v)^2\bigr).
\label{eq:mdpo-hessian}
\end{align}
第二式中 $\dot h=v-\bar v$ 产生方差项；对协方差中的分布求导产生中心三阶混合矩。MDPO 的 Hessian 再除以 $\eta$。当 $p=m$ 时第二项为零；一般位置不能删除它。
\end{proof}

\begin{corollary}[中心化坐标中的回归曲率]
设 $e=C_q\log(p/m)$、$\Lambda=\log\E_me^e$，对 $\E_qv=0$ 的方向有
\begin{align}
\nabla^2L_{\rm Part}[v,v]&=\E_qv^2+\Lambda\Var_pv+(\E_pv)^2,\\
\nabla^2L_{\rm Mean}[v,v]&=\E_qv^2+(\Lambda+\Delta)\Var_pv+(\E_pv)^2.
\end{align}
\end{corollary}
\begin{proof}
在 $\E_qe=0$ 的坐标中，对 $\E_qv=0$ 的方向，
\begin{align}
\nabla^2 L_{\rm Part}[v,v]&=\E_qv^2+\Lambda\Var_pv+(\E_pv)^2,\\
\nabla^2 L_{\rm Mean}[v,v]&=\E_qv^2+(\Lambda+\Delta)\Var_pv+(\E_pv)^2.
\end{align}
这是对 $\Lambda'=\E_pv$、$\Lambda''=\Var_pv$ 使用乘积法则的结果。
\end{proof}

\section{受限表示改变投影}
\begin{proposition}[受限表示下 Part 与 BPO 可分离]
存在完整正支持的三动作锚点、目标及一维共享 logit 学生族，使 BPO 的唯一最优点不是 Part 的驻点。
\end{proposition}
\begin{proof}
取 $q=(1/3,1/3,1/3)$、$m=(0.1,0.7,0.2)$、$\eta=1$，学生族为 $p(z)=\operatorname{softmax}(z,0,0)$。令 $a=\log0.1,b=\log0.7,c=\log0.2$。BPO 为 $\frac12\Var_q(z-a,-b,-c)$，导数为 $\frac29[z-a+(b+c)/2]$，故最优 $z_B=a-(b+c)/2$。

令 $u(z)=\E_q\log[p(z)/m]=z/3-\log(e^z+2)-(a+b+c)/3$，则 $L_{\rm Part}=L_{\rm BPO}+u^2/2$，所以
\begin{equation}
L_{\rm Part}'(z_B)=u(z_B)[1/3-p_1(z_B)]\simeq0.03543884\ne0.
\end{equation}
两个目标在同一受限学生族中具有不同最优点。
\end{proof}

\begin{proposition}[跨前缀共享参数下 MDPO 与 BPO 可分离]
存在两个等权前缀及一个共享相对 logit 参数，使 BPO 的最优参数不满足 MDPO 的驻点方程。
\end{proposition}
\begin{proof}
BPO 与 MDPO 也可分离。取两个等权前缀，锚点第一动作概率分别为 $0.8,0.2$，共享相对 logit 位移 $\theta$；目标相对位移取 $g_1=1,g_2=3$。BPO 的权重均为 $0.16$，故 $\theta_B=2$。MDPO 的驻点方程为
\begin{equation}
\sum_{i=1}^2p_{\theta,i}(1-p_{\theta,i})(\theta-g_i)=0.
\end{equation}
在 $\theta=2$，第一项权重约为 $0.0316$，第二项约为 $0.2278$，和非零。该方程是必要驻点条件，不把任意根直接称为唯一全局解。
\end{proof}

\section{自由 logits 的收敛与理想路径}
\begin{theorem}[BPO 二次结构与梯度下降的稳定区间]
固定 $q,m>0$，独立 logits 中
\begin{equation}
L_{\rm BPO}=\frac{(z-\log m)^\top Q_q(z-\log m)}{2\eta},
\quad\nabla_zL_{\rm BPO}=\eta^{-1}Q_q(z-\log m).
\end{equation}
在 $\mathbf1^\perp$ 上，$q_{\min}I\preceq Q_q\preceq I/2$。常步长梯度下降对所有有限初值收敛，当且仅当 $0<\alpha_{\rm opt}<2\eta/\lambda_{\max}(Q_q)$。
\end{theorem}
\begin{proof}
令 $h=z-\log m$。中心化消去 softmax 配分函数，得到
\begin{equation}
L_{\rm BPO}=\frac{h^\top Q_qh}{2\eta},\quad
\nabla_zL_{\rm BPO}=\frac{Q_qh}{\eta}.
\end{equation}
$Q_q$ 的零空间为常数方向。在欧氏子空间 $\mathbf1^\perp$ 上，
\begin{equation}
q_{\min}\|v\|^2\le v^\top Q_qv\le\tfrac12\|v\|^2.
\end{equation}
下界由 $\Var_qv=\min_c\sum_aq_a(v_a-c)^2\ge q_{\min}\min_c\sum_a(v_a-c)^2$ 得到。上界由方差不超过分量范围平方的 $1/4$，以及 $(v_i-v_j)^2\le2\|v\|^2$ 得到。

普通梯度下降的每个非零谱方向乘以 $1-\alpha_{\rm opt}\lambda_i/\eta$。对全部初值收敛的充要条件是
\begin{equation}
0<\alpha_{\rm opt}<2\eta/\lambda_{\max}(Q_q).
\end{equation}
其中 $\alpha_{\rm opt}<4\eta$ 是统一充分条件。
\end{proof}

\begin{proposition}[两种 Fisher 预条件步的相同概率轨迹]
完整词表、独立 logits 下，MDPO 使用当前 $Q_p^\dagger$，BPO 使用冻结 $Q_q^\dagger$，均产生 $p^+\propto p^{1-\gamma}m^\gamma$，其中 $\gamma=\alpha_{\rm opt}/\eta$。从 $p_0=q$ 出发，若 $0<\gamma_j\le1$，则
\begin{equation}
p_n\propto qe^{\xi_nf},\quad \xi_n=1-\prod_{j<n}(1-\gamma_j),\quad
L_{\rm BPO}(p_n)=(1-\xi_n)^2L_{\rm BPO}(q).
\end{equation}
\end{proposition}
\begin{proof}
MDPO 的 logit 梯度为 $Q_p(z-\log m)/\eta$。正支持下 $Q_p^\dagger Q_p=P_\perp$；BPO 使用 $Q_q^\dagger$ 也得到同一投影。因此两种对应的精确自然梯度步都满足
\begin{equation}
P_\perp h^+=(1-\gamma)P_\perp h,\quad
p^+\propto p^{1-\gamma}m^\gamma,\qquad
\gamma=\alpha_{\rm opt}/\eta.
\end{equation}
BPO 的 Fisher 冻结在 $q$，MDPO 的 Fisher 来自当前 $p$。从 $p_0=q$ 出发，若 $0<\gamma_j\le1$，
\begin{equation}
p_n\propto qe^{\xi_nf},\qquad
\xi_n=1-\prod_{j<n}(1-\gamma_j),\qquad
L_{\rm BPO}(p_n)=(1-\xi_n)^2L_{\rm BPO}(q).
\end{equation}
\end{proof}

\begin{theorem}[理想内层路径的教师散度进展]
对所有凸 $\varphi$、$\varphi(1)=0$，沿 $p_\xi\propto qe^{\xi f}$、$0\le\xi\le1$，$D_\varphi(p_\xi\Vert T)$ 非增。
\end{theorem}
\begin{proof}
令 $r=\log(T/q)$。$f(r)$ 递增且
\begin{equation}
f'(r)=\frac{\eta e^{\kappa r}}{1-\eta+\eta e^{\kappa r}}\le1
\end{equation}
（几何端点为 $\eta$）。因此 $u_\xi(r)=p_\xi/T=e^{-r+\xi f(r)}/Z_\xi$ 非增。可微凸 $\varphi$ 的导数递增，故 $\varphi'(u_\xi(r))$ 非增。微分归一化密度：
\begin{align}
\frac d{d\xi}D_\varphi(p_\xi\Vert T)
&=\Cov_{p_\xi}(\varphi'(u_\xi),f)\\
&=\frac12\sum_{a,b}p_{\xi,a}p_{\xi,b}
[\varphi'(u_{\xi,a})-\varphi'(u_{\xi,b})](f_a-f_b)\le0.
\end{align}
每个乘积非正。有限正支持使所有 $u_{\xi,a}$ 位于紧正区间；对任意有限凸函数用光滑凸函数一致逼近，取极限保留非增结论。
\end{proof}

\section{共享参数、拟合形状与条件 SGD}
\begin{proposition}[共享参数的复合 Hessian]
若 $z_\theta$ 二次可微，$B_\theta=\partial z/\partial\theta$、$h=z-\log m$，则
\begin{equation}
\nabla_\theta^2L_{\rm BPO}
=\eta^{-1}\left[B_\theta^\top Q_qB_\theta+\sum_a(Q_qh)_a\nabla_\theta^2z_a\right].
\end{equation}
所以固定的 logit Hessian 不保证参数空间目标凸。
\end{proposition}
\begin{proof}
链式法则先给 $\nabla_\theta L=B_\theta^\top\nabla_zL$。对第 $i,j$ 个参数再求导，
\begin{equation}
\partial_j\partial_iL
=\sum_{a,b}(\partial_i z_a)(\partial_{z_b}\partial_{z_a}L)(\partial_jz_b)
+\sum_a(\partial_{z_a}L)\partial_j\partial_i z_a.
\end{equation}
代入 $\nabla_zL=Q_qh/\eta$ 与 $\nabla_z^2L=Q_q/\eta$，得到
\begin{equation}
\nabla_\theta^2L_{\rm BPO}
=\eta^{-1}\left[B_\theta^\top Q_qB_\theta
+\sum_a(Q_qh)_a\nabla_\theta^2z_a\right].
\end{equation}
这里 $B_\theta=\partial z/\partial\theta$，即使冻结 $Q_q$，两项仍可变化。自由 logit 的凸性不自动传到参数空间。
\end{proof}

\begin{proposition}[拟合形状与残差二阶矩]
固定目标 $m$。若从当前 $p$ 采样，$x=m/p$、$a_\beta=\E_px^\beta$，则 $g_\beta=(a_\beta-x^\beta)\psi/\beta$ 对 $\nabla_\theta D_\beta(m\Vert p)$ 无偏。若 $0<\beta\le1/2$ 且 $\|\psi\|\le S$，有 $\E_p\|g_\beta\|^2\le S^2D_\beta/\beta$。精确 BPO 的锚点采样梯度则满足 $\E_q\|R_B\psi/\eta\|^2\le(2S^2/\eta)L_{\rm BPO}$。
\end{proposition}
\begin{proof}
也可固定目标 $m_{\kappa,\eta}$ 后另选拟合形状 $\beta$，优化 $D_\beta(m\Vert p)$。从当前学生采样，令 $x=m/p$、$a_\beta=\E_px^\beta$，则
\begin{equation}
g_\beta(a)=\frac{a_\beta-x_a^\beta}{\beta}\psi_a,\qquad
\E_pg_\beta=\nabla_\theta D_\beta(m\Vert p).
\end{equation}
对 $0<\beta\le1/2$、$\|\psi_a\|\le S$，附录~\ref{app:shape} 的矩不等式给出 $\E_p\|g_\beta\|^2\le S^2D_\beta(m\Vert p)/\beta$。$\beta$ 控制拟合，$\kappa$ 控制目标，不能混为一个最优形状判断。BPO 本身也满足
\begin{equation}
\E_q\|R_B\psi/\eta\|^2\le(2S^2/\eta)L_{\rm BPO}.
\end{equation}
\end{proof}

\begin{theorem}[误差相关噪声下的条件 SGD 收敛]
假设所有更新线段上 $L$ 为 $L_{\rm sm}$ 光滑，满足 $\|\nabla L\|^2\ge2\mu(L-L_\star)$，且随机梯度条件无偏、$\E[\|g\|^2\mid\theta]\le CL$。若 $0<\alpha\le\min\{1/\mu,2\mu/(L_{\rm sm}C)\}$，则
\begin{equation}
\E(L_N-L_\star)\le(1-\mu\alpha)^N(L_0-L_\star)
+\frac{L_{\rm sm}C\alpha L_\star}{2\mu}[1-(1-\mu\alpha)^N].
\end{equation}
\end{theorem}
\begin{proof}
设 $L$ 在所有更新线段上为 $L_{\rm sm}$ 光滑，满足 PL 条件 $\|\nabla L\|^2\ge2\mu(L-L_\star)$，随机梯度条件无偏且 $\E[\|g\|^2\mid\theta]\le CL$。对 $0<\alpha\le\min\{1/\mu,2\mu/(L_{\rm sm}C)\}$，光滑性给出
\begin{align}
\E[L(\theta-\alpha g)\mid\theta]
&\le L-\alpha\|\nabla L\|^2+\tfrac12L_{\rm sm}\alpha^2CL\\
&\le L_\star+(1-\mu\alpha)(L-L_\star)
+\tfrac12L_{\rm sm}C\alpha^2L_\star.
\end{align}
取全期望并展开几何递推，得到
\begin{equation}
\E(L_N-L_\star)\le(1-\mu\alpha)^N(L_0-L_\star)
+\frac{L_{\rm sm}C\alpha L_\star}{2\mu}[1-(1-\mu\alpha)^N].
\end{equation}
这是明确假设下的 SGD 结论。模型可实现时 $L_\star=0$，该上界没有固定采样噪声底；数据复用、中心近似或其他优化器需要重新核对假设。
\end{proof}

\section{与条件 OPD 基线的关系}
\begin{proposition}[几何镜像目标与条件 OPD 的关系]
取 $\kappa=0$，$f=\eta\log(T/q)$。则
\begin{equation}
L_{\rm MDPO}=\KL(p\Vert T)+(\eta^{-1}-1)\KL(p\Vert q).
\end{equation}
因此 $\eta=1$ 时 MDPO 等于条件 OPD 目标。固定 $q,T$ 并令 $\eta\to\infty$ 时，其极限为冻结信号目标 $-\E_p\log(T/q)$，最优分布支持在最高信号动作上。
\end{proposition}
\begin{proof}
对当前学生--教师目标 $\KL(p\Vert T)$，令 $r^{\rm cur}=\log(T/p)$。其精确梯度及冻结锚点下的表示为
\begin{equation}
\nabla_\theta\KL(p\Vert T)
=-\E_p[r^{\rm cur}\psi_p]
=-\E_q[(p/q)r^{\rm cur}\psi_p].
\end{equation}
这个重要性比只改变当前前缀下的动作积分，不改变前缀访问测度。

几何端点取 $f=\eta\log(T/q)$，MDPO 恒等式为
\begin{equation}
L_{\rm MDPO}=\KL(p\Vert T)+(\eta^{-1}-1)\KL(p\Vert q).
\end{equation}
因此 $\eta=1$ 时 MDPO 的总体目标处处等于条件 OPD。BPO 在该点仍使用 $\Var_q\log(p/T)/2$，Mean 还包含学生中心平方，二者只在特定位置与 OPD 有相同梯度或共同目标分布。

对固定几何 reward $r=\log(T/q)$，令 $\eta\to\infty$ 删除镜像正则，MDPO 目标趋于 $-\E_pr$，其最优分布支持在 $r$ 的最大值集合，通常不等于 $T$。如果保持当前 OPD 目标，则 reward 必须随学生改为 $r^{\rm cur}=r-U$。两者只在初始化 $p=q$ 处相同。这里的无穷步长讨论仅适用于几何端点，不将正形状的混合权重扩展到概率区间之外。
\end{proof}

\section{标量步长能改变的部分}
\begin{proposition}[固定方向的最优步长与下降能力]
设 $E(\theta)=e^\top He/2$、$H\succeq0$、$E>0$，方向 $d$ 固定。记 $A=\nabla E^\top d$、$B=d^\top Hd$。则 $E(\theta-\alpha d)=E-\alpha A+\alpha^2B/2$。若 $A,B>0$，最优步长为 $A/B$，最大相对下降为 $A^2/(2BE)\le1$；若 $A\le0$，正步长不能下降。
\end{proposition}
\begin{proof}
固定二次超额损失 $E=\frac12e^\top He$，沿方向 $d$ 更新时，
\begin{equation}
E(\theta-\alpha d)=E-\alpha A+\tfrac12\alpha^2B,\quad
A=\nabla E^\top d,\quad B=d^\top Hd.
\end{equation}
若 $A>0,B>0$，最优步长 $\alpha_*=A/B$，最大相对下降为 $\rho_*=A^2/(2BE)\le1$。最后一个界来自 $H$ 半内积的 Cauchy--Schwarz。达到 $0<\rho\le\rho_*$ 的较小正根为
\begin{equation}
\alpha_\rho=\frac{2\rho E}{A+\sqrt{A^2-2B\rho E}}.
\end{equation}
若 $A\le0$，任何正步长都不能降低这个二次目标；标量调度无法修复方向错误。
\end{proof}

\begin{proposition}[谱耗散与精确线搜索的速率]
对固定正谱二次目标的梯度流，以损失谱分量占比为 $w_i$，有 $-d\log E/dt=2\E_w\lambda$，其导数为 $-4\Var_w\lambda$。若谱位于 $[l,u]$、$l>0$，精确最速下降线搜索满足 $E_{n+1}\le[(u-l)/(u+l)]^2E_n$。
\end{proposition}
\begin{proof}
梯度流的谱表达为 $E(t)=\frac12\sum_i\lambda_ie_i(0)^2e^{-2\lambda_it}$。以各项在总损失中的占比为 $w_i$，求导得
\begin{equation}
-\frac d{dt}\log E=2\E_w\lambda,\qquad
\frac d{dt}\left[-\frac d{dt}\log E\right]=-4\Var_w\lambda\le0.
\end{equation}
快方向先消失，使相对下降率变慢。精确线搜索在谱区间 $[l,u]$ 的最坏收缩因子仍为 $[(u-l)/(u+l)]^2$：由 $\lambda^2\le(l+u)\lambda-lu$，最小化 $\E_w\lambda^2$ 上界所对应的比例可得 $\rho_*\ge4lu/(l+u)^2$。
\end{proof}

\begin{proposition}[固定谱下的 Chebyshev 步序]
对固定二次目标，非零谱位于 $0<l<u$。令 $c=(u+l)/2,d=(u-l)/2$，在 $N$ 步使用 $\alpha_j=[c-d\cos((2j-1)\pi/(2N))]^{-1}$，则
\begin{equation}
E_N\le4\left(\frac{\sqrt{u/l}-1}{\sqrt{u/l}+1}\right)^{2N}E_0.
\end{equation}
该结论允许中间误差放大，要求谱和线性更新固定。
\end{proposition}
\begin{proof}
若固定谱端点已知，Chebyshev 步序 $\alpha_j=[c-d\cos((2j-1)\pi/(2N))]^{-1}$，其中 $c=(u+l)/2,d=(u-l)/2$，产生误差多项式 $T_N((c-\lambda)/d)/T_N(c/d)$，故
\begin{equation}
E_N\le4\left(\frac{\sqrt{u/l}-1}{\sqrt{u/l}+1}\right)^{2N}E_0.
\end{equation}
该恒等式由两侧多项式具有相同根和零点处归一化得到；$|T_N|\le1$ 在 $[-1,1]$ 上，再用 $T_N(x)=(v^N+v^{-N})/2$ 得到界。$u=l$ 时单步 $\alpha=1/l$ 即可。这个结论需要固定线性几何，且允许中间误差放大。
\end{proof}

\begin{proposition}[Adam 的常数损失缩放]
在相同原始梯度历史、零初始矩状态下，将全部历史梯度乘以常数 $s>0$，使 Adam 一阶矩乘以 $s$、二阶矩乘以 $s^2$。带偏差修正的方向满足
\begin{equation}
\frac{\widehat m'}{\sqrt{\widehat v'}+\epsilon}
=\frac{\widehat m}{\sqrt{\widehat v}+\epsilon/s}.
\end{equation}
因此常数 loss 缩放通常不同于将学习率乘以 $s$；时变缩放还会改变矩历史。
\end{proposition}
\begin{proof}
记原始梯度历史为 $g_t$，矩初值为零。缩放后的递推为
\begin{align}
m_t'&=\beta_1m_{t-1}'+(1-\beta_1)sg_t,\\
v_t'&=\beta_2v_{t-1}'+(1-\beta_2)s^2g_t^2.
\end{align}
从零初值归纳得到 $m_t'=sm_t$、$v_t'=s^2v_t$。偏差修正分母只依赖时间与 $\beta_1,\beta_2$，故 $\widehat m_t'=s\widehat m_t$、$\widehat v_t'=s^2\widehat v_t$。对更新方向的分子分母同除以正数 $s$，得到命题中的等式。若 $s_t$ 随时间改变，两个矩分别包含历史 $s_j$ 和 $s_j^2$ 的加权和，不能提出同一个常数。
\end{proof}

沿固定优化器方向作三点测量可由 $Q(0),Q(h),Q(2h)$ 恢复二次式的 $A,B$，但非线性网络中必须用实际目标值验收；一次试探不能重复更新优化器矩状态。

\chapter{形状、概率路径与监督矩}
\label{app:shape}
\section{覆盖界与目标半径}
\begin{theorem}[正形状的覆盖与刷新后信号范围]
固定正分布 $q,T$，$0<\kappa\le1,0<\eta<1$。幂均值的归一化常数和目标满足
\begin{equation}
(1-\eta)^{1/\kappa}+\eta^{1/\kappa}\le Z\le1,\qquad
m_a\ge(1-\eta)^{1/\kappa}q_a+\eta^{1/\kappa}T_a.
\end{equation}
因此 $\KL(q\Vert m)\le-\log(1-\eta)/\kappa$、$\KL(T\Vert m)\le-\log\eta/\kappa$。精确刷新 $q^+=m$ 后，
$-c_{\kappa,\eta}\le A(T_a/q_a^+)\le\log(2-\eta)/(\kappa\eta)$。
\end{theorem}
\begin{proof}
对 $0<\kappa\le1$，幂均值不等式给 $h_a\le(1-\eta)q_a+\eta T_a$，所以 $Z\le1$。令 $r=1/\kappa\ge1$，非负 $x,y$ 满足 $(x+y)^r\ge x^r+y^r$；例如对固定 $x$，两侧之差对 $y$ 的导数非负，且在 $y=0$ 为零。因此
\begin{equation}
(1-\eta)^{1/\kappa}+\eta^{1/\kappa}\le Z\le1,\qquad
m_a\ge(1-\eta)^{1/\kappa}q_a+\eta^{1/\kappa}T_a.
\end{equation}
逐坐标取对数并按相应分布平均，得到
\begin{equation}
\KL(q\Vert m)\le-\kappa^{-1}\log(1-\eta),\quad
\KL(T\Vert m)\le-\kappa^{-1}\log\eta.
\end{equation}
固定教师的精确刷新 $q^+=m$ 还给后续概率比 $(T_a/q_a^+)^\kappa\le1/\eta$，所以下一轮
\begin{equation}
-c_{\kappa,\eta}\le A_{\kappa,\eta}(T_a/q_a^+)
\le\frac{\log(2-\eta)}{\kappa\eta}.
\end{equation}
这个双侧界依赖精确刷新产生的覆盖；任意初始锚点只有单侧下界。
\end{proof}

\begin{proposition}[目标到锚点的半径]
固定 $q,T,\kappa$。对任意凸 $\varphi$、$\varphi(1)=0$，$D_\varphi(m_{\kappa,\eta}\Vert q)$ 随 $\eta\in[0,1]$ 非减。若 $q\ne T$ 且 $\varphi$ 严格凸，内部严格递增。
\end{proposition}
\begin{proof}
取 $0<\eta_1<\eta_2<1$，令 $\rho=T/q$、$u_i=m_{\eta_i}/q$。$u_1$ 随 $\rho$ 递增，而
\begin{equation}
\frac{m_{\eta_2}}{m_{\eta_1}}
=\frac{Z_{\eta_1}}{Z_{\eta_2}}
\left(\frac{1-\eta_2+\eta_2\rho^\kappa}
{1-\eta_1+\eta_1\rho^\kappa}\right)^{1/\kappa}
\end{equation}
也严格递增：括号对 $\rho^\kappa$ 的导数为 $(\eta_2-\eta_1)/(1-\eta_1+\eta_1\rho^\kappa)^2>0$。几何端点的比值为常数乘 $\rho^{\eta_2-\eta_1}$。该比值在 $m_{\eta_1}$ 下均值一，故存在交点 $\rho_\star$ 使比值等于一。

记 $u_\star=u_1(\rho_\star)$。低于交点时 $u_2\le u_1\le u_\star$，高于交点时 $u_2\ge u_1\ge u_\star$。取 $b\in\partial\varphi(u_\star)$，则 $\varphi(u)-bu$ 在 $u_\star$ 两侧向外非减，逐点得到 $\varphi(u_2)-bu_2\ge\varphi(u_1)-bu_1$。按 $q$ 求和，因两组 $u$ 的均值均为一，线性项抵消。严格凸与不同目标保证至少一项严格；端点由连续性得到。
\end{proof}

\section{精确目标同时改善凸生成散度}
\begin{theorem}[共同交点的散度进展]
对任意凸 $\varphi$、$\varphi(1)=0$，精确幂均值目标满足
\begin{equation}
D_\varphi(m\Vert T)\le D_\varphi(q\Vert T).
\end{equation}
\end{theorem}
\begin{proof}
写 $s_a=q_a/T_a$、$m_a=T_aG(s_a)/Z$，其中 $G(s)=[(1-\eta)s^\kappa+\eta]^{1/\kappa}$，几何端点为 $G(s)=s^{1-\eta}$。$G$ 递增，$G(s)/s$ 递减，因为
\begin{equation}
\frac d{d\log s}\log[G(s)/s]
=-\frac{\eta s^{-\kappa}}{1-\eta+\eta s^{-\kappa}}<0.
\end{equation}
又有 $Z=\sum_aq_aG(s_a)/s_a$，所以存在 $s_\star$ 使 $G(s_\star)/s_\star=Z$。若 $s_a<s_\star$，则 $s_a\le G(s_a)/Z\le s_\star$；若 $s_a>s_\star$，不等号反向排列。取 $\varphi$ 在 $s_\star$ 的次梯度 $b$，函数 $\varphi(s)-bs$ 向 $s_\star$ 移动时非增。逐项比较后按 $T$ 求和，归一化消去线性项，得到结论。
\end{proof}

\begin{proposition}[精确更新的可逆 Markov 表示]
对固定正分布 $q,T$ 和精确幂均值目标 $m$，存在行随机矩阵 $P$，满足 $TP=T$、$qP=m$，且关于 $T$ 详细平衡：$T_aP_{ab}=T_bP_{ba}$。
\end{proposition}
\begin{proof}
这个更新还可由保持 $T$ 的可逆 Markov 核表示。令 $\delta_a=m_a-q_a$，正质量接收集合为 $L$，输出集合为 $H$，总转移量 $D=\sum_L\delta_a$。若 $D=0$ 取恒等核；否则对 $a\in L,b\in H$，定义
\begin{equation}
w_{ab}=w_{ba}=\frac{\delta_a(-\delta_b)}{D(s_b-s_a)},\qquad
P_{ab}=w_{ab}/T_a,
\end{equation}
其余非对角项为零，对角项补足行和。由共同交点，$\delta_a/T_a\le s_\star-s_a$、$(-\delta_b)/T_b\le s_b-s_\star$，故 $\sum_bw_{ab}\le T_a$、$\sum_aw_{ab}\le T_b$，核非负。详细平衡 $T_aP_{ab}=T_bP_{ba}$ 保证 $TP=T$；而接收坐标有 $(qP)_a-q_a=\sum_bw_{ab}(s_b-s_a)=\delta_a$，输出坐标同理，故 $qP=m$。这是数学表示，无需实际构造词表平方大小的矩阵。
\end{proof}

\section{教师中心的性质}
\begin{proposition}[教师中心的锐统一界]
对 $0<\kappa\le1,0<\eta<1$，任意正分布满足
\begin{equation}
0\le K_T(q,T)\le c_{\kappa,\eta}=-\frac{\log(1-\eta)}{\kappa\eta}.
\end{equation}
至少二动作时，该上界可由严格正分布逼近。
\end{proposition}
\begin{proof}
写 $\rho=T/q$，则 $A(1)=0$、$A'>0$ 且 $A''\le0$。Jensen 给出 $\E_qA(\rho)\le A(\E_q\rho)=0$；结合 $A\ge-c$，
\begin{equation}
0\le K_T=-\E_qA\le c_{\kappa,\eta}.
\end{equation}
上界为锐上确界：取 $q=(\epsilon,1-\epsilon)$、$T=(1-\epsilon^2,\epsilon^2)$。常见动作的 $-A\to c$，稀有动作的贡献为 $O(\epsilon|\log\epsilon|)\to0$。
\end{proof}

\begin{proposition}[中心对两个参数的单调性]
固定 $q,T>0$，教师中心 $K_T$ 随 $\kappa$ 非增、随 $\eta$ 非减，参数范围为 $0<\kappa\le1,0<\eta<1$。该结论不对监督方差给出同样排序。
\end{proposition}
\begin{proof}
固定概率比时，$A$ 随 $\kappa$ 非减：函数 $L(u)=\log(1-\eta+\eta e^{u\log\rho})$ 凸且 $L(0)=0$，故 $L(u)/u$ 在 $u>0$ 非减。另一方面，令 $y=1+\eta(\rho^\kappa-1)$，有
\begin{equation}
\partial_\eta A=\frac{1-y^{-1}-\log y}{\kappa\eta^2}\le0.
\end{equation}
不等式来自 $\log y\ge1-y^{-1}$。因此 $K_T$ 随 $\kappa$ 非增，随 $\eta$ 非减。标量方差没有由此得到同一排序。
\end{proof}

\begin{proposition}[中心的端点与局部尺度]
固定正分布时，$K_T\to(1-\kappa)D_\kappa(T\Vert q)$ 当 $\eta\downarrow0$，而 $\eta\uparrow1$ 时 $K_T\to\KL(q\Vert T)$。若全部 $\rho=T/q$ 一致趋近一，则
\begin{equation}
K_T=\tfrac12[1-\kappa(1-\eta)]\E_q(\rho-1)^2+O(\E_q|\rho-1|^3).
\end{equation}
\end{proposition}
\begin{proof}
固定概率比时，$\log[1+\eta(\rho^\kappa-1)]/\eta\to\rho^\kappa-1$。有限求和可逐项取极限，所以
\begin{align}
\lim_{\eta\downarrow0}K_T
&=\frac{1-\E_q\rho^\kappa}{\kappa}=(1-\kappa)D_\kappa(T\Vert q),\\
K_T&=\frac{1-\kappa(1-\eta)}2\E_q(\rho-1)^2
+O(\E_q|\rho-1|^3).
\end{align}
第一式对 $\kappa=1$ 取零。为证明第二式，直接计算
\begin{equation}
A(1)=0,\qquad A'(1)=1,\qquad A''(1)=\kappa(1-\eta)-1=-c_0.
\end{equation}
在一附近 Taylor 展开为 $A(1+\delta)=\delta-c_0\delta^2/2+O(|\delta|^3)$。按 $q$ 取负期望，一次项由 $\E_q(\rho-1)=0$ 消失，得到局部尺度。$\eta\uparrow1$ 时逐项 $A(\rho)\to\log\rho$，所以 $K_T\to\KL(q\Vert T)$；统一有界性不延续到该端点。
\end{proof}

\section{固定强度的矩界与小步阈值}
\begin{proposition}[固定强度的维数无关二阶矩界]
令 $C_\star=\sup_{u>0}\log^2(1+u)/u<\infty$。对 $0<\kappa\le1,0<\eta<1$，未缩放信号 $f=\eta A$ 满足
\begin{equation}
\E_qf^2\le[\log^2(1-\eta)+C_\star\eta]/\kappa^2.
\end{equation}
标准化信号 $A$ 的相应上界为右侧除以 $\eta^2$。
\end{proposition}
\begin{proof}
定义 $C_\star=\sup_{u>0}\log^2(1+u)/u<\infty$。函数在零与无穷远都趋零，故上确界有限。对 $f=\eta A$，在 $\rho<1$ 上 $|f|\le|\log(1-\eta)|/\kappa$；在 $\rho\ge1$ 上
\begin{equation}
0\le f\le\kappa^{-1}\log(1+\eta\rho^\kappa),\qquad
f^2\le C_\star\eta\rho^\kappa/\kappa^2.
\end{equation}
凹性给出 $\E_q\rho^\kappa\le1$，所以
\begin{equation}
\E_qf^2\le\frac{\log^2(1-\eta)+C_\star\eta}{\kappa^2}.
\label{eq:raw-moment}
\end{equation}
标准化监督 $A$ 的界需再除以 $\eta^2$。
\end{proof}

\begin{theorem}[小步矩阈值]
固定 $r\ge1$、$0<\kappa\le1$。在所有有限正概率分布上取上确界，有
\begin{equation}
\sup_{q,T}\E_q|A|^r
=\Theta_{\kappa,r}(\eta^{\min\{r,1/\kappa\}-r}),\qquad
\sup_{q,T}\E_q|C_qA|^r
=\Theta_{\kappa,r}(\eta^{\min\{r,1/\kappa\}-r})
\end{equation}
当 $\eta\downarrow0$。$\kappa=0$ 对任何正 $\eta$ 都没有有限统一矩界。
\end{theorem}
\begin{proof}
先估计 $f$。负区域为 $O_{\kappa,r}(\eta^r)$。若 $r\kappa\le1$，用 $\log(1+x)\le x$ 和 $\E_q\rho^{r\kappa}\le1$ 得正区域也是 $O(\eta^r)$。若 $r\kappa>1$，令 $s=1/\kappa<r$。函数 $\log^r(1+x)/x^s$ 在零处按 $x^{r-s}$ 趋零，在无穷处也趋零，故有有限上界 $C_{r,s}$。于是
\begin{equation}
\E_q[|f|^r\mathbf1_{\rho\ge1}]
\le\kappa^{-r}C_{r,s}\eta^{1/\kappa}\E_q\rho
=O(\eta^{1/\kappa}).
\end{equation}
除以 $\eta^r$ 得未中心化上界。由 $|X-\E X|^r\le2^{r-1}(|X|^r+|\E X|^r)$ 及 Jensen，中心化矩不超过 $2^r\E|X|^r$。

当 $r\kappa\le1$，任意固定不同的二元正分布使 $A\to(\rho^\kappa-1)/\kappa$，其原始与中心化矩均为严格正常数，给出下界。当 $r\kappa>1$，取 $q_1=\eta^{1/\kappa}/2$、$T_1=1/2$。稀有动作有 $A_1=\log(2-\eta)/(\kappa\eta)$，另一动作 $A_2=O(1)$。所以 $q_1|A_1|^r=\Omega(\eta^{1/\kappa-r})$；中心化后的稀有值为 $(1-q_1)(A_1-A_2)$，具有同阶贡献。几何端点取固定非退化 $q$，令一个教师概率趋零，即使中心化，正质量动作上的对数差也无界。
\end{proof}

二阶矩阈值是 $\kappa\le1/2$，四阶矩阈值是 $\kappa\le1/4$。固定分布的小步 Taylor 展开不能替代这类允许分布随 $\eta$ 改变的最坏情形论证。

\section{误差相关方差与锐常数}
\begin{theorem}[误差相关的监督方差上界]
若 $0<\kappa\le1/2,0<\eta<1$，则
\begin{align}
V_{\kappa,\eta}&\le\frac{1-a_\kappa^2}{\kappa^2(1-\eta)^2},\qquad a_\kappa=\E_q(T/q)^\kappa,\\
V_{\kappa,\eta}&\le\frac{D_\kappa(T\Vert q)}{\kappa(1-\eta)^2}
\le\frac{K_T}{\kappa(1-\kappa)(1-\eta)^2}.
\end{align}
\end{theorem}
\begin{proof}
令 $a_\kappa=\E_q\rho^\kappa$。当 $0<\kappa\le1/2$，函数 $h(x)=\log(1-\eta+\eta x)/(\kappa\eta)$ 的 Lipschitz 常数不超过 $1/[\kappa(1-\eta)]$。独立同分布 $X,X'$ 满足 $\Var h(X)=\frac12\E[h(X)-h(X')]^2$，故
\begin{equation}
V_{\kappa,\eta}\le\frac{\Var_q(\rho^\kappa)}{\kappa^2(1-\eta)^2}
\le\frac{1-a_\kappa^2}{\kappa^2(1-\eta)^2}.
\end{equation}
第二步用 $\E_q\rho^{2\kappa}\le1$。

为得到与教师误差相关的界，令 $a(u)=\E_q\rho^u$。其二阶导数为 $\E_q[\rho^u\log^2\rho]\ge0$，所以对 $2\kappa\in[\kappa,1]$，
\begin{equation}
a(2\kappa)\le\frac{1-2\kappa}{1-\kappa}a(\kappa)+\frac{\kappa}{1-\kappa}.
\end{equation}
减去 $a(\kappa)^2$ 并因式分解：
\begin{equation}
\Var_q(\rho^\kappa)
\le\frac{(1-a_\kappa)[\kappa+(1-\kappa)a_\kappa]}{1-\kappa}
\le\frac{1-a_\kappa}{1-\kappa}.
\end{equation}
代入散度定义得到式~\eqref{eq:error-variance-main}。最后一步利用 $\log(1+\eta x)\le\eta x$ 得 $K_T\ge(1-a_\kappa)/\kappa$。
\end{proof}

\begin{theorem}[相对教师误差的锐小步方差常数]
设 $A_\kappa^0=(\rho^\kappa-1)/\kappa$、$K_T^0=-\E_qA_\kappa^0$。在至少二动作的正分布类中，对 $0<\kappa\le1/2$，
\begin{equation}
\sup_{q\ne T}\frac{\Var_qA_\kappa^0}{D_\kappa(T\Vert q)}=\frac1\kappa,\qquad
\sup_{q\ne T}\frac{\Var_qA_\kappa^0}{K_T^0}=\frac1{\kappa(1-\kappa)}.
\end{equation}
其余形状没有有限统一常数，故两种准则均由 $\kappa=1/2$ 唯一最小化。
\end{theorem}
\begin{proof}
在小步极限 $A^0_\kappa=(\rho^\kappa-1)/\kappa$，
\begin{equation}
\sup_{q\ne T}\frac{\Var_qA^0_\kappa}{D_\kappa(T\Vert q)}=\frac1\kappa,\qquad
\sup_{q\ne T}\frac{\Var_qA^0_\kappa}{K_T^0}=\frac1{\kappa(1-\kappa)}.
\end{equation}
上界已证。锐性取 $q=(v,1-v)$、$T=(v^2,1-v^2)$，此时 $\rho=(v,1+v)$。当 $v\downarrow0$，
\begin{equation}
a_\kappa=1-(1-\kappa)v+o(v),\quad
\Var_qA^0_\kappa=v/\kappa^2+o(v),\quad
D_\kappa=v/\kappa+o(v).
\end{equation}
两比值逼近所列常数。对 $\kappa>1/2$，取 $q=(v,1-v),T=(1/2,1/2)$，方差按 $v^{1-2\kappa}$ 发散；内部幂散度有界，$\kappa=1$ 的 KL 只按 $\log(1/v)$ 增长。$\kappa=0$ 用固定非对称 $q$ 与趋零教师概率，方差按对数平方增长，而 KL 按一次对数增长。故相对准则的唯一最优点为 $\kappa=1/2$。
\end{proof}

\section{绝对最坏方差的不同最优点}
\begin{theorem}[绝对最坏小步方差及其最优形状]
对 $0<\kappa<1/2$，令 $b=(1-2\kappa)/[2(1-\kappa)]$，则
\begin{equation}
\sup_{q,T}\Var_q A_\kappa^0
=\frac{b^{1-2\kappa}(1-b)}{\kappa^2}.
\end{equation}
在 $\kappa=1/2$，上确界为四。该绝对方差准则的唯一最优形状满足
$\log[(1-2\kappa)/(2(1-\kappa))]=-1/\kappa$，数值约为 $0.441928$。
\end{theorem}
\begin{proof}
对 $0<\kappa<1/2$，令 $X=\rho^\kappa$、$r=1/\kappa>2$，则 $\E_qX^r=1$。Hölder 给出
\begin{equation}
\E X^2\le(\E X)^{(r-2)/(r-1)}.
\end{equation}
具体地，将 $X^2=X^{(r-2)/(r-1)}(X^r)^{1/(r-1)}$，取共轭指数 $(r-1)/(r-2)$ 与 $r-1$ 即得。令 $a=\E X$、$\gamma=(1-2\kappa)/(1-\kappa)$，则 $\Var X\le a^\gamma-a^2$。最大值处 $a^{2-\gamma}=\gamma/2$，得到
\begin{equation}
\sup_{q,T}\Var A^0_\kappa
=C_{\rm vel}(\kappa)
=\frac{b^{1-2\kappa}(1-b)}{\kappa^2},
\qquad b=\frac{1-2\kappa}{2(1-\kappa)}.
\end{equation}
取 $q=(b,1-b)$、$T\to(1,0)$，Hölder 与驻点条件同时取等，证明锐性。$\kappa=1/2$ 的上确界为四。对最优 $b$ 使用包络求导，
\begin{equation}
\frac d{d\kappa}\log C_{\rm vel}=-2/\kappa-2\log b.
\end{equation}
唯一根满足 $\log[(1-2\kappa)/(2(1-\kappa))]=-1/\kappa$。左侧加 $1/\kappa$ 严格递减，端点分别趋于正、负无穷，故根唯一，数值为 $\kappa_\star\simeq0.441928$。这个准则最小化绝对最坏方差，不等于相对误差准则，也不等于任务表现。
\end{proof}

\section{平方根流、局部速率与噪声可求和性}
\begin{proposition}[平方根自然梯度流的显式亲和度]
固定正教师，沿 $\kappa=1/2$ 的自然梯度流，令 $a=\sum\sqrt{qT}$。则 $\|v_{1/2}\|_q^2=4(1-a^2)$、$\dot a=1-a^2$。若初值不同于教师，
$a(t)=\tanh[t+\operatorname{artanh}a(0)]$，且 $D(0)e^{-2t}\le D_{1/2}(T\Vert q(t))\le D(0)e^{-t}$。
\end{proposition}
\begin{proof}
平方根点令 $a=\sum\sqrt{qT}$。小步 Fisher 速度满足 $\|v_{1/2}\|_q^2=4(1-a^2)$，且 $D_{1/2}=4(1-a)$。因此沿连续自然梯度流，
\begin{equation}
\dot a=1-a^2,\qquad a(t)=\tanh[t+\operatorname{artanh}a(0)]
\end{equation}
对 $q\ne T$ 成立；在 $q=T$ 处流恒定。又有 $\dot D=-(1+a)D$，故 $D(0)e^{-2t}\le D(t)\le D(0)e^{-t}$。离散精确更新的 Hellinger 收缩由正文式~\eqref{eq:hellinger-contraction-main} 证明。
\end{proof}

\begin{proposition}[幂均值的二阶展开与共同局部率]
固定正教师、$\kappa\in[0,1],0<\eta<1$。令 $q=T+\epsilon h$、$\sum h=0$、$S_h=\sum h_a^2/T_a$，则
\begin{equation}
m_a=T_a+(1-\eta)\epsilon h_a+
\frac{(\kappa-1)(1-\eta)\eta}{2}\epsilon^2
\left(h_a^2/T_a-T_aS_h\right)+O(\epsilon^3).
\end{equation}
对固定非零 $h$ 及一处具有正二阶曲率的三阶光滑 $f$-散度，更新前后教师散度比为 $(1-\eta)^2+O(\epsilon)$。
\end{proposition}
\begin{proof}
固定 $T>0$，写 $q=T+\epsilon h$、$\sum h=0$、$S_h=\sum h_a^2/T_a$。两次二项式展开并除以归一化常数，得到
\begin{equation}
m_a=T_a+(1-\eta)\epsilon h_a
+\frac{(\kappa-1)(1-\eta)\eta}{2}\epsilon^2
\left(\frac{h_a^2}{T_a}-T_aS_h\right)+O(\epsilon^3).
\end{equation}
未归一化展开中二阶项为 $(\kappa-1)(1-\eta)\eta\epsilon^2h_a^2/(2T_a)$；其总和正是需要减去的归一化项。几何端点直接展开 $q^{1-\eta}T^\eta$ 给出相同结果。任何在一处二阶曲率为正的光滑 $f$-散度都满足局部比值 $(1-\eta)^2+O(\epsilon)$。
\end{proof}

\begin{corollary}[精确外层的监督噪声可求和性]
固定教师、$0<\kappa\le1/2,0<\eta<1$，每轮精确刷新。则
\begin{equation}
\sum_{n=1}^\infty V_n\le
\frac{\KL(T\Vert q_0)}{\kappa(1-\kappa)(1-\eta)^2\eta}.
\end{equation}
平方根点还满足 $V_n\le4(1-\eta)^nH^2(q_0,T)/(1-\eta)^2$。该式控制上界包络，不断言实际方差逐步单调。
\end{corollary}
\begin{proof}
最后，固定 $0<\kappa\le1/2,\eta<1$ 并精确刷新，式~\eqref{eq:center-sandwich} 与三点求和给出
\begin{equation}
\sum_{n=1}^\infty V_n
\le\frac{\KL(T\Vert q_0)}{\kappa(1-\kappa)(1-\eta)^2\eta}.
\end{equation}
平方根点可用更直接的 Hellinger 界：
\begin{equation}
V_n\le\frac{4H^2(q_n,T)}{(1-\eta)^2}
\le\frac{4(1-\eta)^nH^2(q_0,T)}{(1-\eta)^2}.
\end{equation}
上界包络下降，不保证实际方差逐步单调；所有这些结论固定同一前缀、教师和参数。
\end{proof}

\section{固定 Fisher 速度下的比较}
\begin{proposition}[共同损失下的最陡形状]
\label{prop:fisher-shape}
固定有限正分布 $q\ne T$，令 $\rho=T/q$、$r=\log\rho$，
$w_\kappa=C_q(\rho^\kappa)/\kappa$，并以连续延拓定义 $w_0=C_qr$。
向量 $v_\kappa=q w_\kappa$ 是小步目标的切向量。将其按 Fisher 范数
$\|v\|_q^2=\sum_a v_a^2/q_a$ 归一化为 $u_\kappa$，则
\begin{align}
-D\KL(q\Vert T)[u_\kappa]
 &=\frac{\E_q[rw_\kappa]}{\sqrt{\E_qw_\kappa^2}}
 \le\sqrt{\Var_qr},\\
-D\KL(T\Vert q)[u_\kappa]
 &=\frac{\E_q[\rho w_\kappa]}{\sqrt{\E_qw_\kappa^2}}
 \le\sqrt{\Var_q\rho}.
\end{align}
第一项由 $\kappa=0$ 达到，第二项由 $\kappa=1$ 达到；不声称达到者唯一。
\end{proposition}
\begin{proof}
切向量满足 $\sum_a v_a=0$。逐坐标微分两个 KL，
其负方向导数分别为 $\sum_a r_av_a$ 与 $\sum_a\rho_av_a$。
代入 $u_\kappa=q w_\kappa/\sqrt{\E_qw_\kappa^2}$ 得等式。
由于 $\E_qw_\kappa=0$，分子可分别写成
$\E_q[(C_qr)w_\kappa]$ 和 $\E_q[(C_q\rho)w_\kappa]$。
Cauchy--Schwarz 不等式给上界；在两个所述端点，相应中心化向量正比例，故取等。
这里统一的是 Fisher 速度和被比较的损失，与统一原始步长、最小化监督方差是不同的比较条件。
\end{proof}

\section{平方根目标与平方根拟合的联合界}
\begin{proposition}[初始拟合系数的统一二阶矩]
\label{prop:root-joint}
取 $\kappa=1/2$ 的目标 $m$，$0<\eta\le1$，并使用平方根散度拟合。
在初始点 $p=q$，令 $a=\sum_i\sqrt{q_iT_i}$、
$b=\sum_i\sqrt{q_im_i}$。若归一化拟合系数为
$G_i=2(b-\sqrt{m_i/q_i})/\eta$，则
\begin{equation}
\E_qG^2=\frac{4(1-a^2)}{Z}
\le\frac4{(1-\eta)^2+\eta^2}\le8,\qquad
Z=(1-\eta)^2+\eta^2+2\eta(1-\eta)a.
\end{equation}
若 score 范数至多为 $S$，相应单动作梯度的二阶矩至多为 $8S^2$。
\end{proposition}
\begin{proof}
由 $\sqrt{m_i}=[(1-\eta)\sqrt{q_i}+\eta\sqrt{T_i}]/\sqrt Z$，
有 $b=(1-\eta+\eta a)/\sqrt Z$。
又 $\E_q(m/q)=1$，所以 $\E_qG^2=4(1-b^2)/\eta^2$。
直接相减得到 $Z-(1-\eta+\eta a)^2=\eta^2(1-a^2)$。
因 $0<a\le1$，丢弃 $Z$ 的非负交叉项并用
$(1-\eta)^2+\eta^2\ge1/2$ 得界。最后逐点使用
$\|G_i\psi_i\|^2\le S^2G_i^2$。
此处拟合形状也是平方根；该系数不是 BPO 的对数残差。
\end{proof}

\begin{proposition}[有限步最坏方差的小步极限]
\label{prop:finite-variance-limit}
设 $0<\kappa\le1/2$，$V_{\kappa,\eta}(q,T)=\Var_q A_{\kappa,\eta}(T/q)$。
以前文的锐常数
$C_{\rm vel}(\kappa)=\sup_{q,T}\Var_q(\rho^\kappa)/\kappa^2$ 记最坏速度方差，则
\begin{equation}
\sup_{q,T}V_{\kappa,\eta}(q,T)
\le\frac{C_{\rm vel}(\kappa)}{(1-\eta)^2},\qquad
\lim_{\eta\downarrow0}\sup_{q,T}V_{\kappa,\eta}(q,T)
=C_{\rm vel}(\kappa).
\end{equation}
固定 $q,T$ 时，平方根端点的极限为
$4(1-a^2)=H^2(q,T)[4-H^2(q,T)]$，其中 $H^2=2-2a$。
\end{proposition}
\begin{proof}
将 $A$ 看成 $x=\rho^\kappa$ 的函数，其导数为
$1/[\kappa(1-\eta+\eta x)]\le1/[\kappa(1-\eta)]$。
对独立同分布的 $X,X'$，恒等式
$\Var h(X)=\E(h(X)-h(X'))^2/2$ 将 Lipschitz 界转为方差界。
固定正分布时逐点 Taylor 展开给
$A_{\kappa,\eta}(\rho)\to(\rho^\kappa-1)/\kappa$。
对任意 $\epsilon>0$，先选择一个速度方差超过
$C_{\rm vel}-\epsilon$ 的固定分布对，再令 $\eta\downarrow0$，
得到下极限至少为 $C_{\rm vel}-\epsilon$；上界给上极限不超过
$C_{\rm vel}$。最后令 $\epsilon\downarrow0$。
平方根公式由 $\E_q\rho=1$ 和 $\E_q\sqrt\rho=a$ 得到。
\end{proof}

\section{小 KL 与标准化不能替代矩分析}
\begin{proposition}[KL 趋零而监督方差发散]
\label{prop:small-kl-large-variance}
存在二元正分布族使 $\KL(q\Vert T)\to0$，但
$\Var_q\log(T/q)\to\infty$。也存在 KL 趋零而非负 KL 估计器
$\phi(\rho)=\rho-1-\log\rho$ 方差发散的族。
\end{proposition}
\begin{proof}
第一族取 $q=(\epsilon,1-\epsilon)$，
$T_1=\exp(-\epsilon^{-3/4})$、$T_2=1-T_1$。
第一坐标的 $r=\log(T/q)$ 为
$-\epsilon^{-3/4}-\log\epsilon$，第二坐标为 $O(\epsilon)$。
所以 $-\E_qr\sim\epsilon^{1/4}\to0$，
而二点方差公式给
$\Var_qr=\epsilon(1-\epsilon)(r_1-r_2)^2\sim\epsilon^{-1/2}$。
对固定正 $\kappa$ 和 $\eta<1$，同一族的 $A(\rho_1)$ 趋于有限下界，
$A(\rho_2)=O(\epsilon)$，方差为 $O(\epsilon)$。

第二族仍取 $q_1=\epsilon$，改为 $T_1=\epsilon^{1/4}$。
第一坐标对 KL 的贡献为 $O(\epsilon|\log\epsilon|)$，
第二坐标贡献渐近为 $\epsilon^{1/4}$。
此时 $\rho_1=\epsilon^{-3/4}$、
$\phi(\rho_1)\sim\epsilon^{-3/4}$，
而 $\rho_2=1-O(\epsilon^{1/4})$ 给 $\phi(\rho_2)=O(\epsilon^{1/2})$。
同一二点方差公式再次给 $\Var_q\phi\sim\epsilon^{-1/2}$。
两族说明，均值散度小不能单独控制单动作估计的二阶矩。
\end{proof}

\begin{proposition}[二元总体标准化消去信号形状]
\label{prop:binary-standardization}
设 $q=(v,1-v)$，$0<v<1$，两点信号差 $D=f_1-f_2\ne0$。
令 $\sigma^2=\Var_qf$，取第一 logit 的 score
$\psi=(1-v,-v)$。则
\begin{equation}
\Var_q\!\left[\psi\,\frac{C_qf}{\sigma}\right]=(1-2v)^2.
\end{equation}
任何严格递增的信号变换均保留逐点标准化值。
\end{proposition}
\begin{proof}
有 $C_qf=((1-v)D,-vD)$、
$\sigma=|D|\sqrt{v(1-v)}$。除以 $\sigma$ 后只剩
$\operatorname{sign}(D)$ 和 $v$。两个梯度取值之差为
$\operatorname{sign}(D)(1-2v)/\sqrt{v(1-v)}$，
乘以二点方差系数 $v(1-v)$ 得结论。
严格递增变换保留 $D$ 的符号。本命题使用总体方差、二元动作和单个 logit；
多动作、有限组经验标准化或共享参数不由此式涵盖。
\end{proof}

\section{共同评价指标下的收敛与稳定性}
\label{sec:shape-comparison}
本节固定有限正分布、同一教师及 $0<\eta<1$，比较精确目标映射。相同 $\eta$ 固定插值强度，不保证不同形状移动相同距离；若比较单位 Fisher 速度的瞬时下降，应使用命题~\ref{prop:fisher-shape}。这里的结论也不直接给出共享参数训练的用时排序。

\begin{theorem}[所有形状的正向 KL 乘法收缩]
\label{thm:forward-kl-contraction}
对 $0\le\kappa\le1$，记幂均值的未归一化坐标为 $h$、$Z=\sum_a h_a$。精确目标满足
\begin{equation}
\KL(T\Vert m_{\kappa,\eta})
\le(1-\eta)\KL(T\Vert q)+\log Z
\le(1-\eta)\KL(T\Vert q).
\end{equation}
因而固定教师、固定 $\eta$ 的精确迭代有 $\KL(T\Vert q_n)\le(1-\eta)^n\KL(T\Vert q_0)$。几何端点的第一步不等式取等。
\end{theorem}
\begin{proof}
设 $b=1-\eta$。对 $\kappa>0$，加权算术--几何均值不等式给 $bq_a^\kappa+\eta T_a^\kappa\ge q_a^{b\kappa}T_a^{\eta\kappa}$，所以 $h_a\ge q_a^bT_a^\eta$。幂均值不等式还给 $h_a\le bq_a+\eta T_a$，故 $Z\le1$。几何端点直接使用 $h_a=q_a^bT_a^\eta$，并由加权算术--几何均值不等式得到同一上界。于是
\begin{equation}
\log m_a\ge b\log q_a+\eta\log T_a-\log Z.
\end{equation}
两边从 $\log T_a$ 中减去、乘 $T_a$ 后求和，即得第一式；$\log Z\le0$ 给第二式。逐轮递推完成证明。这个共同的正向 KL 因子，不意味着所有形状对每个凸生成散度都有相同因子；算术端点的全散度乘法收缩是更强结论。
\end{proof}

\begin{lemma}[平方根目标的双锚点恒等式]
\label{lem:root-two-anchor}
固定正教师 $T$，设 $\eta>0$，两个锚点 $q,\widetilde q$ 的平方根未归一化向量均逐坐标为正。令 $b=1-\eta$、$w(q)=b\sqrt q+\eta\sqrt T$、$Z(q)=\|w(q)\|_2^2$，记 $M(q)=m_{1/2,\eta}(q,T)$。则
\begin{equation}
H^2(M(q),M(\widetilde q))
=\frac{b^2H^2(q,\widetilde q)
-[\sqrt{Z(q)}-\sqrt{Z(\widetilde q)}]^2}
{\sqrt{Z(q)Z(\widetilde q)}}.
\label{eq:root-two-anchor}
\end{equation}
\end{lemma}
\begin{proof}
正性保证 $\sqrt{M(q)}=w(q)/\|w(q)\|$。对两个非零向量 $w,x$，展开内积得到
\begin{equation}
\|w-x\|^2=(\|w\|-\|x\|)^2
+\|w\|\|x\|\left\|\frac{w}{\|w\|}-\frac{x}{\|x\|}\right\|^2.
\end{equation}
代入 $x=w(\widetilde q)$，使用 $w(q)-w(\widetilde q)=b(\sqrt q-\sqrt{\widetilde q})$，移项即得恒等式。平方根正性是该解释的条件，不能用负坐标平方后得到的另一映射代替。
\end{proof}

\begin{proposition}[插值区间的扰动收缩比较]
\label{prop:interpolation-stability}
固定 $0<\eta<1$、$b=1-\eta$。三种形状满足
\begin{align}
d_H(M_0(q),M_0(\widetilde q))&=b\,d_H(q,\widetilde q),\\
H(M_{1/2}(q),M_{1/2}(\widetilde q))&\le\frac b{\sqrt{b^2+\eta^2}}H(q,\widetilde q),\\
\TV(M_1(q),M_1(\widetilde q))&=b\,\TV(q,\widetilde q),\\
H(M_1(q),M_1(\widetilde q))&\le\sqrt b\,H(q,\widetilde q).
\end{align}
平方根的上述 Hellinger 上界常数不大于算术上界常数，当且仅当 $\eta\ge1/2$。
\end{proposition}
\begin{proof}
几何目标的 log-ratio 差为 $b\log(q/\widetilde q)$ 加公共常数，取跨度得到第一式。平方根目标的 $Z(q)=b^2+\eta^2+2b\eta\sum_a\sqrt{q_aT_a}\ge b^2+\eta^2$；在式~\eqref{eq:root-two-anchor} 中丢弃被减去的平方项并开方得到第二式。算术目标差精确等于 $b(q-\widetilde q)$，故 TV 等式成立。

对每个坐标，Cauchy--Schwarz 给
\begin{equation}
\sqrt{(bq_a+\eta T_a)(b\widetilde q_a+\eta T_a)}
\ge b\sqrt{q_a\widetilde q_a}+\eta T_a.
\end{equation}
求和并用 $H^2=2-2\sum\sqrt{pq}$ 得最后一式。比较两正上界常数并平方，$b^2/(b^2+\eta^2)\le b$ 等价于 $b\le b^2+\eta^2$，即 $\eta\ge1/2$。这里只比较已证明的统一界，不推断任意具体锚点对有相同排序。
\end{proof}

\begin{proposition}[平方根插值的角度递推]
\label{prop:root-interpolation-angle}
固定 $0<\eta<1$，令 $\theta=\arccos\sum_a\sqrt{q_aT_a}\in[0,\pi/2)$。目标的夹角满足
\begin{equation}
\theta^+=\arctan\frac{(1-\eta)\sin\theta}{\eta+(1-\eta)\cos\theta}.
\end{equation}
在 $\eta=1/2$ 时，$\theta^+=\theta/2$，因而 $H^2(q_n,T)=4\sin^2(\theta_0/2^{n+1})$。
\end{proposition}
\begin{proof}
当 $q\ne T$，写 $\sqrt q=\sqrt T\cos\theta+w\sin\theta$，其中 $w$ 为与 $\sqrt T$ 正交的单位向量。未归一化目标的两个正交系数分别为 $\eta+(1-\eta)\cos\theta$ 和 $(1-\eta)\sin\theta$。前者为正，二者之比即夹角正切。$\eta=1/2$ 时用 $\sin\theta/(1+\cos\theta)=\tan(\theta/2)$ 得精确二分；迭代后使用 $H^2=4\sin^2(\theta/2)$。$q=T$ 时连续延拓给零角度。
\end{proof}

\section{形状优劣依赖比较准则}
\label{sec:shape-criteria-summary}
三个端点分别在 log-ratio、平方根和概率坐标中插值。它们共享定理~\ref{thm:forward-kl-contraction} 的正向 KL 全局界与局部 $(1-\eta)^2$ 的 KL 比例，但只在算术端点直接得到所有凸生成散度的同因子收缩。教师覆盖系数分别为：几何端点无统一正下界，平方根至少 $\eta^2T_a$，算术至少 $\eta T_a$。

监督信号采用相同标准化 $A=f/\eta$。其小步极限依次为 $\log\rho$、$2(\sqrt\rho-1)$、$\rho-1$。平方根方差为 $4(1-c^2)\le4$，其中 $c=\sum\sqrt{qT}$；算术方差为 $\chi^2(T\Vert q)$，几何对数方差也无统一有限上界。固定正 $\eta<1$ 的算术信号仍有有限统一界，不能与其最坏方差在 $\eta\downarrow0$ 时发散混淆。这些矩结论及最优形状的不同准则已经在本附录证明；它们约束采样信号，不直接给出实际优化时间排序。

\begin{proposition}[共同 KL 指标不存在固定端点排序]
\label{prop:shape-order-counterexample}
在 $\eta=1/2$ 下，存在正二动作分布使三种形状的更新后逆向 KL 按 $1,1/2,0$ 从小到大排列，也存在按 $0,1/2,1$ 排列的分布。因此平方根不能在这个共同指标下逐实例支配另外两个端点。
\end{proposition}
\begin{proof}
给定 $q=(x,1-x)$、$T=(t,1-t)$，记目标第一坐标为 $p_\kappa$。三式分别为
\begin{align}
p_0&=\frac{\sqrt{xt}}{\sqrt{xt}+\sqrt{(1-x)(1-t)}},\\
p_{1/2}&=\frac{(\sqrt x+\sqrt t)^2}
{2+2[\sqrt{xt}+\sqrt{(1-x)(1-t)}]},\qquad p_1=\frac{x+t}{2}.
\end{align}
第一组 $x=1/20,t=1/5$ 代入后有 $0.10290<p_0<0.10291$、$0.11411<p_{1/2}<0.11412$、$p_1=0.125<t$。第二组 $x=3/5,t=19/20$ 有 $p_1=0.775$、$0.80678<p_{1/2}<0.80679$、$0.84223<p_0<0.84224<t$；这些有理数夹界由对正根式两边平方即可核对。二元 KL 对 $p$ 的导数是 $\log[p(1-t)/(t(1-p))]$，在 $p<t$ 时严格为负，故两组概率次序给出相反的严格 KL 排序。代入二元 KL 后的近似值如下。
\begin{center}
\begin{tabular}{ccrrr}
\toprule
$x$ & $t$ & $\kappa=0$ & $\kappa=1/2$ & $\kappa=1$\\
\midrule
$0.05$ & $0.2$ & $0.034382$ & $0.026311$ & $0.019660$\\
$0.6$ & $0.95$ & $0.079878$ & $0.129358$ & $0.180628$\\
\bottomrule
\end{tabular}
\end{center}
两个严格反向排序否定固定排序；相同 $\eta$ 并没有使各形状移动相同距离，故该例不替代单位 Fisher 速度的命题~\ref{prop:fisher-shape}。
\end{proof}

在可实现教师附近，有限 BPO、MDPO 的共同参数谱由定理~\ref{thm:joint-spectrum} 描述。平方根的采样矩控制不能改写为改善了固定参数 Fisher 矩阵的最小特征值。选择形状时须分别说明评价散度、动作积分方式、内层预算和优化器；这些条件不同，最优性判断也不同。


\chapter{信息访问与条件积分估计}
\label{app:estimation}
\section{四个积分对象与可观测信息}
\begin{proposition}[双向 KL 的数值与梯度积分]
有限词表上，$P=P_\theta>0$ 可微且归一化，$Q>0$ 冻结。令 $\ell=\log(P/Q)$、$\psi=\nabla_\theta\log P$，则
\begin{equation}
\KL(P\Vert Q)=\E_P\ell,\quad \nabla_\theta\KL(P\Vert Q)=\E_P(\ell\psi),\quad
\nabla_\theta\KL(Q\Vert P)=-\E_Q\psi.
\end{equation}
\end{proposition}
\begin{proof}
本附录局部记当前学生为 $P=P_\theta$、冻结参考为 $Q$、动作提案为 $\beta$；三者职责不同。令 $\ell_a=\log(P_a/Q_a)$、$\psi_a=\nabla_\theta\log P_a$。双向 KL 及其梯度为
\begin{equation}
K=\E_P\ell,\quad g=\E_P[\ell\psi],\qquad
K^{\rm F}=-\E_Q\ell,\quad g^{\rm F}=-\E_Q\psi.
\label{eq:four-integrals}
\end{equation}
$g$ 的求导先产生 $\E_P[(\ell+1)\psi]$，再利用 $\E_P\psi=0$。参考为锚点时可代入 $Q=q$；BPO 学生中心 $\KL(q\Vert p)$ 对应这里的正向积分。
\end{proof}

完整词表概率允许精确有限求和。头部查询只给出集合 $S$ 上的 $Q_a$ 和尾部总质量；实际动作查询再给出该动作的一项概率。若 $S$ 是教师真正的 Top-$K_{\rm head}$，未知尾部还满足排名约束 $Q_a\le\min_{b\in S}Q_b$。接口是否确定某个量，应检查它在全部符合这些观测的分布上是否相同。

只保留桶质量时，不能恢复桶内形状。若实际动作落在尾部，其额外概率观测必须计入信息；也不能把一次尾部查询当成全部尾部已知。固定支持集的微分与每步重新选 Top-K 的程序是不同对象，以下推导均冻结分区。

\section{粗粒化、下界与不可识别性}
\begin{theorem}[分区链式分解与细化单调性]
设 $\mathcal C$ 为固定词表分区，$P,Q>0$，桶质量与条件分布分别为 $P_C,Q_C,P^C,Q^C$。则
\begin{equation}
\KL(P\Vert Q)=\sum_CP_C\log(P_C/Q_C)+\sum_CP_C\KL(P^C\Vert Q^C).
\end{equation}
交换 $P,Q$ 给另一方向；细化分区使粗粒度 KL 非减，完整 singleton 分区恢复完整 KL。
\end{theorem}
\begin{proof}
设 $\mathcal C$ 为词表分区，定义 $P_C=\sum_{a\in C}P_a$、$Q_C=\sum_{a\in C}Q_a$、$\ell_C=\log(P_C/Q_C)$、$P^C_a=P_a/P_C$、$Q^C_a=Q_a/Q_C$。逐点有
\begin{equation}
\ell_a=\ell_C+\log(P_a^C/Q_a^C).
\end{equation}
乘 $P_a$，先桶内求和再跨桶求和，得到
\begin{equation}
\KL(P\Vert Q)=K_{\mathcal C}+\sum_CP_C\KL(P^C\Vert Q^C),
\qquad K_{\mathcal C}=\sum_CP_C\ell_C.
\label{eq:partition-chain}
\end{equation}
交换 $P,Q$ 给出正向版本。细化分区时，新旧粗化 KL 的差为每个旧桶内细桶标签 KL 的加权和，故非负。全 singleton 分区恢复完整值。
\end{proof}

\begin{proposition}[桶质量信息下的 KL 下包络]
固定 $P>0$ 与各桶正质量 $Q_C$，仅施加桶质量约束时，$\inf_Q\KL(P\Vert Q)=\sum_CP_C\log(P_C/Q_C)$，唯一最优补全为 $Q_a=Q_CP_a^C$。交换 KL 方向仍由同一补全达到相应粗粒度下界。加入 Top-K 排名后，该下界不一定仍可达到。
\end{proposition}
\begin{proof}
在只知道 $Q_C$ 的信息类上，$K_{\mathcal C}$ 是完整 KL 的锐下包络，由 $Q_a=Q_CP_a^C$ 唯一达到；两个方向同时成立。加入真实 Top-K 排名后该构造可能不可行。例如 $P=(0.1,0.8,0.1)$，已知 $Q_1=0.4,Q_2+Q_3=0.6$ 且动作一是 Top-1；质量约束最优构造给 $Q_2=8/15>0.4$，违反排名。因此 Tail-aware 仍是下界，但不必是完整信息类上的锐下界。
\end{proof}

\begin{proposition}[稀疏头部信息不控制完整逆向 KL 上界]
当至少两个尾部概率仍未知时，存在具有相同真实 Top-1 概率和尾部质量的正参考分布族，使对同一个固定 $P$ 的 $\KL(P\Vert Q)$ 无上界。
\end{proposition}
\begin{proof}
相同头部信息也不控制逆向 KL 上界。固定 $P=(0.4,0.2,0.4)$ 与 $Q_\epsilon=(0.7,\epsilon,0.3-\epsilon)$、$0<\epsilon<0.15$，Top-1 和尾部质量完全相同，而 $0.2\log(0.2/\epsilon)\to\infty$。该反例依赖至少两个仍未知的尾部坐标。
\end{proof}

\begin{proposition}[粗粒度 KL 的梯度残差]
固定分区和参考 $Q$，令 $\ell_C=\log(P_C/Q_C)$、$\bar\psi_C=\E_{P^C}\psi$，则
\begin{align}
\nabla K_{\mathcal C}&=\sum_CP_C\ell_C\bar\psi_C,\\
g-\nabla K_{\mathcal C}
&=\sum_CP_C\{\Cov_{P^C}(\ell-\ell_C,\psi)+\KL(P^C\Vert Q^C)\bar\psi_C\}.
\end{align}
标量粗化误差的单调改进不一般推出梯度偏差范数的单调改进。
\end{proposition}
\begin{proof}
记 $\bar\psi_C=\nabla\log P_C=\E_{P^C}\psi$。对粗化目标求导：
\begin{align}
\nabla K_{\mathcal C}
&=\sum_C(\ell_C+1)\nabla P_C
=\sum_CP_C\ell_C\bar\psi_C,\\
g-\nabla K_{\mathcal C}
&=\sum_CP_C\{\Cov_{P^C}(\ell-\ell_C,\psi)
+\KL(P^C\Vert Q^C)\bar\psi_C\}.
\end{align}
最后一步对桶内乘积期望分解。标量误差只含条件 KL，梯度误差还依赖 score 相关性；标量下界改善不推出梯度偏差范数单调改善。
\end{proof}

\section{精确积分加残差抽样}
\begin{theorem}[统一估计与协方差]
设目标 $I_\nu(h)=\sum_a\nu_ah_a$，$h_a$ 可为向量。选择可精确积分的代理 $b_a$，令 $R_a=h_a-b_a$。若 $\nu_a\|R_a\|>0$ 时 $\beta_a>0$，则对 $A\sim\beta$，
\begin{equation}
\widehat I=I_\nu(b)+\frac{\nu_A}{\beta_A}R_A
\end{equation}
无偏，且
\begin{equation}
\Cov(\widehat I)=\sum_{\beta_a>0}\frac{\nu_a^2}{\beta_a}R_aR_a^\top
-I_\nu(R)I_\nu(R)^\top.
\label{eq:master-cov}
\end{equation}
\end{theorem}
\begin{proof}
随机修正的均值为 $\sum_{\beta_a>0}\nu_aR_a=I_\nu(R)$，支持条件使缺失动作贡献为零；与确定积分相加即得目标。随机项外积的期望为式~\eqref{eq:master-cov} 的求和项，减去均值外积得到协方差。取迹得到向量方差，无偏性使其等于均方误差。$M$ 个条件独立样本的平均协方差除以 $M$。
\end{proof}

\begin{proposition}[积分项、权重和支持的统一偏差]
记 $R=h-b$、$\mathcal D=\{a:\beta_a>0\}$。若实际估计器为 $I_\nu(b)+\delta+\widetilde w_AR_A$，则其偏差为
\begin{equation}
\delta+\sum_{\mathcal D}(\beta_a\widetilde w_a-\nu_a)R_a
-\sum_{\mathcal D^c}\nu_aR_a.
\end{equation}
\end{proposition}
\begin{proof}
若确定积分有误差 $\delta$，使用错误权重 $\widetilde w$ 或缺失支持 $\mathcal D^c$，偏差恰为
\begin{equation}
\delta+\sum_{a\in\mathcal D}(\beta_a\widetilde w_a-\nu_a)R_a
-\sum_{a\notin\mathcal D}\nu_aR_a.
\end{equation}
这是对实际期望与目标直接相减的结果。精确覆盖的零残差动作可以没有采样支持；未知残差则不能由有限权重恢复。
\end{proof}

\begin{proposition}[固定残差下的迹方差最优提案]
固定代理与残差 $R$，若 $L_R=\sum_a\nu_a\|R_a\|>0$，则最优提案为 $\beta_a^\star=\nu_a\|R_a\|/L_R$，最小迹方差为 $L_R^2-\|I_\nu(R)\|^2$。若 $L_R=0$，目标可确定性精确返回。
\end{proposition}
\begin{proof}
固定代理时，迹方差最优提案为
\begin{equation}
\beta_a^\star=\frac{\nu_a\|R_a\|}{\sum_b\nu_b\|R_b\|},\quad
V_{\min}=\left(\sum_a\nu_a\|R_a\|\right)^2-\|I_\nu(R)\|^2.
\end{equation}
证明对 $\sum_a[\nu_a\|R_a\|/\sqrt{\beta_a}]\sqrt{\beta_a}$ 用 Cauchy--Schwarz；取比例提案时等号成立。全部残差为零时直接返回确定积分。这个 oracle 提案通常需要未知尾部值，不能仅凭公式称其可实现。
\end{proof}

\begin{theorem}[多控制的 Gram 最优解]
设 $U_j$ 为二阶可积零均值控制，估计器为 $X-\sum_jc_jU_j$。令 $\Gamma_{ij}=\E\langle U_i,U_j\rangle$、$d_j=\E\langle U_j,X-\E X\rangle$。则全部迹方差最优系数为 $\Gamma^\dagger d+\ker\Gamma$，相对最优值的方差代价为 $(c-\Gamma^\dagger d)^\top\Gamma(c-\Gamma^\dagger d)$。
\end{theorem}
\begin{proof}
若有多个已知积分控制 $h_j$，令 $X=(\nu_A/\beta_A)h_A$、$U_j=(\nu_A/\beta_A)h_{j,A}-I_\nu(h_j)$、$\Gamma_{ij}=\E\langle U_i,U_j\rangle$、$d_j=\E\langle U_j,X-\E X\rangle$。估计器 $X-\sum_jc_jU_j$ 的方差为
\begin{equation}
V_0-2c^\top d+c^\top\Gamma c.
\end{equation}
若 $z\in\ker\Gamma$，则 $\sum_jz_jU_j=0$ 几乎处处，所以 $z^\top d=0$。因此 $d\in\operatorname{range}\Gamma$，最优系数为 $\Gamma^\dagger d+\ker\Gamma$，超额方差为 $(c-\Gamma^\dagger d)^\top\Gamma(c-\Gamma^\dagger d)$。这最小化协方差迹，不一般同时最小化每个方向。
\end{proof}

\section{头尾估计器作为统一定理的实例}
\begin{proposition}[三种确定性头部目标及其偏差]
固定非空头尾分区，令 $H=\sum_SP_a\ell_a$、$l=\log(\tau_P/\tau_Q)$、$B=\tau_P\KL(P^{\rm tail}\Vert Q^{\rm tail})$。完整逆向 KL 为 $H+\tau_Pl+B$。Truncated 与 Tail-aware 的偏差分别为 $-\tau_Pl-B,-B$；Normalized 优化头部条件 KL $H/P_S+\log(Q_S/P_S)$。
\end{proposition}
\begin{proof}
令 $\tau_P=P(S^c)$、$\tau_Q=Q(S^c)$、$J_A=\mathbf1\{A\notin S\}$，并定义
\begin{equation}
H=\sum_{a\in S}P_a\ell_a,\quad
l=\log(\tau_P/\tau_Q),\quad
B=\tau_P\KL(P^{\rm tail}\Vert Q^{\rm tail}).
\end{equation}
则 $K=H+\tau_Pl+B$。Truncated 返回 $H$，偏差为 $-\tau_Pl-B$；Tail-aware 返回 $H+\tau_Pl$，偏差为 $-B$；Normalized 返回
\begin{equation}
K_{\rm norm}=H/P_S+\log(Q_S/P_S),
\end{equation}
它是头部条件分布的 KL。后者改变积分权重，不能当成只删除尾部项。
\end{proof}

\begin{proposition}[标量尾部控制族的精确方差]
令 $A\sim P$，控制 $c$ 在抽样前固定。估计器 $\widehat K(c)=H+\tau_Pc+J_A(\ell_A-c)$ 无偏，且
\begin{equation}
\Var\widehat K(c)=\tau_P\Var_{\rm tail}\ell+P_S\tau_P(\E_{\rm tail}\ell-c)^2.
\end{equation}
因此最优常数为尾部条件均值；HETS 与桶残差校正分别取 $c=0,l$。
\end{proposition}
\begin{proof}
在 $A\sim P$ 下，标量控制族
\begin{equation}
\widehat K(c)=H+\tau_Pc+J_A(\ell_A-c)
\end{equation}
无偏。设 $\mu_\ell=\E_{P^{\rm tail}}\ell$、$V_\ell=\Var_{P^{\rm tail}}\ell$，全方差公式给出
\begin{equation}
\Var\widehat K(c)=\tau_PV_\ell+P_S\tau_P(\mu_\ell-c)^2.
\end{equation}
头部精确、尾部采样（Head-Exact Tail-Sampled，HETS）取 $c=0$，桶残差修正取 $c=l$，oracle 取 $c=\mu_\ell$。残差修正优于 HETS 当且仅当 $|\mu_\ell-l|<|\mu_\ell|$。
\end{proof}

\begin{proposition}[向量尾部控制族的协方差]
令 $h_g=\sum_SP_a\ell_a\psi_a$、$\mu_g=\E_{\rm tail}(\ell\psi)$，固定向量 $b$。则 $\widehat g(b)=h_g+\tau_Pb+J_A(\ell_A\psi_A-b)$ 对完整梯度无偏，协方差为
\begin{equation}
\tau_P\Cov_{\rm tail}(\ell\psi)+P_S\tau_P(\mu_g-b)(\mu_g-b)^\top.
\end{equation}
\end{proposition}
\begin{proof}
对梯度令 $h_g=\sum_SP_a\ell_a\psi_a$、$\mu_g=\E_{P^{\rm tail}}(\ell\psi)$、$\bar\psi_{\rm tail}=\nabla\log\tau_P$、$b_{\rm TA}=l\bar\psi_{\rm tail}$。统一向量控制为
\begin{equation}
\widehat g(b)=h_g+\tau_Pb+J_A(\ell_A\psi_A-b),
\end{equation}
\begin{equation}
\Cov\widehat g(b)=\tau_P\Cov_{\rm tail}(\ell\psi)
+P_S\tau_P(\mu_g-b)(\mu_g-b)^\top.
\end{equation}
HETS 为 $b=0$，ResOPD 为 $b=b_{\rm TA}$。常向量只改变头尾事件间的均值差，不改变尾部内部协方差。
\end{proof}

\begin{proposition}[正交双控制的最优值与排序]
on-policy 抽样下，取 $\chi=l(\psi-\bar\psi_{\rm tail})$。估计器 $\widehat g(b,\vartheta)=h_g+\tau_Pb+J_A(\ell_A\psi_A-b-\vartheta\chi_A)$ 无偏。若 $V_\chi>0$，最优值为 $b=\mu_g,\vartheta=D_\chi/V_\chi$，超额方差为
\begin{equation}
P_S\tau_P\|b-\mu_g\|^2+\tau_PV_\chi(\vartheta-D_\chi/V_\chi)^2.
\end{equation}
其中 $V_\chi=\E_{\rm tail}\|\chi\|^2$、$D_\chi=\E_{\rm tail}\langle\ell\psi-\mu_g,\chi\rangle$。SC-TA 优于 ResOPD 当且仅当 $2D_\chi>V_\chi$。
\end{proposition}
\begin{proof}
再加入尾部条件零均值方向 $\chi_a=l(\psi_a-\bar\psi_{\rm tail})$，定义
\begin{equation}
\widehat g(b,\vartheta)=h_g+\tau_Pb+J_A(\ell_A\psi_A-b-\vartheta\chi_A).
\end{equation}
尾部内部与事件控制正交，因为 $\E_{\rm tail}\chi=0$。因此最优 $b=\mu_g$，最优 $\vartheta=D_\chi/V_\chi$，其中
\begin{equation}
V_\chi=\E_{\rm tail}\|\chi\|^2,\qquad
D_\chi=\E_{\rm tail}\langle\ell\psi-\mu_g,\chi\rangle.
\end{equation}
配方给出精确代价
\begin{equation}
V(b,\vartheta)-V_{\min}
=P_S\tau_P\|b-\mu_g\|^2
+\tau_PV_\chi(\vartheta-D_\chi/V_\chi)^2.
\end{equation}
$V_\chi=0$ 时内部控制不起作用。SC-TA 对应 $b=b_{\rm TA},\vartheta=1$；它比 ResOPD 方差小当且仅当 $2D_\chi>V_\chi$。
\end{proof}

\begin{proposition}[精确头部不保证逐输入降方差]
对向量 integrand $u$、$I=\E_Pu$、$H_u=\sum_SP_au_a$，有
\begin{equation}
\Var(u_A)-\Var(H_u+J_Au_A)
=\sum_SP_a\|u_a\|^2-2\langle I,H_u\rangle+\|H_u\|^2.
\end{equation}
右侧可以为负；真实条件期望的 Rao--Blackwell 化则保证方差不增。
\end{proposition}
\begin{proof}
即使头部精确，方差也不必比直接抽样小。对任意 integrand $u$，令 $I=\E_Pu$、$H_u=\sum_SP_au_a$，直接展开二阶矩：
\begin{equation}
\Var(u_A)-\Var(H_u+J_Au_A)
=\sum_SP_a\|u_a\|^2-2\langle I,H_u\rangle+\|H_u\|^2.
\end{equation}
右侧可负。数值反例 $P=(0.18,0.819,0.001)$、$Q=(0.08,0.67,0.25)$、$S=\{2\}$（按一开始编号），直接 log-ratio 方差约为 $0.088913$，HETS 为 $0.129131$。真实 Rao--Blackwell 化应取条件期望 $\E_P[\ell\mid C]=\ell_C+\KL(P^C\Vert Q^C)$；只知道桶质量时缺少第二项。
\end{proof}

\section{提案变化、分层采样与正向 KL}
\begin{proposition}[一般提案下的最优尾部常控制]
提案 $\beta$ 覆盖尾部。令 $W=\sum_{S^c}P_a^2/\beta_a$、$U=\sum_{S^c}(P_a^2/\beta_a)h_a$、$T_h=\sum_{S^c}P_ah_a$。若 $W>\tau_P^2$，最优控制为 $(U-\tau_PT_h)/(W-\tau_P^2)$；若等号成立，$\beta=P^{\rm tail}$，控制逐样本抵消，$M$ 次独立抽样的方差为 $\tau_P^2\Var_{\rm tail}(h)/M$。
\end{proposition}
\begin{proof}
从 $\beta$ 采样时，所有残差改为 $J_A(P_A/\beta_A)(h_A-b)$。二阶矩权重是 $P^2/\beta$，所以原有 on-policy 最优控制通常改变。令
\begin{equation}
W=\sum_{S^c}P_a^2/\beta_a,\quad U=\sum_{S^c}(P_a^2/\beta_a)h_a,\quad
T_h=\sum_{S^c}P_ah_a,
\end{equation}
展开关于 $b$ 的二次式，最优 $b=(U-\tau_PT_h)/(W-\tau_P^2)$。若 $W=\tau_P^2$，Cauchy--Schwarz 等号要求 $\beta=P^{\rm tail}$，控制项逐样本抵消。此时直接尾部条件采样的方差为 $\tau_P^2\Var_{\rm tail}(h)/M$，不超过同样本数的 on-policy 头尾估计方差。
\end{proof}

\begin{proposition}[分层预算的最优分配]
各桶独立采样 $M_C$ 次，$\sigma_C^2=\Var_{P^C}h$。总估计方差为 $\sum_CP_C^2\sigma_C^2/M_C$。在连续预算 $\sum_CM_C=M$、$\sum_CP_C\sigma_C>0$ 下，最优分配为 $M_C\propto P_C\sigma_C$。控制若由数据估计，条件无偏结论还要求控制拟合与评估动作之间满足对应独立性。
\end{proposition}
\begin{proof}
更一般，分区内独立采样 $M_C$ 次，方差为 $\sum_CP_C^2\sigma_C^2/M_C$。固定连续预算 $\sum_CM_C=M$，Cauchy--Schwarz 给出最优分配 $M_C\propto P_C\sigma_C$；零方差桶直接返回常值，整数预算需另行取整。控制拟合与评估共用随机动作时，$\E U=0$ 不保证 $\E[c(A)U]=0$；应条件化于独立校准数据，或使用交叉拟合并保留折间相关性。
\end{proof}

\begin{proposition}[正向 KL 的头尾校正]
令 $H^{\rm F}=-\sum_SQ_a\ell_a$，则 $K^{\rm F}=H^{\rm F}-\tau_Ql+\tau_Q\KL(Q^{\rm tail}\Vert P^{\rm tail})$。若 $\beta$ 覆盖未知 $Q$ 尾部，则 $H^{\rm F}+\tau_Qc+J_A(Q_A/\beta_A)(-\ell_A-c)$ 无偏。粗粒度正向 KL 的真实梯度为 $-\sum_CQ_C\bar\psi_C$。
\end{proposition}
\begin{proof}
正向 KL 使用同一积分定理取 $\nu=Q,h=-\ell$ 或 $-\psi$。令 $H^{\rm F}=-\sum_SQ_a\ell_a$，则
\begin{equation}
K^{\rm F}=H^{\rm F}-\tau_Ql+\tau_Q\KL(Q^{\rm tail}\Vert P^{\rm tail}),
\end{equation}
无偏标量控制为 $H^{\rm F}+\tau_Qc+J_A(Q_A/\beta_A)(-\ell_A-c)$。其支持要求覆盖未知的 $Q$ 尾部。直接微分粗化正向 KL 得 $-\sum_CQ_C\bar\psi_C$，其中 $\bar\psi_C$ 仍是 $P^C$ 下的 score 均值。
\end{proof}

\section{数值估计与反传的独立核验}
\begin{proposition}[固定样本反传与积分导数]
设 $\phi$ 可微，$A\sim P_\theta$、$c_A=P_\theta(A)/\sg(P_\theta(A))$。固定样本反传满足
\begin{equation}
\nabla\phi(\ell_A)=\phi'(\ell_A)\psi_A,\qquad
\nabla[c_A\phi(\ell_A)]=[\phi(\ell_A)+\phi'(\ell_A)]\psi_A.
\end{equation}
第二式的期望为 $\nabla\E_P\phi(\ell)$。因此数值无偏与裸反传梯度无偏是不同条件；下表给出 K1--K4 与 K3$^{++}$ 的逐样本结果。
\end{proposition}
\begin{proof}
令 $A\sim P$，$\ell=\ell_A$，并定义 $c_A=P_\theta(A)/\sg(P_\theta(A))$。固定动作求导有
\begin{equation}
\nabla\phi(\ell)=\phi'(\ell)\psi,\qquad
\nabla[c_A\phi(\ell)]=[\phi(\ell)+\phi'(\ell)]\psi.
\end{equation}
第二式的期望由式~\eqref{eq:measure-derivative} 恰好恢复积分导数。数值与反传分类如下，均先讨论无截断的解析对象：
\begin{center}
\small
\begin{tabular}{lll}
\toprule
程序 & 前向值 & 固定样本反传\\
\midrule
K1 & $\ell$ & $\psi$\\
K2 & $\ell^2/2$ & $\ell\psi$\\
K3 & $e^{-\ell}+\ell-1$ & $(1-e^{-\ell})\psi$\\
K4：$c_A\sg(\ell)$ & $\ell$ & $\ell\psi$\\
live K4：$c_A\ell$ & $\ell$ & $(\ell+1)\psi$\\
K3$^{++}$：$c_A(e^{-\ell}+\ell-1)$ & $e^{-\ell}+\ell-1$ & $\ell\psi$\\
\bottomrule
\end{tabular}
\end{center}
K1、K3 的数值对逆向 KL 无偏，因为 $\E_Pe^{-\ell}=1$；K2 一般有偏，其反传却是逆向 KL 的无偏 score 梯度。K3 的裸反传期望是 $-\E_Q\psi$，对应正向 KL。K4 与 K3$^{++}$ 的逐样本反传相同，前向方差不同。Top-K 和相应 score 构造的外部背景见 EMA-PG\footfullcite{ema2026}；这里按明确公式区分命名。
\end{proof}

\begin{proposition}[K5 的双重无偏性]
完整正支持且 $A\sim P$ 时，$K5=-\sg(e^{-\ell_A})\ell_A+\log P_A-\sg(\log P_A)$ 的前向期望为 $\KL(Q\Vert P)$，固定样本反传期望为其真实参数梯度。
\end{proposition}
\begin{proof}
若需要正向 KL 的数值与 score 反传同时匹配，可取
\begin{equation}
K5=-\sg(e^{-\ell})\ell+\log P_A-\sg(\log P_A).
\end{equation}
其前向为 $-e^{-\ell}\ell$，期望为 $K^{\rm F}$；反传为 $(1-e^{-\ell})\psi$，期望为 $g^{\rm F}$。正支持和实际采样测度是这些恒等式的必要条件。
\end{proof}

\begin{proposition}[确定性头部导数与头尾常数匹配]
固定支持 $S$，令 $d_S=\nabla P_S$，则
\begin{equation}
\nabla H=h_g+d_S,\quad \nabla K_{\rm TA}=h_g-ld_S,\quad
\nabla K_{\rm norm}=h_g/P_S-Hd_S/P_S^2.
\end{equation}
对任意固定 $\upsilon$，$h_g+\upsilon d_S+J_A(\ell_A+\upsilon)\psi_A$ 无偏；只在头部或尾部保留该常数会留下对应的 $\pm d_S$ 偏差。
\end{proposition}
\begin{proof}
确定性头部数值的真实导数也应保留概率权重项。令 $d_S=\nabla P_S$，则
\begin{equation}
\nabla H=h_g+d_S,\quad
\nabla(H+\tau_Pl)=h_g-ld_S,\quad
\nabla K_{\rm norm}=\frac{h_g}{P_S}-\frac{H}{P_S^2}d_S.
\end{equation}
最后一式对商 $H/P_S$ 求导并与 $-\log P_S$ 的导数抵消得到，分母明确为标量 $P_S^2$。

头尾 score 常数必须匹配：$h_g+\upsilon d_S+J_A(\ell_A+\upsilon)\psi_A$ 对任意固定 $\upsilon$ 无偏，因为 $\E_PJ_A\psi_A=-d_S$。只在头部或尾部保留一个常数项，会分别留下 $\pm d_S$ 偏差。
\end{proof}

\begin{proposition}[固定样本势与非保守更新场]
冻结参考、头部和抽样动作。HETS 与 ResOPD 的场分别是 $H-P_S+J_A\ell_A^2/2$ 和 $K_{\rm TA}+J_A(\ell_A^2-l^2)/2$ 的普通梯度。SC-TA 的场一般非保守，其交叉偏导差为
\begin{equation}
-J_A[(\partial_i l)\psi_j-(\partial_j l)\psi_i].
\end{equation}
逐样本的保守性与总体梯度无偏性不相互替代。
\end{proposition}
\begin{proof}
可以用 $\sg(F)+L-\sg(L)$ 让前向显示 $F$、反向使用 $L$ 的导数；这不意味着前向的普通导数就是更新场。HETS 与 ResOPD 的固定样本势分别为
\begin{equation}
H-P_S+\tfrac12J_A\ell_A^2,\qquad
K_{\rm TA}+\tfrac12J_A(\ell_A^2-l^2).
\end{equation}
SC-TA 的场 $\nabla K_{\rm TA}+J_A(\ell_A-l)\nabla\ell_A$ 一般非保守，其交叉偏导差是
\begin{equation}
-J_A[(\partial_i l)\psi_j-(\partial_j l)\psi_i].
\end{equation}
取三动作 softmax、$S=\{1\},A=2$，$(i,j)=(1,2)$ 时该差为 $P_1P_3/(1-P_1)>0$。总体期望仍可为真实梯度，因为参数依赖采样测度的导数补上了额外项。
\end{proof}

\section{配分函数估计与可控展开}
\begin{proposition}[对数配分函数的 plug-in 偏差]
若 $\widehat Z>0$、$\E\widehat Z=Z>0$、$\E|\log\widehat Z|<\infty$，则 $\E\log\widehat Z\le\log Z$，非退化时严格。这个结论不排除使用其他随机构造无偏估计 $\log Z$。
\end{proposition}
\begin{proof}
固定 $q,f$，$Z=\E_qe^f$。精确头部加尾部样本均值无偏估计 $Z$，但若 $\widehat Z>0$、$\E|\log\widehat Z|<\infty$，Jensen 给
\begin{equation}
\E\log\widehat Z\le\log Z,
\end{equation}
非退化时严格。这个结论只针对 plug-in，不排除其他信息条件下无偏估计 $\log Z$。例如有限词表完全覆盖事件 $E$ 的概率 $c>0$ 已知时，$\mathbf1_E\log Z/c$ 是无偏构造，方差为 $(c^{-1}-1)(\log Z)^2$，但通常成本很高。
\end{proof}

\begin{theorem}[log-of-mean 的偏差与方差展开]
固定有限正分布与头部，$\widehat Z_M$ 为有界独立尾部变量均值加确定头部，具有与 $M$ 无关的正下界。若尾部单样本方差为 $V$，则
\begin{align}
\E\log\widehat Z_M-\log Z&=-V/(2MZ^2)+O(M^{-3/2}),\\
\Var\log\widehat Z_M&=V/(MZ^2)+O(M^{-3/2}).
\end{align}
\end{theorem}
\begin{proof}
对固定有限正分布，设 $\widehat Z_M$ 是有界独立尾部变量的均值加确定头部，且有与 $M$ 无关的正下界。令 $\epsilon=\widehat Z_M-Z$，尾部单样本方差为 $V$。则
\begin{equation}
\E\epsilon=0,\quad \E\epsilon^2=V/M,\quad
\E\epsilon^4=O(M^{-2}),\quad \E|\epsilon|^3=O(M^{-3/2}).
\end{equation}
第四矩展开仅留下四索引相同和两两配对项，分别为 $O(M)$、$O(M^2)$，除以 $M^4$ 得结论；第三绝对矩由 Cauchy--Schwarz 控制。$\log$ 的三阶导数在统一正区间有界，因此全区间 Taylor 余项满足 $|R_3|\le C|\epsilon|^3$。取期望得到
\begin{equation}
\E\log\widehat Z_M-\log Z=-\frac{V}{2MZ^2}+O(M^{-3/2}).
\end{equation}
方差中线性项贡献 $V/(MZ^2)$，其与二次项的交叉由第三矩控制，二次项平方由第四矩控制，故 $\Var\log\widehat Z_M=V/(MZ^2)+O(M^{-3/2})$。
\end{proof}

\begin{theorem}[累积量截断的可控偏差]
固定有限正分布，令 $\varphi(t)=\log\E_qe^{tr}$。用 $M\ge L$ 个独立样本无偏估计前 $L$ 阶累积量，再构造 $\widehat L=\sum_{j=1}^L\eta^j\widehat\kappa_j/j!$，则
\begin{equation}
\E\widehat L-\log Z=-\frac{\eta^{L+1}}{L!}\int_0^1(1-u)^L\varphi^{(L+1)}(u\eta)\,du.
\end{equation}
固定阶数下方差为 $O(M^{-1})$。有限阶公式不要求无限 Taylor 级数收敛。
\end{theorem}
\begin{proof}
几何信号 $f=\eta r$ 时可改用累积量展开。令 $\varphi(t)=\log\E_qe^{tr}$，指数倾斜分布为 $q_t\propto qe^{tr}$。逐项微分得 $\varphi^{(j)}(t)$ 为 $q_t$ 下第 $j$ 阶累积量。有限阶 Taylor 公式给出
\begin{equation}
\varphi(\eta)=\sum_{j=1}^L\frac{\eta^j}{j!}\varphi^{(j)}(0)
+\frac{\eta^{L+1}}{L!}\int_0^1(1-u)^L\varphi^{(L+1)}(u\eta)\,du.
\end{equation}
每阶累积量可按集合分拆 $\pi$ 写成
\begin{equation}
\kappa_j(r)=\sum_{\pi\in\mathcal P_j}(|\pi|-1)!(-1)^{|\pi|-1}
\prod_{B\in\pi}\E r^{|B|}.
\end{equation}
用互异样本索引的乘积平均替代每个矩乘积，即得到无偏估计，要求样本数 $M\ge j$。独立性逐项保证均值正确。固定阶有界核的两项乘积只在索引交叠时可能有协方差，交叠概率为 $O(1/M)$，故方差及各阶协方差均为 $O(1/M)$。截断估计的偏差正是上式积分余项的负值；无限级数还需位于收敛半径内。
\end{proof}

\begin{proposition}[头部配分函数的尾部级数估计]
若确定头部 $Z_H>0$、尾部贡献 $\tau=\E_qu$ 满足 $\tau/Z_H<1$，则以互异样本乘积无偏估计 $\tau^j$ 所得 $L$ 阶 log 尾部展开，仅保留确定截断偏差，其绝对值不超过 $(\tau/Z_H)^{L+1}/(L+1)$。样本数需 $M\ge L$。
\end{proposition}
\begin{proof}
另一种展开从确定头部 $Z_H>0$ 开始。令 $\tau=\E_qu$ 为尾部贡献，若 $\tau/Z_H<1$，则
\begin{equation}
\log(Z_H+\tau)=\log Z_H+\sum_{j=1}^L\frac{(-1)^{j+1}}j(\tau/Z_H)^j+R_{L+1},
\quad |R_{L+1}|\le\frac{(\tau/Z_H)^{L+1}}{L+1}.
\end{equation}
用互异索引乘积的 U-统计量无偏估计 $\tau^j$，样本数 $M\ge L$，留下可控的确定截断偏差。尾部均值小不自动提供逐样本小上界，实际方差仍要由核的矩计算。
\end{proof}

\section{截断、归一化与尾部匹配的差别}
\begin{proposition}[截断与 Tail-aware 的平方误差比较]
\label{prop:truncation-mse}
固定正分布 $P,Q$ 和非空真子集 $S$，令
$H=\sum_{a\in S}P_a\log(P_a/Q_a)$，
$\tau_P=P(S^c)$、$\tau_Q=Q(S^c)$，
$M=\tau_P\log(\tau_P/\tau_Q)$、
$B=\tau_P\KL(P^{S^c}\Vert Q^{S^c})\ge0$。
则完整 KL 为 $K=H+M+B$，截断值 $H$ 与 Tail-aware 值 $H+M$ 满足
\begin{equation}
(H-K)^2-(H+M-K)^2=M(M+2B).
\end{equation}
因此 Tail-aware 的下界性质不保证其平方误差小于截断值。
\end{proposition}
\begin{proof}
在尾部逐点分解
$\log(P_a/Q_a)=\log(\tau_P/\tau_Q)+
\log(P_a^{S^c}/Q_a^{S^c})$，求和得 $K=H+M+B$。
两个误差分别为 $-(M+B)$ 和 $-B$，平方相减即得公式。
$M$ 可正可负。例如取 $P=(0.4,0.2,0.4)$、
$Q=(0.2,0.4,0.4)$、$S=\{1\}$，
有 $M=0.6\log(3/4)<0$、
$B=0.2\log(2/3)+0.4\log(4/3)>0$，
且 $M+2B<0$，此时差值为正。
再保持头尾质量不变并令一个 $Q$ 尾概率趋零，$B$ 无界增长，
最终 $M+2B>0$，差值转负。两种排序均可发生。
\end{proof}

\begin{proposition}[归一化头部 KL 对尾 logit 的不敏感性]
\label{prop:tail-logit}
令 $P=\operatorname{softmax}(z)$，$Q,S$ 固定，
$P_S=\sum_{a\in S}P_a$。
归一化头部 KL $K_{\rm norm}=\KL(P^S\Vert Q^S)$ 与截断值 $H$ 满足
\begin{equation}
\frac{\partial K_{\rm norm}}{\partial z_j}=0,\qquad
\frac{\partial H}{\partial z_j}=-P_j(H+P_S),\qquad j\notin S.
\end{equation}
\end{proposition}
\begin{proof}
$P_a^S=e^{z_a}/\sum_{i\in S}e^{z_i}$ 不含尾 logit，故第一式成立。
对 $a\in S,j\notin S$，
$\partial P_a/\partial z_j=-P_aP_j$、
$\partial\log(P_a/Q_a)/\partial z_j=-P_j$。
对 $H$ 使用乘积法则并求和得第二式。
归一化会消去头部总质量，而截断求和仍通过 softmax 分母依赖尾部。
\end{proof}

\begin{proposition}[尾形状匹配时的残差梯度噪声]
\label{prop:tailmatch-noise}
设 $P^{S^c}=Q^{S^c}$，$\ell=\log(\tau_P/\tau_Q)$，
$\psi_a=\nabla\log P_a$、
$\bar\psi_C=\E_{P^{S^c}}\psi$。
前文标量桶残差修正（取 $c=\ell$）与 SC-TA 的尾形状残差在此点逐样本为零。
ResOPD 的零均值残差为
$J_A\ell(\psi_A-\bar\psi_C)$，$A\sim P$，
其协方差为
\begin{equation}
\tau_P\ell^2\Cov_{P^{S^c}}(\psi).
\end{equation}
该量可严格为正，即使 Tail-aware 的值和梯度已精确。
\end{proposition}
\begin{proof}
尾形状匹配给尾 log-ratio 恒为 $\ell$，所以
标量桶残差修正的值残差及 SC-TA 的形状残差因子均为零。
ResOPD 保留 $\ell$ 乘中心化尾 score；其尾条件均值为零。
对二阶矩按 $J_A=\mathbf1\{A\notin S\}$ 条件化，即得协方差。
例取 $P=(0.4,0.2,0.4)$、$Q=(0.7,0.1,0.2)$、
$S=\{1\}$。两者尾条件分布均为 $(1/3,2/3)$，
$\tau_P=3/5$、$\ell=\log2$。分类 softmax 的
$\psi_a=e_a-P$ 给尾协方差迹 $1-(1/3)^2-(2/3)^2=4/9$，
故残差协方差迹为 $4(\log2)^2/15>0$。
\end{proof}

\section{异分布采样与自动微分语义}
\begin{proposition}[异分布下的双控制项一般耦合]
\label{prop:offpolicy-control}
令 $A\sim\beta$，$\beta$ 为正分布，$w=P/\beta$，
$J_a=\mathbf1\{a\notin S\}$。取尾向量 $\chi_a$ 满足
$\sum_{a\notin S}P_a\chi_a=0$，定义零均值控制项
$U_{0,j}=(J_Aw_A-\tau_P)e_j$、$U_1=J_Aw_A\chi_A$。
则其 Gram 交叉项为
\begin{equation}
\E_\beta\langle U_{0,j},U_1\rangle
=\sum_{a\notin S}\frac{P_a^2}{\beta_a}\langle e_j,\chi_a\rangle.
\end{equation}
在 $\beta=P$ 时该项为零，一般 $\beta$ 下不为零。
\end{proposition}
\begin{proof}
$\E_\beta J_Aw_A=\tau_P$、
$\E_\beta J_Aw_A\chi_A=0$ 保证两个控制项均值为零。
展开内积，常数 $-\tau_P$ 对应项因后一等式消去，
其余项中的 $\beta_aw_a^2=P_a^2/\beta_a$ 给所述结果。
例如在两点全尾集合上，取 $P=(1/2,1/2)$、
$\beta=(1/4,3/4)$ 和一维 $\chi=(1,-1)$，
交叉项为 $1-1/3=2/3$。因此两个系数的联合最优值
应解完整 Gram 线性系统，不能分别按正交控制项优化。
\end{proof}

\begin{proposition}[停止梯度的双参数表示]
\label{prop:stopgrad-diagonal}
设样本 $a$ 固定，$S(\theta,\bar\theta,a)$ 对两参数可微。
冻结第二参数后在对角线取值的反传场定义为
\begin{equation}
B(\theta,a)=\partial_1S(\theta,\bar\theta,a)|_{\bar\theta=\theta}.
\end{equation}
而前向数值函数 $s(\theta,a)=S(\theta,\theta,a)$ 的普通导数为
\begin{equation}
\nabla s(\theta,a)=B(\theta,a)
+\partial_2S(\theta,\bar\theta,a)|_{\bar\theta=\theta}.
\end{equation}
若动作分布 $P_\theta$ 也随参数变化，则
\begin{equation}
\nabla\E_{P_\theta}s
=\E_{P_\theta}[\nabla s+s\nabla\log P_\theta].
\end{equation}
\end{proposition}
\begin{proof}
第一式是复合映射 $\theta\mapsto(\theta,\theta)$ 的链式法则。
第二式在有限词表上直接对 $\sum_aP_\theta(a)s(\theta,a)$ 使用乘积法则。
这两个额外项分别来自被停止的参数路径和采样测度。
因此前向数值相同、固定样本反传相同、总体梯度相同是三个需要分别核对的命题。
即使 $B$ 连续，也需满足相应可积条件才能将它称为某个标量目标的梯度。
\end{proof}

\begin{proposition}[前缀混合下的误差与方差口径]
\label{prop:prefix-mse}
设前缀 $s\sim d$，估计器 $\widehat K$ 的条件偏差为
$c(s)=\E[\widehat K\mid s]-K(s)$，条件方差为 $V(s)$。
只要二阶矩有限，
\begin{align}
\E[\widehat K-K(s)]&=\E_dc(s),\\
\E[(\widehat K-K(s))^2]&=\E_d[V(s)+c(s)^2],\\
\Var(\widehat K)&=\E_dV(s)+\Var_d(K(s)+c(s)).
\end{align}
\end{proposition}
\begin{proof}
条件化后误差均值为 $c(s)$，其平方均值为 $V(s)+c(s)^2$，
对 $d$ 平均得前两式。对估计器本身使用全方差公式得第三式。
跨前缀原始数值方差包含真实目标 $K(s)$ 的变化，
不能用它替代估计误差的均方值或条件采样方差。
\end{proof}


\chapter{条件中心与有限样本拟合}
\label{app:centers}
\section{中心估计代入损失后的对象}
固定前缀、锚点和候选参数，记 $h_a=U_a-f_a$、$R=C_qh$、$\psi_a=\nabla\log p_a$。BPO 梯度为 $\eta^{-1}\Cov_q(h,\psi)$。这里中心化是损失定义的一部分；用于 score 估计的额外 baseline 是另一个选择。

对 Mean，设 $\widehat c_f$ 与主动作 $A\sim q$ 独立，且其分布不依赖学生参数，$\E\widehat c_f=\E_qf$。平方展开给出
\begin{equation}
\E\left[\tfrac12(U_A-f_A+\widehat c_f)^2\right]
=L_{\rm Mean}+\tfrac12\Var(\widehat c_f).
\end{equation}
独立零均值误差使交叉项消失，最后一项与学生无关，因此梯度无偏。只知道数值估计无偏而没有独立性时，不能做此抵消。

对 BPO，令主动作 $A$ 和 $M$ 个辅助动作 $B_j$ 独立同分布于 $q$。记 $\bar h_B=M^{-1}\sum_jh_{B_j}$。独立性给出
\begin{equation}
\E(h_A-\bar h_B)^2=(1+M^{-1})\Var_qh.
\end{equation}
因此 $\frac{M}{M+1}(h_A-\bar h_B)^2/(2\eta)$ 对损失和直接反传梯度均无偏。与 Mean 不同，辅助中心含有学生参数，额外方差项的导数不能丢掉。同前缀共 $M\ge2$ 个动作的留一中心对应倍率 $(M-1)/M$。

\section{无中心样本方差与 U-统计量}
两个独立动作给出
\begin{equation}
L_{\rm BPO}=\frac1{4\eta}\E_{a,b\sim q}(h_a-h_b)^2,\qquad
\bm G=\nabla L_{\rm BPO}
=\E_{a,b}\frac{(h_a-h_b)(\psi_a-\psi_b)}{2\eta}.
\label{eq:pairwise}
\end{equation}
逐项展开利用 $\E h_a h_b=(\E h)^2$ 即证。若 $p=\operatorname{softmax}(z)$，$h_a-h_b=z_a-z_b-\log(q_a/q_b)-(f_a-f_b)$，学生配分函数与两个条件中心全部消失。

对 $M\ge2$ 个独立动作，定义
\begin{equation}
\widehat L_M=\frac{\sum_i(h_i-\bar h)^2}{2\eta(M-1)},\qquad
\nabla\widehat L_M=\frac{\sum_i(h_i-\bar h)\psi_i}{\eta(M-1)}.
\end{equation}
恒等式 $\sum_{i<j}(h_i-h_j)^2=M\sum_i(h_i-\bar h)^2$ 使其等于成对核的 U-统计量平均。每一对的均值正确，故整体无偏；各对相互相关，不能把方差简单除以对数。

\begin{theorem}[消去中心的精确方差代价]
令成对梯度核 $H(a,b)=(h_a-h_b)(\psi_a-\psi_b)/(2\eta)$ 的迹方差为 $V_2$。令精确双中心估计器
\begin{equation}
g_c(a)=\frac{(h_a-\E_qh)(\psi_a-\E_q\psi)}{\eta},
\qquad V_c=\Var_qg_c.
\end{equation}
则
\begin{equation}
\Var(\nabla\widehat L_M)
=\frac{(M-2)V_c+2V_2}{M(M-1)},\qquad
\Var(\nabla\widehat L_M)-\frac{V_c}{M}
=\frac{2V_2-V_c}{M(M-1)}\ge0.
\end{equation}
\end{theorem}
\begin{proof}
对 $b$ 条件积分，展开乘积得到 $\E_bH(a,b)=[g_c(a)+\bm G]/2$。令
\begin{equation}
H_1(a)=\E_bH(a,b)-\bm G=(g_c(a)-\bm G)/2,\quad
H_2(a,b)=H(a,b)-\bm G-H_1(a)-H_1(b).
\end{equation}
有 $\E H_1=0$、$\E[H_2(a,b)\mid a]=\E[H_2(a,b)\mid b]=0$。因此一阶项与二阶余项正交，两不同配对的二阶余项即使共用一个索引，其协方差也由条件零均值消失。单对核给出 $V_2=2\Var H_1+\Var H_2=V_c/2+\Var H_2$。

样本平均核的中心化表达为
\begin{equation}
\frac2M\sum_iH_1(a_i)+\binom M2^{-1}\sum_{i<j}H_2(a_i,a_j).
\end{equation}
第一项方差为 $4\Var H_1/M$，第二项为 $\Var H_2/\binom M2$，交叉项为零。代入 $\Var H_1=V_c/4$ 与 $\Var H_2=V_2-V_c/2$ 并整理，即得公式及非负性。
\end{proof}

\begin{proposition}[排除重复动作的半正定改进]
先抽 $a\sim q$，再抽 $b\sim q(\cdot\mid b\ne a)$，使用加权核 $H_{\rm nr}=(1-q_a)H(a,b)$。则均值仍为 $\bm G$，且
\begin{equation}
\Cov_{q\otimes q}H-\Cov_{\rm nr}H_{\rm nr}
=\E_{q\otimes q}[q_aH(a,b)H(a,b)^\top]\succeq0.
\end{equation}
\end{proposition}
\begin{proof}
有序对概率为 $q_aq_b/(1-q_a)$，且 $H(a,a)=0$。乘以一个 $(1-q_a)$ 后求和恢复独立抽样均值；外积含两个权重，故二阶矩变为 $\E_{q\otimes q}[(1-q_a)HH^\top]$。两个均值相同，协方差之差等于二阶矩之差，每项均为非负权重的外积。正支持且至少二动作保证 $1-q_a>0$。
\end{proof}

若两个独立动作来自另一个提案 $\beta$，成对损失需乘 $(q_a/\beta_a)(q_b/\beta_b)$。先对单动作 residual 乘权重再平方，会产生错误的权重平方。

\section{教师中心的有符号与非负估计}
记 $\rho=T/q$、$A=A_{\kappa,\eta}$、$\phi(\rho)=\rho-1-A(\rho)$。两个无偏估计为
\begin{equation}
\widehat K_-=-M^{-1}\sum_iA(\rho_i),\qquad
\widehat K_+=M^{-1}\sum_i\phi(\rho_i),\qquad a_i\overset{\rm iid}\sim q.
\end{equation}
$\phi$ 凸且在一处值与导数为零，所以第二个逐样本非负。第一种方差为 $V_{\kappa,\eta}/M$。非负不保证更小方差：取 $q=(\epsilon,1-\epsilon)$、$T=(1/2,1/2)$，固定正 $\kappa,\eta<1$。稀有项的 $A=O(\log(1/\epsilon))$，有符号估计方差趋零；$\phi\sim1/(2\epsilon)$，非负估计方差按 $1/(4\epsilon)$ 发散。

局部情况不同。若 $\rho=1+\epsilon u$、$\E_qu=0$ 且 $u$ 有界，令 $c_0=1-\kappa(1-\eta)$，则
\begin{equation}
-A=-\epsilon u+\tfrac12c_0\epsilon^2u^2+O(\epsilon^3),\qquad
\phi=\tfrac12c_0\epsilon^2u^2+O(\epsilon^3).
\end{equation}
两种方差分别为 $\epsilon^2\E_qu^2+O(\epsilon^3)$ 与 $\frac14c_0^2\epsilon^4\Var_qu^2+O(\epsilon^5)$。这需要比率一致接近一，不能只用 KL 小替代。

同动作估计中心尤其需要核对：若用该动作的 $\phi$ 补到 $A$ 上，逐样本有 $A+\phi=\rho-1$，形状变换被完全抵消。若用同一组 $M$ 个样本减去经验均值，
\begin{equation}
\E\left[M^{-1}\sum_i(A_i-\bar A)\psi_i\right]
=\frac{M-1}{M}\Cov_q(A,\psi).
\end{equation}
展开双重和，同索引项 $M$ 个、异索引项 $M(M-1)$ 个，即得到该系数；$M=1$ 时更新为零。

固定前缀的精确中心化不改变标量条件方差。跨冻结前缀测度 $d$，记 $V_s=\Var_{q(\cdot\mid s)}A_s$，则全方差公式给出 $\Var_{d,q}A=\E_dV_s+\Var_dK_T(s)$，而 $\Var_{d,q}(A+K_T)=\E_dV_s$。中心化移除不同前缀的均值差异；这不直接推出乘以 score 后的参数梯度方差排序。

\section{有界包络与头尾区间}
当 $\kappa>0,\eta<1$，$c=-\log(1-\eta)/(\kappa\eta)>1$，且 $0\le\phi(x)\le x+c-1$。定义
\begin{equation}
\mu_c=(1-1/c)q+T/c,\qquad Z(a)=q_a\phi(\rho_a)/\mu_c(a).
\end{equation}
则 $\E_{\mu_c}Z=K_T$、$0\le Z\le c$。由于 $Z^2\le cZ$，
\begin{equation}
\Var_{\mu_c}Z\le K_T(c-K_T)\le c^2/4.
\end{equation}
独立平均 $M$ 次后方差除以 $M$，Hoeffding 给 $\Pr(|\widehat K-K_T|\ge\epsilon)\le2e^{-2M\epsilon^2/c^2}$。整个分布类上 $K_T$ 可逼近 $c$，所以任何无偏估计器的统一样本上界不能低于 $c$。

一般混合 $\mu_\zeta=(1-\zeta)q+\zeta T$ 的逐样本上确界为
\begin{equation}
C_\zeta=\max\{(c-1)/(1-\zeta),1/\zeta\}.
\end{equation}
两个值分别由 $\rho\to0,\infty$ 逼近。使统一范围最小的唯一权重是 $\zeta=1/c$。固定输入的方差则为 $\E_q[\phi(\rho)^2/(1+\zeta(\rho-1))]-K_T^2$，二阶导数为
\begin{equation}
2\E_q\frac{\phi(\rho)^2(\rho-1)^2}{[1+\zeta(\rho-1)]^3}\ge0.
\end{equation}
所以输入方差最优权重可经一维凸优化得到，但一般不同于最小统一范围权重。

保留头部 $S$，尾部质量为 $\delta=q(S^c)$、$\tau=T(S^c)$。令 $K_S^T=-\sum_Sq_aA(\rho_a)$。尾部条件均值 $\E_{q^{\rm tail}}\rho=\tau/\delta$，由 $-A$ 凸以及 $-A\le c$，
\begin{equation}
K_S^T-\delta A(\tau/\delta)\le K_T\le K_S^T+\delta c.
\label{eq:center-interval}
\end{equation}
扩大嵌套头部时，下界由 Jensen 非减，上界每加入动作减少 $q_a[c+A(\rho_a)]\ge0$，故非增。尾部为空时直接取精确头部值，不计算零质量比。

还能构造始终落在区间内的无偏估计。令 $r_0=\tau/\delta$、$H=\sum_Sq\phi$，
\begin{equation}
L=H+\delta\phi(r_0),\quad U=H+\tau+(c-1)\delta,\quad w=U-L.
\end{equation}
取凸余项 $g(x)=\phi(x)-\phi(r_0)-\phi'(r_0)(x-r_0)$，令 $u=A'(r_0)>0$、$v=c+A(r_0)-r_0u\ge0$。$v\ge0$ 来自凹函数 $A$ 的切线在零处不低于 $A(0)=-c$。于是 $0\le g(x)\le ux+v$，且
\begin{equation}
\sum_{S^c}q_ag(\rho_a)=K_T-L,\quad u\tau+v\delta=w.
\end{equation}
若 $w>0$，在尾部取 $\nu_a=(uT_a+vq_a)/w$、$Y=q_ag(\rho_a)/\nu_a$，则 $0\le Y\le w$。因此
\begin{equation}
\widehat K=L+M^{-1}\sum_iY_i,\quad
\E\widehat K=K_T,\quad L\le\widehat K\le U,\quad
\Var\widehat K\le(K_T-L)(U-K_T)/M.
\end{equation}
证明分别由期望相消和 $Y^2\le wY$ 得到。$w=0$ 时直接返回 $L$。该方案要求能从条件尾部分布采样并查询抽中动作概率；只有头部和尾部总质量时不能恢复未知尾部形状。

\section{中心误差、baseline 与更新场}
设学生、教师中心误差分别为 $e_S=\widehat K_S-K_S$、$e_T=\widehat K_T-K_T$。原始 score 形式的梯度误差为
\begin{equation}
\widehat g-g=(e_S/\eta-e_T)\E_q\psi.
\label{eq:center-gradient-bias}
\end{equation}
在 $p=q$ 时该均值为零，一般内层不为零。若中心与主动作相关，需保留 $\E[e(A)\psi_A]$，不能把误差提出期望。

若学生中心精确、教师误差为冻结常数 $e_T$，直接微分平方损失时，残差变为 $R-\eta e_T$。由于 $\E_qR=0$，损失仅增加 $\eta e_T^2/2$，梯度不变；对应的 score 应为 $(R/\eta-e_T)(\psi-\E_q\psi)$。在目标点 $p=m$，其采样噪声为
\begin{equation}
e_T^2\E_q\|\psi-\E_q\psi\|^2.
\end{equation}
完整求和精确抵消该噪声，单动作抽样则保留。用学生反馈 $y=C_qU/\eta-A$ 估计同一个教师中心时，$\E_qy=K_T$ 且 $\Var_qy=2L_{\rm BPO}/\eta$，噪声可随拟合消失。

若 $e_S$ 随参数变化，式~\eqref{eq:center-gradient-bias} 的场含 $-(e_S/\eta-e_T)\nabla K_S$。其反对称 Jacobian 一般不为零，因为 $\nabla e_S$ 与 $\nabla K_S$ 未必平行。直接微分含近似中心的平方损失则多出 $\E_q\widehat R\,\nabla\widehat K_S/\eta$，得到另一个保守场。有限中心近似后，两种反传方式需要分别定义。

在锚点 $p=q$，action-independent baseline $b$ 不改变 $\E_q[(A-b)\psi]$ 的均值。最小迹方差的选择为
\begin{equation}
b_\star=\frac{\E_q[A\|\psi\|^2]}{\E_q\|\psi\|^2},\quad
V_g(b)-V_g(b_\star)=\E_q\|\psi\|^2(b-b_\star)^2.
\end{equation}
这是对加权二次式配方的结果，不一般等于 $\E_qA$。重要性采样估计器改用权重 $(p/q)^2\|\psi_p\|^2$。对普通 BPO 的 $q$ 测度，离开锚点后随意改 baseline 会改变均值，除非保留补偿或使用精确中心化 score。

独立随机 baseline 额外增加其方差乘以相应 score 二阶矩。同样本 baseline 即使数值无偏也可能有偏：取 baseline 等于该动作的信号，整个更新逐样本为零。跨前缀 token 不能替代同前缀独立辅助动作。

\chapter{外层动力学、扰动与序列边界}
\label{app:dynamics}
\section{幂均值的 escort 轨迹}
\begin{proposition}[escort 坐标中的精确直线轨迹]
设 $0<\kappa\le1$，$E_\kappa(q)_a=q_a^\kappa/S_\kappa(q)$、$S_\kappa(q)=\sum_aq_a^\kappa$。幂均值更新满足
\begin{equation}
E_\kappa(m)=(1-b)E_\kappa(q)+bE_\kappa(T),\quad
b=\frac{\eta S_\kappa(T)}{(1-\eta)S_\kappa(q)+\eta S_\kappa(T)}.
\end{equation}
固定教师的全部精确外层迭代位于初始 escort 分布与教师 escort 分布的同一条直线上。
\end{proposition}
\begin{proof}
对 $\kappa>0$，定义 $S_\kappa(q)=\sum_aq_a^\kappa$ 和 escort 分布 $E_\kappa(q)_a=q_a^\kappa/S_\kappa(q)$。它的逆为 $q_a\propto E_\kappa(q)_a^{1/\kappa}$。对幂均值目标取 $\kappa$ 次幂，公共归一化消去，得
\begin{equation}
E_\kappa(m)=(1-b)E_\kappa(q)+bE_\kappa(T),\quad
b=\frac{\eta S_\kappa(T)}
{(1-\eta)S_\kappa(q)+\eta S_\kappa(T)}.
\label{eq:escort}
\end{equation}
因此固定教师的精确轨迹在 escort 坐标中位于同一条直线：
\begin{equation}
E_\kappa(q_n)-E_\kappa(T)
=\prod_{j<n}(1-b_j)[E_\kappa(q_0)-E_\kappa(T)].
\end{equation}
\end{proof}

\begin{theorem}[escort 距离收缩与近似输出误差]
固定正教师、有限词表及 $0<\kappa\le1,0<\eta<1$。令
$c=(1-\eta)V^{1-\kappa}/[(1-\eta)V^{1-\kappa}+\eta S_\kappa(T)]<1$。
令 $d_\kappa(P,Q)=\|E_\kappa(P)-E_\kappa(Q)\|_1$。若 $d_\kappa(q_{n+1},m(q_n,T))\le\epsilon_n$，则
\begin{equation}
d_\kappa(q_n,T)\le c^nd_\kappa(q_0,T)+\sum_{j<n}c^{n-1-j}\epsilon_j.
\end{equation}
若 $\limsup\epsilon_n\le\epsilon_\infty$，稳态距离不超过 $\epsilon_\infty/(1-c)$；精确更新收敛到 $T$。
\end{theorem}
\begin{proof}
词表大小为 $V$，Jensen 给 $S_\kappa(q)\le V^{1-\kappa}$。令
\begin{equation}
c=\frac{(1-\eta)V^{1-\kappa}}
{(1-\eta)V^{1-\kappa}+\eta S_\kappa(T)}<1.
\end{equation}
定义距离
\begin{equation}
d_\kappa(q,T)=\|E_\kappa(q)-E_\kappa(T)\|_1.
\end{equation}
精确更新满足 $d_\kappa(q_{n+1},T)\le c\,d_\kappa(q_n,T)$。设近似输出的误差满足
\begin{equation}
d_\kappa(q_{n+1},m(q_n,T))\le\epsilon_n.
\end{equation}
三角不等式给
\begin{equation}
d_\kappa(q_n,T)\le c^nd_\kappa(q_0,T)+\sum_{j<n}c^{n-1-j}\epsilon_j.
\label{eq:escort-errors}
\end{equation}
若 $\limsup\epsilon_n\le\epsilon_\infty$，拆开有限前缀与充分晚的尾和，得到 $\limsup d_\kappa(q_n,T)\le\epsilon_\infty/(1-c)$。这是到固定教师的定向收缩，不声称两任意锚点间有相同收缩常数。
\end{proof}

\begin{proposition}[自然梯度流的 escort 表达与几何极限]
沿幂散度的自由自然梯度流，
\begin{equation}
\frac d{dt}E_\kappa(q)=\frac{S_\kappa(T)}{S_\kappa(q)}[E_\kappa(T)-E_\kappa(q)].
\end{equation}
因此正时间重参数化后是线性趋近教师的流。固定正 $q$ 时，
$\frac V\kappa(E_\kappa(q)_a-1/V)\to\log q_a-V^{-1}\sum_b\log q_b$。
\end{proposition}
\begin{proof}
连续流也有同一轨迹。令 $a_\kappa=\E_q(T/q)^\kappa$，由 $\dot q_a=q_a[(T_a/q_a)^\kappa-a_\kappa]/\kappa$ 得
\begin{equation}
\frac d{dt}q_a^\kappa=T_a^\kappa-a_\kappa q_a^\kappa,\qquad
\frac d{dt}E_\kappa(q)=\frac{S_\kappa(T)}{S_\kappa(q)}
[E_\kappa(T)-E_\kappa(q)].
\end{equation}
重新定义时间 $d\sigma/dt=S_\kappa(T)/S_\kappa(q)>0$，轨迹即为指数趋近教师的直线。$\kappa\downarrow0$ 时 escort 本身趋均匀，需要重标度：
\begin{equation}
\frac V\kappa\left(E_\kappa(q)_a-\frac1V\right)
\longrightarrow\log q_a-\frac1V\sum_b\log q_b.
\end{equation}
\end{proof}

\section{几何端点的完整稳定区间}
这一节允许 $\eta>0$，但只讨论 $\kappa=0$ 的几何更新及后面相应的 Mean 更新。正形状的概率插值仍取 $0<\eta<1$。

\begin{theorem}[几何端点的闭式轨迹与稳定分类]
固定 $T>0$，几何更新为 $q^+\propto q^{1-\eta}T^\eta$。同一教师下，Hilbert 距离精确乘以 $|1-\eta|$。变步长轨迹满足
\begin{equation}
q_k(a)\propto T_a e^{\chi_k\log(q_{0,a}/T_a)},\qquad
\chi_k=\prod_{j<k}(1-\eta_j).
\end{equation}
非教师初值收敛当且仅当 $\chi_k\to0$。常步长在 $0<\eta<2$ 收敛，$\eta=2$ 两周期，$\eta>2$ 跨度发散。
\end{theorem}
\begin{proof}
固定正教师 $T$，令 $q^+\propto q^{1-\eta}T^\eta$。归一化常数在动作差中消失，所以
\begin{equation}
d_H(q^+,T)=|1-\eta|d_H(q,T),\qquad
d_H(q^+,\widetilde q^+)=|1-\eta|d_H(q,\widetilde q).
\end{equation}
对变步长，记 $\chi_k=\prod_{j<k}(1-\eta_j)$、$u=\log(q_0/T)$，归纳得到
\begin{equation}
q_k(a)=\frac{T_ae^{\chi_ku_a}}{\sum_bT_be^{\chi_ku_b}}.
\end{equation}
归纳步骤中上一轮的配分函数仅产生公共乘数，归一化后消去。非教师初值下，收敛当且仅当 $\chi_k\to0$。常步长 $0<\eta<2$ 收敛，$\eta=1$ 一步到教师，$\eta=2$ 为两周期，$\eta>2$ 的跨度发散。
\end{proof}

\begin{proposition}[几何变步长的充要收敛条件]
若 $\eta_k\in(0,2)$ 且没有一步等于一，则非教师初值收敛到教师，当且仅当 $\sum_k\eta_k(2-\eta_k)=\infty$。若某一步等于一，此后精确保持教师。
\end{proposition}
\begin{proof}
若所有 $\eta_k\in(0,2)$ 且没有一步等于一，令 $a_k=|1-\eta_k|\in(0,1)$。乘积趋零当且仅当 $\sum_k(1-a_k)=\infty$：一边由 $-\log a\ge1-a$，另一边在 $a$ 接近一时由 $-\log a\le2(1-a)$ 得到。再用 $1-a_k\le1-a_k^2\le2(1-a_k)$，得到等价判据 $\sum_k\eta_k(2-\eta_k)=\infty$。
\end{proof}

\begin{proposition}[几何路径的双向 KL 积分]
令 $u=\log(q_0/T)$、$\Psi(\chi)=\log\E_Te^{\chi u}$、$q^{[\chi]}\propto Te^{\chi u}$。则
\begin{align}
\KL(q^{[\chi]}\Vert T)&=\int_0^\chi s\Psi''(s)\,ds,\\
\KL(T\Vert q^{[\chi]})&=\int_0^\chi(\chi-s)\Psi''(s)\,ds,\qquad
\Psi''(\chi)=\Var_{q^{[\chi]}}u.
\end{align}
非教师初值、$0<\eta<2$ 且 $\eta\ne1$ 时，几何迭代的局部 KL 比例为 $(1-\eta)^2$；全局有 $\KL(q_k\Vert T)\le\chi_k^2(\osc u)^2/8$。
\end{proposition}
\begin{proof}
定义 $\Psi(\chi)=\log\sum_aT_ae^{\chi u_a}$。有限求和微分给 $\Psi'=\E_{q^{[\chi]}}u$、$\Psi''=\Var_{q^{[\chi]}}u$。于是
\begin{align}
\KL(q^{[\chi]}\Vert T)&=\chi\Psi'(\chi)-\Psi(\chi)
=\int_0^\chi s\Psi''(s)\,ds,\\
\KL(T\Vert q^{[\chi]})&=\Psi(\chi)-\chi\Psi'(0)
=\int_0^\chi(\chi-s)\Psi''(s)\,ds.
\end{align}
正、负半轴上两种 KL 各自随离零距离增加而增加，但跨越两侧的值不能仅由 $|\chi|$ 比较。教师附近有 KL 比例极限 $(1-\eta)^2$；全局只有绝对界 $\KL(q_k\Vert T)\le\chi_k^2(\osc u)^2/8$，由 $\Psi''\le(\osc u)^2/4$ 积分得到。
\end{proof}

\begin{proposition}[几何更新的三点恒等式]
几何更新对任意 $\eta>0$ 满足
\begin{equation}
\eta\KL(q\Vert T)=\eta\KL(q^+\Vert T)+(1-\eta)\KL(q^+\Vert q)+\KL(q\Vert q^+).
\end{equation}
若 $0<\eta\le1$，教师逆向 KL 单步非增，且常步长下 $\sum_k\KL(q_k\Vert q_{k+1})\le\eta\KL(q_0\Vert T)$。稳定区间内的 $\eta>1$ 不一般保证 KL 单步下降。
\end{proposition}
\begin{proof}
由 $q_a^+=q_a^{1-\eta}T_a^\eta/Z$，
\begin{align}
\log(q_a^+/q_a)&=-\eta\log(q_a/T_a)-\log Z,\\
\log(q_a^+/T_a)&=(1-\eta)\log(q_a/T_a)-\log Z.
\end{align}
第一式在 $q$ 下取负期望，得到 $\KL(q\Vert q^+)=\eta\KL(q\Vert T)+\log Z$。两式在 $q^+$ 下取期望并按 $\eta,1-\eta$ 组合，得到 $\eta\KL(q^+\Vert T)+(1-\eta)\KL(q^+\Vert q)=-\log Z$。相加消去配分常数，即得
\begin{equation}
\eta\KL(q\Vert T)=\eta\KL(q^+\Vert T)
+(1-\eta)\KL(q^+\Vert q)+\KL(q\Vert q^+).
\end{equation}
当 $0<\eta\le1$ 时，这给单步 KL 下降及 $\sum_k\KL(q_k\Vert q_{k+1})\le\eta\KL(q_0\Vert T)$。$\eta>1$ 时不能把跨度稳定性换成每步 KL 下降。例如 $q=(10^{-12},1-10^{-12})$、$T=(0.1,0.9)$、$\eta=1.1$，逆向 KL 从约 $0.10536$ 增至 $0.71385$。
\end{proof}

\section{Mean 的单轮界与精确外层}
本节取 $r=\log(T/q)$、$d=\KL(q\Vert T)$、$G=\eta(r+d)$。Mean 的 KKT 为
\begin{equation}
\log w_a+\lambda w_a=\eta(r_a+d),\quad
w_a=p_a^{\rm M}/q_a,\quad \E_qw=1,\quad
\lambda=\KL(q\Vert p^{\rm M}).
\end{equation}
\begin{proposition}[Mean 归一化乘子的上下界]
设 $G=\eta C_q\log(T/q)$，精确 Mean 乘子为 $\lambda$。令 $A_1=\E_qe^G,A_2=\E_qe^{2G}$，则
\begin{equation}
\frac{A_1(A_1-1)}{A_2}\le\lambda\le\log A_1.
\end{equation}
若 $q=T$，两界均为零；否则 $\lambda>0$。
\end{proposition}
\begin{proof}
令 $X=e^G$、$A_1=\E_qX$、$A_2=\E_qX^2$。由 $\lambda=\E_qW_0(\lambda X)$，$W_0$ 的凹性给 $\lambda\le\log A_1$。其凹性可由 $W_0'(x)=e^{-W_0(x)}/(1+W_0(x))$ 再求导看出。

下界使用 $W_0(z)\ge z/(1+z)$：写 $z=we^w$，等价于 $(1-w)e^w\le1$；该函数在 $w\ge0$ 非增。Cauchy--Schwarz 再给
\begin{equation}
\lambda\ge\E_q\frac{\lambda X}{1+\lambda X}
\ge\frac{\lambda A_1^2}{A_1+\lambda A_2},
\qquad
\frac{A_1(A_1-1)}{A_2}\le\lambda\le\log A_1.
\end{equation}
非恒定信号时 $\lambda>0$ 可以约去；恒定信号两边为零。代入 $A_1=e^{\eta d}Z_\eta,A_2=e^{2\eta d}Z_{2\eta}$，得到原稿的 Rényi 归一化界，其中 $Z_\nu=\sum q^{1-\nu}T^\nu$。
\end{proof}

\begin{proposition}[精确 Mean 的教师 KL 下降]
固定 $q,T>0$，令 $d=\KL(q\Vert T)$，$p^{\rm M}$ 是信号 $f=\eta\log(T/q)$ 的精确 Mean 解。则
\begin{equation}
\KL(p^{\rm M}\Vert T)\le d-(\eta^{-1}-1)\KL(p^{\rm M}\Vert q)
-\frac{\lambda}{2\eta}\chi^2(p^{\rm M}\Vert q).
\end{equation}
当 $0<\eta\le1$ 且 $q\ne T$ 时，教师 KL 严格下降。
\end{proposition}
\begin{proof}
Mean 最大化式~\eqref{eq:mean-mixed}，将它与可行点 $q$ 比较，再除以 $\eta$，得到
\begin{equation}
\KL(p^{\rm M}\Vert T)
\le d-(\eta^{-1}-1)\KL(p^{\rm M}\Vert q)
-\frac{\lambda}{2\eta}\chi^2(p^{\rm M}\Vert q).
\end{equation}
当 $0<\eta\le1,q\ne T$，最后一项严格为正，所以教师 KL 严格下降。
\end{proof}

\begin{theorem}[Mean 的跨度稳定性]
固定正教师和正初值，逐轮精确 Mean 求解并同步。$0<\eta\le2$ 时跨度严格下降到零；$0<\eta<2$ 时
\begin{equation}
b_{k+1}\le C_\eta(b_k)b_k,\quad
C_\eta(b)=\max\{1-\eta/(1+\eta b),\eta-1\}<1
\end{equation}
对 $b>0$ 成立。$\eta>2$ 时教师局部不稳定。
\end{theorem}
\begin{proof}
在一轮中固定 $\lambda,d$，把 $w$ 看成 $r$ 的函数。隐式求导给 $d\log w/dr=\eta/(1+\lambda w)$。因此相对教师的新 log-ratio $v(r)=\log w(r)-r$ 满足
\begin{equation}
v'(r)=-1+\eta/(1+\lambda w(r)).
\end{equation}
非教师状态 $\lambda>0$，当 $\eta\le2$ 时 $-1<v'<1$。在有限 reward 区间上连续，故均值定理给严格跨度收缩。为得到显式系数，归一化保证 $w_{\max}\ge1$；最大 reward 处 KKT 给
\begin{equation}
\lambda w_{\max}=\eta(r_{\max}+d)-\log w_{\max}
\le\eta(r_{\max}-r_{\min})=\eta b.
\end{equation}
所以导数落在 $[-1+\eta/(1+\eta b),\eta-1)$，绝对值由 $C_\eta(b)$ 控制。该系数随 $b$ 非减，故 $\eta<2$ 时由初始跨度给出统一几何界。

在 $\eta=2$，显式系数只有一。跨度有界蕴含 $T_ae^{-b_0}\le q_a\le T_ae^{b_0}$，迭代位于正单纯形内的紧集。归一化根严格单调且连续，教师附近 $\lambda\le\log\E_qe^G\to0$，故更新也连续。若跨度极限 $b_\infty>0$，该紧集上远离教师的更新跨度差处处正且连续，有统一正下界，与跨度收敛矛盾。故极限为零。最后，下节线性化为 $(1-\eta)y$，$\eta>2$ 时其模大于一，局部不稳定。
\end{proof}

\section{Mean 的三次项与边界率}
\begin{theorem}[Mean 的三阶局部展开]
固定正教师和 $\eta>0$。令 $y=C_T\log(q/T)$、$V=\E_Ty^2$，则当 $\|y\|_\infty\to0$，
\begin{equation}
y^+=(1-\eta)y+\tfrac12\eta^3Vy+O(\|y\|_\infty^4).
\end{equation}
$\eta=1$ 时，$V^+=V^3/4+O(V^{7/2})$，且非平凡趋近满足
$\KL(p^{\rm M}\Vert T)/\KL(q\Vert T)^3\to1$。
\end{theorem}
\begin{proof}
定义教师中心坐标 $y=C_T\log(q/T)$，$q=Te^y/\E_Te^y$，$V=\E_Ty^2$、$\delta=\|y\|_\infty$。固定正教师使 $T_{\min}\delta^2\le V\le\delta^2$。Taylor 展开给
\begin{equation}
\log\E_Te^y=V/2+O(\delta^3),\quad
\E_qy=V+O(\delta^3),\quad d=V/2+O(\delta^3).
\end{equation}
此时 $G=\eta(-y+\E_qy)$。归一化方程
\begin{equation}
\Phi(\lambda,y)=\E_q\exp[G-W_0(\lambda e^G)]-1=0
\end{equation}
在原点对 $\lambda$ 的导数为 $-1$，隐函数定理给解析根。展开 $W_0(z)=z+O(z^2)$，
\begin{equation}
0=\E_qe^G-1-\lambda\E_qe^{2G}+O(\lambda^2),\qquad
\E_qe^G-1=\eta^2V/2+O(\delta^3).
\end{equation}
因 $\E_qe^{2G}$ 在邻域内远离零，先得 $\lambda=O(\delta^2)$，再得 $\lambda=\eta^2V/2+O(\delta^3)$。

将 KKT 输出按教师中心化：
\begin{align}
y^+&=(1-\eta)y-C_TW_0(\lambda e^G),\\
C_TW_0(\lambda e^G)&=\lambda C_Te^G+O(\delta^4)
=-\eta\lambda y+O(\delta^4).
\end{align}
故
\begin{equation}
y^+=(1-\eta)y+\frac{\eta^3}{2}Vy+O(\delta^4).
\label{eq:mean-cubic}
\end{equation}
$\eta=1$ 时 $V^+=V^3/4+O(V^{7/2})$，两次使用 $\KL(q\Vert T)=V/2+O(V^{3/2})$，得到 $\KL(p^{\rm M}\Vert T)/\KL(q\Vert T)^3\to1$。
\end{proof}

\begin{proposition}[Mean 的全局三次跨度界]
对 $\eta=1$ 的精确 Mean 更新，令 $b=d_H(q,T)$，则
\begin{equation}
d_H(p^{\rm M},T)\le \frac{e^bb^2}{8+e^bb^2}\,b.
\end{equation}
\end{proposition}
\begin{proof}
还可给出全局三次跨度上界。$\eta=1$ 时 $\lambda\le d$、$w_{\max}\le e^b$。令 $F(t)=\log\E_qe^{t(r-\E_qr)}$，有 $F(0)=F'(0)=0$、$F''(t)\le b^2/4$；两次积分得 $d=F(1)\le b^2/8$。因此
\begin{equation}
b^+\le\frac{\lambda w_{\max}}{1+\lambda w_{\max}}b
\le\frac{e^bb^2}{8+e^bb^2}b.
\end{equation}
\end{proof}

\begin{theorem}[Mean 边界步长的渐近率]
固定正教师，在 $\eta=2$ 下精确 Mean 同步。若没有有限步到达教师，则 $kV_k\to1/8$、$k\KL(q_k\Vert T)\to1/16$；若某步已经到达教师，后续误差为零。
\end{theorem}
\begin{proof}
在 $\eta=2$，式~\eqref{eq:mean-cubic} 给 $V^+=V-8V^2+O(V^{5/2})$。若不在有限步精确到达教师，则 $V_k>0$ 且由全局稳定性 $V_k\to0$，
\begin{equation}
V_{k+1}^{-1}-V_k^{-1}=8+O(\sqrt{V_k})\to8.
\end{equation}
累加并除以轮数得 $kV_k\to1/8$，进而 $k\KL(q_k\Vert T)\to1/16$。几何更新在相同步长是两周期，Mean 的三次阻尼改变了边界行为。
\end{proof}

\begin{proposition}[Mean 的固定锚点大步长极限]
固定不同的正分布 $q,T$，令 $h=C_q\log(T/q)$、$c=\E_q[h]_+>0$。当 $\eta\to\infty$ 时，Mean 的归一化乘子满足 $\lambda/\eta\to c$，最优分布满足 $p_a^{\rm M}\to q_a[h_a]_+/c$。该极限不同于 Gibbs 的最高 reward 支持集极限。
\end{proposition}
\begin{proof}
固定锚点令 $\eta\to\infty$ 时，Mean 还有不同于 Gibbs 的极限。令 $h=C_qr$、$c=\E_q[h]_+>0$，则 $\lambda/\eta\to c$、$p_a^{\rm M}\to q_a[h_a]_+/c$。证明如下：最大 $h_a=M>0$ 处 KKT 与 $p_a/q_a\le1/q_a$ 给 $\lambda\ge q_a\eta M-q_a\log(1/q_a)$；上界 $\lambda\le\log\E_qe^{\eta h}\le\eta M$。所以 $\lambda/\eta$ 有严格正常数界。沿任一收敛子列，$W_0(\lambda e^{\eta h_a})/\eta\to[h_a]_+$，正、零、负三种情形分别使用 $W+\log W=\log\lambda+\eta h_a$、$W_0(\lambda)=O(\log\eta)$、$W_0(z)\le z$。有限求和归一化强制子列极限为 $c$，从而整个序列收敛。Gibbs 则只保留最高 reward 动作；两个极限不能互换。
\end{proof}

\section{有限 inner 的误差接口}
\begin{proposition}[四种内层损失到输出误差的接口]
若 $q_a\ge q_{\min}>0$，Part 或 Mean 对各自正确目标的超额损失 $\delta$ 蕴含 $d_H(p,p^\star)\le2\sqrt{\delta/q_{\min}}$；BPO 蕴含 $d_H(p,m)\le2\sqrt{\eta L_{\rm BPO}/q_{\min}}$。MDPO 则有 $\KL(p\Vert m)=\eta\delta$ 和 $\TV(p,m)\le\sqrt{\eta\delta/2}$，但无额外概率下界时不能控制统一 log-ratio。
\end{proposition}
\begin{proof}
由式~\eqref{eq:mean-excess}，Mean 超额损失 $\delta$ 控制 $\E_q\log^2(p/p^{\rm M})\le2\delta$；Part 超额损失正好等于该平方项（目标换成 $m$）。若 $q_a\ge q_{\min}$，附录~\ref{app:measure} 给
\begin{equation}
d_H(p,p^\star)\le2\sqrt{\delta/q_{\min}},\qquad
d_H(p,m)\le2\sqrt{\eta L_{\rm BPO}/q_{\min}}.
\end{equation}
MDPO 给 $\KL(p\Vert m)=\eta\delta$ 和 TV 界，但没有统一 log-ratio 界。取 $m=(1-\epsilon,\epsilon)$、$p=(1-\epsilon^2,\epsilon^2)$，KL 趋零，第二坐标的 log-ratio 为 $\log\epsilon\to-\infty$。
\end{proof}

\begin{proposition}[有限拟合的教师 KL 证书]
固定 $p,m,T>0$，令 $\varepsilon=\KL(p\Vert m)$、$H=\osc\log(m/T)$。则
\begin{align}
\KL(p\Vert T)&\le\KL(m\Vert T)+\varepsilon+H\sqrt{\varepsilon/2},\\
\KL(p\Vert T)&\le(1+1/u)\varepsilon+D^{\rm R}_{1+u}(m\Vert T),\qquad u>0.
\end{align}
\end{proposition}
\begin{proof}
直接检查教师逆向 KL 时，记 $\varepsilon=\KL(p\Vert m)$、$\ell=\log(m/T)$、$H=\osc\ell$。恒等式
\begin{equation}
\KL(p\Vert T)=\varepsilon+\KL(m\Vert T)+\sum_a(p_a-m_a)\ell_a
\end{equation}
配合区间中点和 Pinsker，得到上界 $\KL(m\Vert T)+\varepsilon+H\sqrt{\varepsilon/2}$。另一种按概率加权的证书是：对 $u>0$，取 $r_u\propto m^{1+u}T^{-u}$，展开 $\KL(p\Vert r_u)\ge0$ 得
\begin{equation}
\KL(p\Vert T)
\le(1+1/u)\KL(p\Vert m)+D^{\rm R}_{1+u}(m\Vert T).
\end{equation}
可在多个 $u$ 上取最小值；这里只用其精确有限 $u$ 形式作为保证。
\end{proof}

\begin{proposition}[条件无偏 SGD 的平均梯度范数界]
设损失 $L_{\rm sm}$ 光滑、下界为 $L_\star$，随机梯度条件无偏且方差不超过 $\sigma^2/M$。若 $0<\alpha\le1/L_{\rm sm}$，则
\begin{equation}
\frac1N\sum_{n<N}\E\|\nabla L(\theta_n)\|^2
\le\frac{2(L_0-L_\star)}{\alpha N}+\frac{L_{\rm sm}\alpha\sigma^2}{M}.
\end{equation}
该界本身不提供分布距离控制。
\end{proposition}
\begin{proof}
梯度范数界还需额外条件才能连接到这些输出证书。若随机梯度条件无偏、方差不超过 $\sigma^2/M$，且损失 $L_{\rm sm}$ 光滑，SGD 步长 $\alpha\le1/L_{\rm sm}$ 满足
\begin{equation}
\frac1N\sum_{n<N}\E\|\nabla L(\theta_n)\|^2
\le\frac{2(L_0-L_\star)}{\alpha N}+\frac{L_{\rm sm}\alpha\sigma^2}{M}.
\end{equation}
证明由光滑性、$\E\|g\|^2=\|\nabla L\|^2+\Var g$，每步下降至少 $\alpha\|\nabla L\|^2/2$，再求和。此式本身不控制超额损失或分布距离。
\end{proof}

\begin{proposition}[数据复用与平均覆盖的两个限制]
固定候选参数下的估计无偏，不一般蕴含数据依赖迭代参数下的条件无偏。另一方面，$\E_d\delta_s\le\bar\delta$ 不一般控制逐前缀误差；若另有 $d(s)\ge d_{\min}>0$，才可推出 $\delta_s\le\bar\delta/d_{\min}$。
\end{proposition}
\begin{proof}
复用数据还会破坏条件无偏前提。例如单个零均值样本 $X$，样本梯度 $\theta-X$ 在固定 $\theta$ 时无偏。从零更新得 $\theta_1=\alpha X$ 后，
\begin{equation}
\E[\theta_1-X\mid\theta_1]=(1-\alpha^{-1})\theta_1\ne\theta_1.
\end{equation}
因此独立动作的估计恒等式不能直接用于数据依赖的反复更新参数。平均前缀误差也不等于逐前缀误差；若另有 $d(s)\ge d_{\min}>0$ 才能由 $\E_d\delta_s\le\bar\delta$ 得 $\delta_s\le\bar\delta/d_{\min}$。
\end{proof}

\section{输出扰动、噪声与同步}
\begin{theorem}[确定性输出扰动的跨度递推]
令 $b_k=d_H(q_k,T_k)$。几何近似更新 $q_{k+1}\propto q_k^{1-\eta_k}T_k^{\eta_k}e^{e_k}$ 满足
\begin{equation}
b_{k+1}\le|1-\eta_k|b_k+\osc e_k+d_H(T_k,T_{k+1}).
\end{equation}
对 $0<\eta\le1$ 的 Mean 更新，若逐坐标输出 log 误差不超过 $\epsilon_k\le\bar\epsilon<1/2$，则 $b_{k+1}\le b_k-\eta b_k/(1+\eta b_k)+2\epsilon_k$，并有不变上界 $\max\{b_0,2\bar\epsilon/[\eta(1-2\bar\epsilon)]\}$。
\end{theorem}
\begin{proof}
几何端点的近似输出写为 $q_{k+1}\propto q_k^{1-\eta_k}T_k^{\eta_k}e^{e_k}$。取相对 $T_{k+1}$ 的对数并用跨度三角不等式：
\begin{equation}
b_{k+1}\le|1-\eta_k|b_k+\osc e_k+d_H(T_k,T_{k+1}).
\end{equation}
固定 $0<\eta<2$ 时，误差按系数 $c=|1-\eta|$ 作几何卷积；有界扰动上限 $B$ 给稳态界 $B/(1-c)$，趋零扰动给收敛。对 Mean 的 $0<\eta\le1$，若输出逐坐标 log 误差不超过 $\epsilon_k$，则
\begin{equation}
b_{k+1}\le b_k-\frac{\eta b_k}{1+\eta b_k}+2\epsilon_k.
\end{equation}
当 $\epsilon_k\le\bar\epsilon<1/2$，映射 $F(b)=b-\eta b/(1+\eta b)+2\bar\epsilon$ 递增，不动点为 $b_\infty=2\bar\epsilon/[\eta(1-2\bar\epsilon)]$。从 $\max(b_0,b_\infty)$ 起的比较递推单调降到 $b_\infty$，给出不变跨度与稳态上界。
\end{proof}

\begin{theorem}[几何更新的随机协方差传播]
固定正教师，$y_k=C_T\log(q_k/T)$、$\xi_k=C_Te_k$ 满足精确递推 $y_{k+1}=(1-\eta)y_k+\xi_k$。若初值确定，$\xi_k$ 为二阶可积鞅差且条件协方差为确定 $\Sigma_k$，则
\begin{equation}
\Cov(y_K)=\sum_{j<K}(1-\eta)^{2(K-1-j)}\Sigma_j.
\end{equation}
常协方差及 $0<\eta<2$ 时，极限为 $\Sigma/[\eta(2-\eta)]$；固定偏置 $\mu$ 的均值极限为 $\mu/\eta$。KL 噪声近似另需小噪声高阶矩条件。
\end{theorem}
\begin{proof}
固定教师，定义 $y_k=C_T\log(q_k/T)$、$\xi_k=C_Te_k$，几何更新精确满足
\begin{equation}
y_{k+1}=(1-\eta)y_k+\xi_k.
\end{equation}
若 $y_0$ 确定，$\xi_k$ 关于历史为鞅差且条件协方差为确定矩阵 $\Sigma_k$，不同时间交叉项由塔式期望为零，所以
\begin{equation}
\Cov(y_K)=\sum_{j<K}(1-\eta)^{2(K-1-j)}\Sigma_j.
\end{equation}
常协方差时极限为 $\Sigma/[\eta(2-\eta)]$。固定均值误差 $\mu$ 给均值极限 $\mu/\eta$。只有在另有小噪声高阶矩控制时，才能用局部 KL 展开转换成
\begin{equation}
\E\KL(q_\infty\Vert T)
=\frac{\mu^\top\diag(T)\mu}{2\eta^2}
+\frac{\operatorname{tr}(\diag(T)\Sigma)}{2\eta(2-\eta)}
+O(\epsilon^3),
\end{equation}
其中例如高斯误差的均值与标准差同时按 $\epsilon$ 缩小即可。若注入的是 reward 误差，输出噪声多乘 $\eta$，协方差多乘 $\eta^2$；不能混用两种步长依赖。
\end{proof}

\begin{proposition}[概率 soft sync 的局部稳定性]
固定正教师，几何、幂均值或相应 Mean 精确内层输出后执行概率混合 $q^+=(1-\zeta)q+\zeta p_{\rm out}$。其教师中心坐标的一阶系数为 $1-\zeta\eta$；若 $0<\zeta\eta<2$，教师局部渐近稳定。该命题不把参数 EMA 视作相同概率映射。
\end{proposition}
\begin{proof}
概率 soft sync 定义为 $q^+=(1-\zeta)q+\zeta p_{\rm out}$。在教师附近，几何、幂均值与 Mean 的一阶输出均为 $T+(1-\eta)(q-T)$，因此中心坐标变为 $(1-\zeta\eta)y+O(\|y\|^2)$。$0<\zeta\eta<2$ 给局部稳定，边界需高阶分析。参数 EMA 通常不是这个概率混合。
为给出稳定性的定量步骤，令 $a=|1-\zeta\eta|<1$。在充分小邻域内，二阶余项范数不超过 $(1-a)\|y\|/2$，故 $\|y^+\|\le(1+a)\|y\|/2$。该邻域不变，反复迭代几何趋零。
\end{proof}

\begin{proposition}[保持外层速率所需的内层精度]
若内层输出误差不超过 $C\rho_{\rm in}^{N_k}$，其中 $0<\rho_{\rm in}<1$，则要保证其不超过 $B\rho_{\rm out}^k$，充分取
\begin{equation}
N_k\ge\max\{0,[\log(C/B)+k\log(1/\rho_{\rm out})]/\log(1/\rho_{\rm in})\}.
\end{equation}
Mean 在 $\eta=1$ 的三次局部率要求输出误差为 $O(b_k^3)$；基于平方损失证书，一个充分条件是超额损失 $O(q_{\min,k}b_k^6)$。
\end{proposition}
\begin{proof}
若理想外层收缩因子为 $c$、内层输出误差 $\le C\rho_{\rm in}^{N_k}$，希望输出误差 $\le B\rho_{\rm out}^k$，充分条件是
\begin{equation}
N_k\ge\max\left\{0,\frac{\log(C/B)+k\log(1/\rho_{\rm out})}{\log(1/\rho_{\rm in})}\right\}.
\end{equation}
内层指数率必须先有独立依据。Mean 在 $\eta=1$ 的三次局部率需要输出跨度误差 $O(b_k^3)$；若使用损失证书，则充分要求超额损失为 $O(q_{\min,k}b_k^6)$。
\end{proof}

\section{共享参数的自然梯度投影}
\begin{proposition}[共享参数自然梯度的正交投影]
固定有限前缀和正权重，$B=\partial q_\theta/\partial\theta$、$W=\diag(d_s/q_{\theta,sa})$、$\mathcal I=B^\top WB$。若自由负自然梯度为 $v$，则参数自然梯度流诱导 $\dot q=Pv$，其中 $P=B\mathcal I^\dagger B^\top W$ 为 $W$ 正交投影。因此 $\|\dot q\|_W\le\|v\|_W$、$\dot L=-\|Pv\|_W^2\le0$。
\end{proposition}
\begin{proof}
固定有限前缀及正权重 $d_s$，令 $B=\partial q_\theta/\partial\theta$、$W=\diag(d_s/q_{\theta,sa})$、$\mathcal I=B^\top WB$。逐前缀自由负自然梯度为 $v$，有 $\nabla_\theta L=-B^\top Wv$。参数自然梯度流诱导
\begin{equation}
\dot q=B\mathcal I^\dagger B^\top Wv=Pv.
\end{equation}
令 $C=W^{1/2}B$，则 $W^{1/2}PW^{-1/2}=C(C^\top C)^\dagger C^\top$ 是欧氏正交投影。因此 $\|\dot q\|_W\le\|v\|_W$、$\dot L=-\|Pv\|_W^2\le0$。平方根自由速度不超过二，投影后也不超过二；这控制分布的 Fisher 速度，不控制参数欧氏速度。还可能 $Pv=0$ 而 $q_\theta\ne T$，故不保证消除表示误差。
\end{proof}

\section{条件更新与完整序列}
\begin{proposition}[条件改善不保证序列 KL 改善]
对有限自回归序列，KL 等于生成策略访问测度下的条件 KL 之和。即使每个前缀都执行精确幂均值更新、条件教师 KL 均不增加，完整序列逆向 KL 仍可能增加。
\end{proposition}
\begin{proof}
固定长度序列分布 $Q(y)=\prod_tq(y_t\mid y_{<t})$，教师序列分布为 $\mathsf T$。链式法则给
\begin{equation}
\KL(Q\Vert\mathsf T)=\sum_t\E_{s_t\sim Q}\KL(q(\cdot\mid s_t)\Vert T(\cdot\mid s_t)).
\end{equation}
更新条件策略同时改变访问权重。一个两步反例取根节点 $q_0=(0.99,0.01),T_0=(0.01,0.99)$；分支 A 中二者均为 $(1/2,1/2)$；分支 B 中 $q_B=(1-10^{-6},10^{-6})$、$T_B=(10^{-30},1-10^{-30})$。对全部条件分布同时做 $\eta=0.1$ 的精确更新，条件双向 KL 均不增加，但序列逆向 KL 由约 $5.19399$ 变为 $6.07445$（几何）、$7.09792$（平方根）或 $10.44581$（算术）。学生进入高误差分支的概率上升。
\end{proof}

\begin{proposition}[条件化与序列混合的差异]
完整序列算术混合的条件教师权重为 $\eta_s=\eta\mathsf T(s)/[(1-\eta)Q(s)+\eta\mathsf T(s)]$，一般依赖前缀。完整序列几何混合的条件概率为
\begin{equation}
\mathsf M(a\mid s)=q(a\mid s)^{1-\eta}T(a\mid s)^\eta Z_\eta(sa)/Z_\eta(s).
\end{equation}
因此两种序列混合一般都不同于逐前缀直接使用固定权重的条件混合。
\end{proposition}
\begin{proof}
对完整序列算术混合，先将前缀概率边缘化，再条件化：
\begin{equation}
\mathsf M(a\mid s)
=\frac{(1-\eta)Q(s)q(a\mid s)+\eta\mathsf T(s)T(a\mid s)}
{(1-\eta)Q(s)+\eta\mathsf T(s)}.
\end{equation}
分别提取两个条件分布的系数，教师权重为
\begin{equation}
\eta_s=\frac{\eta\mathsf T(s)}{(1-\eta)Q(s)+\eta\mathsf T(s)}.
\end{equation}
所以它一般不同于每个前缀使用固定 $\eta$ 的条件混合。序列几何混合也含后缀配分因子：
\begin{equation}
\mathsf M(a\mid s)=q(a\mid s)^{1-\eta}T(a\mid s)^\eta
\frac{Z_\eta(sa)}{Z_\eta(s)},\quad
Z_\eta(s)=\sum_zQ(z\mid s)^{1-\eta}\mathsf T(z\mid s)^\eta.
\end{equation}
具体地，未归一化的前缀 $sa$ 质量等于 $Q(s)^{1-\eta}\mathsf T(s)^\eta q(a\mid s)^{1-\eta}T(a\mid s)^\eta Z_\eta(sa)$；除以前缀 $s$ 的质量，公共前缀因子消去，即得到上式。
\end{proof}

\begin{proposition}[序列 KL 梯度中的后缀成本]
在有限长度、正支持、可微策略下，固定教师序列分布 $\mathsf T$，有
\begin{equation}
\nabla\KL(Q\Vert\mathsf T)=
\E_Q\sum_t\psi_t\sum_{j\ge t}\log[q(a_j\mid s_j)/T(a_j\mid s_j)].
\end{equation}
令后缀 KL 为 $V^q(s)$，对应条件负信号包含 $\log(T(a\mid s)/q(a\mid s))-V^q(sa)$，而不只有即时 log-ratio。
\end{proposition}
\begin{proof}
直接对序列 KL 求导，先用 $\nabla\log Q=\sum_t\psi_t$，再用条件 score 均值为零消去过去成本，得到
\begin{equation}
\nabla\KL(Q\Vert\mathsf T)
=\E_Q\sum_t\psi_t\sum_{j\ge t}\log\frac{q(a_j\mid s_j)}{T(a_j\mid s_j)}.
\end{equation}
记后缀值 $V^q(s)=\KL(Q(\cdot\mid s)\Vert\mathsf T(\cdot\mid s))$，即时负梯度信号应含 $\log(T(a\mid s)/q(a\mid s))-V^q(sa)$。固定测度的条件目标只保留即时项。若将幂均值问题直接定义在完整有限序列空间，前面的分布定理仍成立，但目标表达、采样和投影成本是另一问题。
\end{proof}

\section{旧锚点与前缀测度漂移}
\begin{proposition}[旧锚点引起的目标偏移]
\label{prop:stale-anchor}
固定正教师 $T$ 和几何目标 $m(q)\propto q^{1-\eta}T^\eta$，$\eta>0$。
对两个正锚点 $q_j,q_k$，
\begin{equation}
d_H(m(q_j),m(q_k))=|1-\eta|d_H(q_j,q_k).
\end{equation}
若输出 $p$ 对旧目标的拟合误差不超过 $\epsilon$，
则 $d_H(p,m(q_k))\le\epsilon+|1-\eta|d_H(q_j,q_k)$。
\end{proposition}
\begin{proof}
两个目标的对数比等于
$(1-\eta)\log(q_j/q_k)$ 加动作无关常数。
跨度消去常数，并将标量乘法转换为绝对值，得到等式。
再用 Hilbert 距离的三角不等式即得输出界。
这里的误差来自目标构造所用锚点的改变；即使精确积分，仍会存在。
采样分布与积分分布不一致导致的重要性加权问题需另行处理。
\end{proof}

\begin{proposition}[前缀测度改变的值与梯度代价]
\label{prop:prefix-drift}
设 $d,d'$ 为同一前缀空间上的概率测度，
标量或向量函数 $h$ 满足 $\sup_s\|h(s)\|\le C$。
则
\begin{equation}
\|\E_{d'}h-\E_dh\|\le2C\TV(d',d).
\end{equation}
对于条件目标 $J$，
\begin{align}
\E_{d'}J(q',T)-\E_dJ(q,T)
={}&\E_d[J(q',T)-J(q,T)]\\
&+(\E_{d'}-\E_d)J(q',T).
\end{align}
所以固定前缀测度下的进展，还需覆盖第二项才能推出总体进展。
\end{proposition}
\begin{proof}
有限前缀时，对 $\sum_s[d'(s)-d(s)]h(s)$ 用三角不等式，
得到 $C\sum_s|d'-d|=2C\TV(d',d)$；一般情形用总变差测度积分。
目标差中加减 $\E_dJ(q',T)$ 得分解。
若 $0\le J\le C$，标量项可用中心化 $J-C/2$ 将系数收紧为 $C$。
同样的论证适用于梯度，但需要独立的逐前缀梯度范数界；
目标值有界本身不提供梯度界。
\end{proof}

\section{有偏拟合与误差衰减速度}
\begin{proposition}[固定对数偏差产生倾斜稳态]
\label{prop:biased-fixed-point}
考虑几何更新
$q_{k+1,a}\propto q_{k,a}^{1-\eta}T_a^\eta e^{e_a}$，
其中 $T>0$、$e$ 固定、$0<\eta<2$。
其极限为
\begin{equation}
q_{\infty,a}=\frac{T_ae^{e_a/\eta}}{\sum_bT_be^{e_b/\eta}},
\qquad d_H(q_\infty,T)=\frac{\osc e}{\eta}.
\end{equation}
若 $1<\eta<2$ 而误差允许逐轮交替，统一输入界
$\limsup d_H(q_k,T)\le\sup_k\osc e_k/(2-\eta)$ 可以达到。
\end{proposition}
\begin{proof}
令 $y_k=C_T\log(q_k/T)$，则归一化常数被中心化消去，递推为
$y_{k+1}=(1-\eta)y_k+C_Te$。
因 $|1-\eta|<1$，
$y_k=(1-\eta)^ky_0+[1-(1-\eta)^k]C_Te/\eta$，
取极限并归一化得所述稳态和跨度。
为验证交替误差的界，取
$e_k=(-1)^{k+1}e$ 和
$y_k=(-1)^kC_Te/(2-\eta)$。
直接代入递推，两侧相等，且每步跨度为 $\osc e/(2-\eta)$。
任意初值与该特解之差按 $|1-\eta|^k$ 衰减。
固定偏差与最坏时变偏差因此有不同稳态常数。
\end{proof}

\begin{proposition}[几何衰减误差的卷积与同速因子]
\label{prop:equal-rate-errors}
设 $b_{k+1}\le c b_k+E_k$，$0<c<1$，
$0\le E_k\le B\rho^k$、$0<\rho<1$。则
\begin{equation}
b_K\le c^Kb_0+
\begin{cases}
B(c^K-\rho^K)/(c-\rho),&c\ne\rho,\\
BKc^{K-1},&c=\rho.
\end{cases}
\end{equation}
若 $c=0$，则直接有 $b_K\le E_{K-1}$，$K\ge1$。
\end{proposition}
\begin{proof}
迭代不等式得
$b_K\le c^Kb_0+B\sum_{j=0}^{K-1}c^{K-1-j}\rho^j$。
当 $c\ne\rho$，乘以 $c-\rho$ 后相邻项抵消，
得到 $c^K-\rho^K$；当 $c=\rho$，每一项均为 $c^{K-1}$，
共有 $K$ 项。该式解释了误差与精确收缩同速时出现的额外线性因子。
\end{proof}

\begin{proposition}[Mean 的小噪声局部误差底]
\label{prop:mean-noise-floor}
取几何信号、Mean 更新及 $\eta=1$，令
$b_k=d_H(q_k,T)$。若输出对精确 Mean 解的跨度误差至多为 $E$，
前文局部界写成
\begin{equation}
b_{k+1}\le h(b_k)b_k+E,\qquad
h(b)=\frac{e^bb^2}{8+e^bb^2}.
\end{equation}
若 $b_0\le2E$ 且 $h(2E)\le1/2$，则
\begin{equation}
\sup_k b_k\le2E,\qquad
\limsup_k b_k\le\frac{E}{1-h(2E)}=E+O(E^3).
\end{equation}
对趋零的局部轨迹，误差 $E_k=o(b_k^3)$ 足以保留精确更新的三次首项。
\end{proposition}
\begin{proof}
$h$ 在非负半轴非减。若 $b_k\le2E$，则
$b_{k+1}\le h(2E)2E+E\le2E$，归纳得到不变邻域。
在此邻域将 $h(b_k)$ 换为固定 $h(2E)<1$，迭代并取上极限，
得到 $E/[1-h(2E)]$。由 $h(2E)=O(E^2)$ 得最后展开。
若近似输出与精确输出的 Hilbert 距离为 $E_k$，
三角不等式的两个方向给其到教师距离之差的绝对值至多为 $E_k$，
故 $o(b_k^3)$ 不改变存在的三次首项。
以 Mean 超额损失 $\delta_k$ 提供残差证书时，
$E_k\le2\sqrt{\delta_k/q_{\min,k}}$，
所以 $\delta_k=o(q_{\min,k}b_k^6)$ 是一个充分条件。
该条件依赖原有局部轨迹与精确解证书，不声称对任意共享参数更新成立。
\end{proof}

\section{局部收缩率与全局边界}
\begin{proposition}[几何更新在稳定区间之外的两个极限]
\label{prop:geometric-unstable-limit}
固定 $q_0\ne T$、$T>0$、$\eta>2$，
令 $u_a=\log(q_{0,a}/T_a)$，最大与最小值集合分别为
$I_+$、$I_-$。精确几何迭代的偶数子列趋于 $T$ 在 $I_+$ 上的条件分布，
奇数子列趋于 $T$ 在 $I_-$ 上的条件分布。
此时 $\KL(T\Vert q_k)\to\infty$，但
$\KL(q_k\Vert T)\le\log(1/T_{\min})$ 始终有限。
\end{proposition}
\begin{proof}
闭式轨迹为 $q_k(a)\propto T_ae^{(1-\eta)^ku_a}$。
偶数时指数系数趋正无穷，除以最大指数项后，
集合外的指数趋零，集合内的比例仍为 $T_a/T_b$。
奇数时同理保留最小值集合。
$u$ 非常数，两个集合均为真子集，因此每个子列都有正教师质量被学生压到零；
$-\sum_aT_a\log q_k(a)$ 发散，而 $\sum_aT_a\log T_a$ 固定。
反向 KL 则满足
$\sum_aq_a\log q_a-\sum_aq_a\log T_a\le\log(1/T_{\min})$，
因为第一项非正。跨度发散与两个方向的 KL 有不同关系。
\end{proof}

\begin{proposition}[几何逆向 KL 没有统一严格乘法收缩]
\label{prop:no-global-kl-factor}
固定正二元教师 $T$ 与 $0<\eta<1$。
不存在仅依赖 $T,\eta$ 的 $c<1$，使所有正 $q$ 都满足
$\KL(m(q)\Vert T)\le c\KL(q\Vert T)$。
\end{proposition}
\begin{proof}
取 $q_\epsilon=(1-\epsilon,\epsilon)$。
几何目标的第二坐标与第一坐标之比为
$(\epsilon/(1-\epsilon))^{1-\eta}(T_2/T_1)^\eta\to0$，
所以 $q_\epsilon$ 与 $m(q_\epsilon)$ 都趋于 $(1,0)$。
对固定正教师，逆向 KL 在闭单纯形上连续，约定 $0\log0=0$。
两者 KL 均趋于 $\log(1/T_1)>0$，比值趋一，
与任意统一 $c<1$ 矛盾。
该反例不影响教师附近的局部比例 $(1-\eta)^2$，
也不影响全局 Hilbert 收缩。
\end{proof}


\chapter{TOP-D 目标教师的退化与性质}
\label{app:topd}
\section{参数对应与精确可达性}
本附录比较 TOP-D 的目标教师及相关分布性质，依据其 v1 第 3.1 节、定理 4.2、4.4 和附录 C.1、C.2\footfullcite{topd2026}。固定同一前缀，TOP-D 的教师混合权重记为 $\alpha$，教师为 $T$，构造混合时的学生记为 $q$。
\begin{theorem}[TOP-D 目标教师的精确退化]
固定同一条件分布空间的 $q,T>0$。取 $\kappa=1,\eta=\alpha\in(0,1)$ 时，镜像目标教师精确等于 TOP-D 的算术教师 $(1-\alpha)q+\alpha T$，且 $Z=1$。该分布唯一最小化加权正向 KL，超额损失为 $\KL(m\Vert p)$。
\end{theorem}
\begin{proof}
\begin{equation}
\boxed{\kappa=1,\qquad \eta=\alpha,\qquad
m_{1,\eta}=(1-\eta)q+\eta T.}
\label{eq:topd-match}
\end{equation}
这在目标分布上是精确等式，无需小步近似，也无需对两分布接近作假设。两者必须作用于同一个条件分布空间。

目标教师由正文的算术端点最优化问题产生：
\begin{equation}
m_{1,\eta}=\argmin_{p\in\Delta}
\{(1-\eta)\KL(q\Vert p)+\eta\KL(T\Vert p)\}.
\end{equation}
其概率梯度为 $-[(1-\eta)q_a+\eta T_a]/p_a$。约束乘子使 $p_a$ 与分子成比例；分子总和已经为一，所以 $p=m$。完整误差恒等式是
\begin{equation}
F(p)-F(m)=\KL(m\Vert p).
\end{equation}
因此镜像解与该重心解都等于 TOP-D 的目标教师。
\end{proof}

\begin{corollary}[算术教师的总体可达性]
在完整支持、精确期望和目标可实现时，MDPO、Part、BPO 的总体最优分布都等于算术目标教师。Mean 对非恒定信号一般具有不同的最优分布。
\end{corollary}
\begin{proof}
将上一退化定理的算术 $m$ 代入命题~\ref{prop:inner-roots}，MDPO 的 KL gap、Part 的零残差和 BPO 的常数残差条件分别唯一确定 $p=m$。同一命题的 Mean KKT 论证排除非恒定信号下的 Gibbs 解。可达性指总体最优分布；有限参数更新是否到达，应由拟合误差证书判断。
\end{proof}

\section{监督信号、中心与教师目标}
\begin{proposition}[算术端点的信号、中心与教师目标]
取 $\kappa=1$，则
\begin{equation}
f=\log(1-\eta+\eta T/q),\quad A=f/\eta,\quad K_T=\KL(q\Vert m)/\eta.
\end{equation}
本文反推的教师目标为 $J_{1,\eta}(p;T)=\JS_\eta(p,T)/[\eta(1-\eta)]$，其在 $q$ 处步长为 $\eta$ 的负熵镜像解等于算术教师。
\end{proposition}
\begin{proof}
算术端点的配分函数恒为一，故
\begin{equation}
f_{1,\eta}=\log(m/q)=\log(1-\eta+\eta T/q),\qquad
A_{1,\eta}=f_{1,\eta}/\eta,\qquad
K_T=\KL(q\Vert m)/\eta.
\label{eq:topd-signal}
\end{equation}
第一项恢复 TOP-D 式（5）的未缩放信号。第二项是本文内层归一化使用的信号，比较方差时需保留 $1/\eta^2$。

本文由逆向积分定义的教师目标在此端点为
\begin{equation}
J_{1,\eta}(p;T)=
\frac{(1-\eta)\KL(p\Vert M_p)+\eta\KL(T\Vert M_p)}
{\eta(1-\eta)},\qquad M_p=(1-\eta)p+\eta T.
\end{equation}
推导已在附录~\ref{app:objective} 给出。这里 $M_p$ 随候选学生 $p$ 变化，用于定义同一个教师目标；本轮目标 $m=M_q$ 在内层冻结。对 $J$ 在 $q$ 处执行步长 $\eta$ 的负熵镜像下降，恰好得到算术目标。这是本文对该目标教师的变分解释，不把它与其他同名目标混用。
\end{proof}

\section{覆盖与散度收缩}
\begin{proposition}[算术教师的覆盖与信号下界]
设 $m=(1-\eta)q+\eta T$、$0<\eta<1$，则 $m_a\ge(1-\eta)q_a$、$m_a\ge\eta T_a$，并有
\begin{equation}
f\ge\log(1-\eta),\quad
\KL(q\Vert m)\le-\log(1-\eta),\quad \KL(T\Vert m)\le-\log\eta.
\end{equation}
\end{proposition}
\begin{proof}
由算术混合逐坐标有 $m_a\ge(1-\eta)q_a$、$m_a\ge\eta T_a$，因此
\begin{equation}
f_{1,\eta}\ge\log(1-\eta),\quad
\KL(q\Vert m)\le-\log(1-\eta),\quad
\KL(T\Vert m)\le-\log\eta.
\end{equation}
前者控制教师极小概率造成的负监督尾部，后两者控制到两端分布的覆盖代价。
\end{proof}

\begin{proposition}[算术教师的精确收缩与闭式外层]
固定正教师和 $0<\eta<1$，精确算术更新满足 $\|m-T\|_1=(1-\eta)\|q-T\|_1$。反复更新的闭式为 $q_k=T+(1-\eta)^k(q_0-T)$。
\end{proposition}
\begin{proof}
对于精确更新，
\begin{equation}
m-T=(1-\eta)(q-T),\qquad
\|m-T\|_1=(1-\eta)\|q-T\|_1.
\end{equation}
反复精确更新的闭式为 $q_k=T+(1-\eta)^k(q_0-T)$。所以原文定理 4.4 所用的算术收缩结构在这个条件分布空间直接恢复。
\end{proof}

\begin{proposition}[算术端点的凸散度收缩]
对任意凸函数 $\varphi$、$\varphi(1)=0$，
\begin{equation}
D_\varphi(m\Vert T)\le(1-\eta)D_\varphi(q\Vert T).
\label{eq:topd-all-f}
\end{equation}
\end{proposition}
\begin{proof}
逐坐标有 $m_a/T_a=(1-\eta)q_a/T_a+\eta$。凸性给
\begin{equation}
\varphi(m_a/T_a)
\le(1-\eta)\varphi(q_a/T_a)+\eta\varphi(1).
\end{equation}
乘以 $T_a$ 并求和，第二项为零。取 $\varphi(x)=x\log x$、$-\log x$、$(\sqrt x-1)^2$、$|x-1|/2$，分别得到两个 KL 方向、平方 Hellinger 和 TV 的界。
\end{proof}

式~\eqref{eq:topd-all-f} 是算术映射自身的凸性结论。一般幂均值保留所有凸生成散度的非增性，但不自动具有同一个 $(1-\eta)$ 全局因子；平方根点的 Hellinger 因子由它的球面结构另行给出。

\section{有限拟合的误差累积}
\begin{theorem}[算术外层的拟合误差累积]
若 $\epsilon_k=\|q_{k+1}-m(q_k,T)\|_1$，则
\begin{equation}
\|q_K-T\|_1\le(1-\eta)^K\|q_0-T\|_1+\sum_{k<K}(1-\eta)^{K-1-k}\epsilon_k.
\end{equation}
若 $\limsup\epsilon_k\le\epsilon_\infty$，则 $\limsup\|q_K-T\|_1\le\epsilon_\infty/\eta$。固定测度的条件距离平均满足同一界。
\end{theorem}
\begin{proof}
令 $\epsilon_k=\|q_{k+1}-m(q_k,T)\|_1$。三角不等式给
\begin{equation}
\|q_{k+1}-T\|_1\le(1-\eta)\|q_k-T\|_1+\epsilon_k.
\end{equation}
归纳展开：
\begin{equation}
\|q_K-T\|_1\le(1-\eta)^K\|q_0-T\|_1
+\sum_{k=0}^{K-1}(1-\eta)^{K-1-k}\epsilon_k.
\label{eq:topd-errors}
\end{equation}
若 $\limsup\epsilon_k\le\epsilon_\infty$，对任意 $\delta>0$ 取足够晚的 $k_0$ 使 $\epsilon_k\le\epsilon_\infty+\delta$。有限前缀和趋零，后段几何和不超过 $(\epsilon_\infty+\delta)/\eta$。令 $\delta\downarrow0$，得到 $\limsup\|q_K-T\|_1\le\epsilon_\infty/\eta$。这恢复原文误差累积与稳态上界的形式。

固定前缀测度下可以逐条件取平均，再用同一个线性递推。如果用 MDPO 的超额损失 $\delta_k$，Pinsker 给 $\epsilon_k\le\sqrt{2\eta\delta_k}$；若用 BPO，则先应用附录~\ref{app:dynamics} 的跨度证书。本文在这里补上的是不同 inner loss 到外层输出误差的接口。
\end{proof}

\section{方差定理的尺度对应}
\begin{theorem}[算术监督的方差界与小步尺度]
固定 $0<\eta<1$，$f=\log(1-\eta+\eta T/q)$，且 $\|\psi\|\le S$。令 $C_\star=\sup_{u>0}\log^2(1+u)/u$，则
\begin{equation}
\Var_q(f\psi)\le S^2[\log^2(1-\eta)+C_\star\eta].
\end{equation}
对正分布类取上确界，$\sup\Var_qf=\Theta(\eta)$，而 $A=f/\eta$ 满足 $\sup\Var_qA=\Theta(\eta^{-1})$ 当 $\eta\downarrow0$。
\end{theorem}
\begin{proof}
TOP-D 定理 4.2 在有限词表和 score 范数不超过 $S$ 的条件下，给未缩放信号的梯度方差上界
\begin{equation}
S^2|\mathcal V|\max\{\log^2(1-\eta),C_\star\eta\}.
\end{equation}
在同一条件下，利用概率归一化可以得到无词表大小因子的界。对 $f=\log(1-\eta+\eta\rho)$，将 $\rho<1$ 与 $\rho\ge1$ 分开：
\begin{align}
\E_q[f^2\mathbf1_{\rho<1}]&\le\log^2(1-\eta),\\
\E_q[f^2\mathbf1_{\rho\ge1}]
&\le C_\star\eta\E_q[\rho\mathbf1_{\rho\ge1}]
\le C_\star\eta.
\end{align}
其中 $C_\star=\sup_{u>0}\log^2(1+u)/u$，和原文单变量上界使用同一个常数。因此
\begin{equation}
\Var_q(f\psi)\le S^2[\log^2(1-\eta)+C_\star\eta].
\end{equation}
这个推导控制条件单动作梯度；前缀加权后只要使用同一 score 上界，二阶矩界仍成立。

未缩放信号的最坏二阶矩在 $\eta\downarrow0$ 时为 $\Theta(\eta)$。上界由前式，下界取 $q_1=\eta/2,T_1=1/2$，稀有动作 $f_1=\log(2-\eta)$，贡献至少 $\eta\log^2(2-\eta)/2$。对二元分布，中心化方差 $q_1q_2(f_1-f_2)^2$ 也有同阶，因为另一项 $f_2=O(\eta)$。

所以本文标准化监督满足
\begin{equation}
\sup_{q,T}\Var_qA_{1,\eta}=\Theta(\eta^{-1}).
\end{equation}
这与未缩放信号随小 $\eta$ 衰减相容：两者相差 $\eta^2$。它也说明算术端点虽有固定强度方差界，却不属于统一小步标准化方差有界的 $0<\kappa\le1/2$ 区间。
\end{proof}

\section{比较结论的作用范围}
\begin{center}
\small
\begin{tabular}{p{38mm}p{93mm}}
\toprule
对象 & 精确对应或进一步性质\\
\midrule
目标教师 & $\kappa=1,\eta=\alpha$ 时完全一致，$Z=1$\\
未缩放信号 & $f=\log(1-\eta+\eta T/q)$ 完全一致\\
最优化来源 & 加权正向 KL 重心，超额损失为 $\KL(m\Vert p)$\\
教师目标解释 & 本文的 $J_{1,\eta}$ 是归一化加权 JS，其熵镜像步等于混合教师\\
精确与近似外层 & 算术 $L_1$ 收缩、几何误差累积、$\epsilon_\infty/\eta$ 稳态界\\
其他教师散度 & 凸性给任意凸生成散度的 $(1-\eta)$ 上界\\
监督方差 & 未缩放 $f$ 与标准化 $A=f/\eta$ 分开比较\\
\bottomrule
\end{tabular}
\end{center}

上述比较针对相同的条件分布空间。若直接在完整序列空间定义算术混合，同一代数也成立；逐 token 固定权重混合一般不等于完整序列混合，不能把两个空间的距离直接认作同一指标。这里保留目标教师与定理层面的对应。

\chapter*{来源对应与复现}
\addcontentsline{toc}{chapter}{来源对应与复现}
本稿对五组本地理论来源按数学依赖重新组织。外部散度、PMD-Mean 和 TOP-D 资料在对应位置引用；本地来源保留原文件，以便追溯定义及历史语境。同一命题在多个版本中出现时，本稿只维护一个完整证明。

随源文件提供的 \path{source-map.csv} 记录原文件、命题位置、源文件校验和及新稿主题归属。主题合并可能是多对一关系；对应表用于追溯材料，不将重复版本计为独立理论结果。原文没有命题环境的推导另以文件条目登记。

\begin{longtable}{p{24mm}p{57mm}p{47mm}}
\toprule
来源目录 & 主要材料 & 新稿归属\\
\midrule
\endhead
\path{alpha-divergence/} & 密度比矩、散度端点、score 梯度、几何重加权测度 & 附录 A、B、D；统一教师记号与 Hellinger 常数\\
\path{book/} & 重心、$J$、四种 inner、信息访问、估计、动力学与边界 & 正文及附录 A--H 的主要数学来源\\
\path{book-old/} & Mean 修订、KL 控制族、几何与 Mean 的同步推导 & 附录 C、E、F、G；保留历史差异的数学内容\\
\path{Notes/} & 整合证明、成对 BPO、无重复动作、方差与误差连接 & 附录 B、D、F、G，特别是 U-统计量完整分解\\
\path{paper/} & 参数化叙事、曲率、受限投影与有限步论证 & 正文结构及附录 C、D；精确展开曲率与驻点证明\\
\bottomrule
\end{longtable}

来源中的纯实验配置、旧生产模块路径与任务成绩保留在原材料中。本稿的反例和验证针对所写数学对象。来源名相同不表示统计口径相同：头部可能由教师或锚点选出，条件动作采样可能来自当前学生或旧锚点，必须按本稿每条定理的条件代入。

一些跨主题内容的归属为：多教师重心与通用单调信号在附录 B；独立拟合形状、共享参数投影和标量谱调度在 C、G；最坏矩、相对噪声和绝对噪声准则在 D；配分函数、自动微分与非保守场在 E、F；教师漂移、同步及序列反例在 G。幂散度梯度中的几何重加权分布 $p(T/p)^\kappa$ 与本轮幂均值目标 $m_{\kappa,\eta}$ 是不同对象，前者用于积分测度解释，后者由重心问题定义。

\section*{数学核验的范围}
数值核验采用有限词表的确定性求和，对照解析恒等式、自动微分和独立的概率公式。重点包括重心超额损失、$J$ 的梯度与 JS 端点、四种 inner 的梯度和 Hessian、成对估计的协方差、形状端点、TOP-D 算术收缩以及条件与序列的反例。测试入口为 \path{tests/reference/test_core_monograph.py}。

2026 年 10 月 15 日的 CPU 核验中，该入口的 46 个参数化测试全部通过；仓库全套测试为 316 项通过。积分目标采用 96 点 Gauss--Legendre 求积与解析梯度对照，主要恒等式容差为 $10^{-11}$ 或 $10^{-12}$。U-统计量使用三动作词表、$M=2,3,4$ 的全部有序样本组合；无重复抽样直接比较协方差矩阵，检查半正定差值。局部三次率采用逐渐接近教师的固定方向，不以有限差分噪声替代极限系数。运行环境与命令保存于 \path{validation.json}。

数值测试用于暴露符号、尺度和实现错误；有限实例通过不构成定理证明。理论正文中的结论以相应解析证明和显式条件为依据。测试结果不用于声称共享参数大模型的训练收敛或任务准确率。


在仓库根执行 \texttt{make theory-core} 构建完整稿和正文阅读版；在本目录执行 \texttt{make} 等价。构建复用相邻 \path{../book/elegantbook.cls}，使用 XeLaTeX、latexmk 与 Biber。正文阅读版由同一主源通过条件装配生成，附录引用保留完整稿的编号。

\fi
\backmatter
\printbibliography[heading=bibintoc,title={参考文献}]
\end{document}
